%=====================================================================
%  主文件：中文精译 + 学术审校版
%  原文：A. Balasubramanian, J. Teschner,
%        "Supersymmetric field theories and geometric Langlands:
%         The other side of the coin", arXiv:1702.06499v2
%        （String-Math 2016 会议文集）
%  编译：XeLaTeX（必须使用 XeLaTeX 才能加载 ctex 中文字库）
%=====================================================================
\documentclass[12pt,reqno]{article}

% !TEX program = xelatex
%=====================================================================
%  个人宏包与版式设定（preamble）
%  说明：本文件是完整可编译文档的导言区，供主文件 \input 使用。
%        宏包顺序不可随意打乱；hyperref 必须最后加载。
%=====================================================================

%---------------------------------------------------------------------
% 一、版心
%---------------------------------------------------------------------
% 页面尺寸与边距（参照原文版式，便于与原文逐页对照）
\setlength{\textheight}{210mm}
\setlength{\textwidth}{160mm}
\setlength{\oddsidemargin}{0mm}
\setlength{\evensidemargin}{0mm}
\setlength{\topmargin}{0mm}
\setlength{\headheight}{12pt}
\setlength{\headsep}{4mm}
\setlength{\parindent}{2em}
\setlength{\parskip}{0.5ex}
\linespread{1.1}
\sloppy
\flushbottom

%---------------------------------------------------------------------
% 二、数学宏包
%---------------------------------------------------------------------
\usepackage{amsmath}
\usepackage{amsthm}
\usepackage{amsfonts}
\usepackage{amssymb}
\usepackage{mathtools}

%---------------------------------------------------------------------
% 三、中文支持（XeLaTeX + ctex）
%     若本机无 Windows 中文字体，可将 fontset 改为 fandol
%---------------------------------------------------------------------
\usepackage[UTF8]{ctex}

%---------------------------------------------------------------------
% 四、图形、颜色、表格、链接
%---------------------------------------------------------------------
\usepackage{graphicx}
\usepackage{xcolor}
\usepackage{booktabs}
\usepackage{longtable}
\usepackage{array}
\usepackage{multirow}
\usepackage{enumitem}
\setlist{nosep,leftmargin=*}
\usepackage[unicode,colorlinks=true,linkcolor=blue!55!black,
            citecolor=green!45!black,urlcolor=blue!60!black]{hyperref}

%---------------------------------------------------------------------
% 五、定理类环境（与原文对应）
%---------------------------------------------------------------------
\theoremstyle{plain}
\newtheorem{thm}{定理}
\newtheorem{claim}{断言}
\newtheorem{propn}{命题}
\newtheorem{lem}{引理}
\newtheorem{cor}{推论}
\newtheorem{conj}{猜想}

\theoremstyle{definition}
\newtheorem{defn}{定义}

\theoremstyle{remark}
\newtheorem{rem}{注}

%---------------------------------------------------------------------
% 六、记号缩写（与原文宏定义一一对应，保证公式逐字无损）
%---------------------------------------------------------------------
% 希腊字母缩写
\newcommand{\al}{\alpha}
\newcommand{\be}{\beta}
\newcommand{\ga}{\gamma}
\newcommand{\Ga}{\Gamma}
\newcommand{\de}{\delta}
\newcommand{\De}{\Delta}
\newcommand{\ep}{\epsilon}
\newcommand{\la}{\lambda}
\newcommand{\om}{\omega}
\newcommand{\Om}{\Omega}
\newcommand{\si}{\sigma}
\newcommand{\up}{\Upsilon}
\newcommand{\vf}{\varphi}
\newcommand{\vp}{\varpi}
\newcommand{\ze}{\zeta}
\newcommand{\tpi}{\tilde{\pi}}

% 常用算符与结构
\newcommand{\mc}{\mathcal}
\newcommand{\on}{\operatorname}
\newcommand{\id}{\mathrm{id}}
\newcommand{\pa}{\partial}
\newcommand{\ot}{\otimes}
\newcommand{\ra}{\to}
\newcommand{\ti}{\times}
\newcommand{\fr}[2]{\textstyle\frac{#1}{#2}}
\newcommand{\oh}{\frac{1}{2}}
\newcommand{\st}{\scriptstyle}
\newcommand{\sst}{\scriptscriptstyle}
\newcommand{\dst}{\displaystyle}
\newcommand{\tst}{\textstyle}
\newcommand{\ew}{\hspace*{-1mm}}
\newcommand{\ppe}{\hspace*{-2.5mm}}
\newcommand{\bb}{\hspace*{-1mm}}
\newcommand{\no}[1]{\!:\!#1\!:\!}

% 黑板体（原文 \C, \R, \Z 等）
\newcommand{\C}{\mathbb{C}}
\newcommand{\BR}{\mathbb{R}}
\newcommand{\R}{\mathbb{R}}
\newcommand{\BZ}{\mathbb{Z}}
\newcommand{\Z}{\mathbb{Z}}
\newcommand{\BQ}{\mathbb{Q}}
\newcommand{\BN}{\mathbb{N}}
\newcommand{\BP}{\mathbb{P}}
\newcommand{\BT}{\mathbb{T}}
\newcommand{\BS}{\mathbb{S}}
\newcommand{\BA}{\mathbb{A}}
\newcommand{\BB}{\mathbb{B}}
\newcommand{\BD}{\mathbb{D}}
\newcommand{\BH}{\mathbb{H}}
\newcommand{\BJ}{\mathbb{J}}
\newcommand{\BI}{\mathbb{I}}
\newcommand{\BF}{\mathbb{F}}
\newcommand{\BC}{\mathbb{C}}

% Fraktur 与算子
\newcommand{\g}{\mathfrak{g}}
\newcommand{\fg}{\mathfrak{g}}
\newcommand{\hfg}{\hat{\mathfrak{g}}}
\newcommand{\hfb}{\hat{\mathfrak{b}}}
\newcommand{\fsl}{\mathfrak{sl}}
\newcommand{\hfsl}{\widehat{\mathfrak{sl}}}
\newcommand{\FB}{\mathfrak{B}}
\newcommand{\one}{\mathfrak{1}}
\newcommand{\two}{\mathfrak{2}}
\newcommand{\three}{\mathfrak{3}}
\newcommand{\Bun}{\operatorname{Bun}}
\newcommand{\Dmod}{\mathrm{D}_k\text{-}\mathrm{mod}}
\newcommand{\ppart}{(\!(t)\!)}

% 大写花体（原文 \CA … \CZ）
\newcommand{\CA}{\mathcal{A}}
\newcommand{\CB}{\mathcal{B}}
\newcommand{\CC}{\mathcal{C}}
\newcommand{\CD}{\mathcal{D}}
\newcommand{\CE}{\mathcal{E}}
\newcommand{\CF}{\mathcal{F}}
\newcommand{\CG}{\mathcal{G}}
\newcommand{\CH}{\mathcal{H}}
\newcommand{\CI}{\mathcal{I}}
\newcommand{\CJ}{\mathcal{J}}
\newcommand{\CK}{\mathcal{K}}
\newcommand{\CL}{\mathcal{L}}
\newcommand{\CM}{\mathcal{M}}
\newcommand{\CN}{\mathcal{N}}
\newcommand{\CO}{\mathcal{O}}
\newcommand{\CP}{\mathcal{P}}
\newcommand{\CS}{\mathcal{S}}
\newcommand{\CU}{\mathcal{U}}
\newcommand{\CV}{\mathcal{V}}
\newcommand{\CW}{\mathcal{W}}
\newcommand{\CZ}{\mathcal{Z}}

% 大写无衬线体（原文 \SA … \SZ，用于算子、空间与标签）
\newcommand{\SA}{\mathsf{A}}
\newcommand{\SB}{\mathsf{B}}
\newcommand{\SC}{\mathsf{C}}
\newcommand{\SD}{\mathsf{D}}
\newcommand{\SE}{\mathsf{E}}
\newcommand{\SF}{\mathsf{F}}
\newcommand{\SG}{\mathsf{G}}
\newcommand{\SH}{\mathsf{H}}
\newcommand{\SI}{\mathsf{I}}
\newcommand{\SJ}{\mathsf{J}}
\newcommand{\SK}{\mathsf{K}}
\newcommand{\SL}{\mathsf{L}}
\newcommand{\SM}{\mathsf{M}}
\newcommand{\SN}{\mathsf{N}}
\newcommand{\SO}{\mathsf{O}}
\newcommand{\SP}{\mathsf{P}}
\newcommand{\SQ}{\mathsf{Q}}
\newcommand{\SR}{\mathsf{R}}
\renewcommand{\SS}{\mathsf{S}}
\newcommand{\ST}{\mathsf{T}}
\newcommand{\SU}{\mathsf{U}}
\newcommand{\SV}{\mathsf{V}}
\newcommand{\SW}{\mathsf{W}}
\newcommand{\SX}{\mathsf{X}}
\newcommand{\SY}{\mathsf{Y}}
\newcommand{\SZ}{\mathsf{Z}}

% 小写无衬线体
\newcommand{\se}{\mathsf{e}}
\renewcommand{\sf}{\mathsf{f}}
\newcommand{\spp}{\mathsf{p}}
\newcommand{\mss}{\mathsf{s}}
\newcommand{\sv}{\mathsf{v}}
\newcommand{\su}{\mathsf{u}}
\newcommand{\sx}{\mathsf{x}}

% Dirac 记号
\newcommand{\bra}[1]{\langle #1|}
\newcommand{\ket}[1]{|#1\rangle}

% 算符
\DeclareMathOperator*{\Res}{Res}
\DeclareMathOperator{\Tr}{Tr}
\DeclareMathOperator{\Sym}{Sym}
\DeclareMathOperator{\Rep}{Rep}
\DeclareMathOperator{\Loc}{Loc}
\DeclareMathOperator{\Hom}{Hom}
\DeclareMathOperator{\End}{End}
\DeclareMathOperator{\Aut}{Aut}
\DeclareMathOperator{\Pic}{Pic}
\DeclareMathOperator{\Crit}{Crit}
\DeclareMathOperator{\Vol}{Vol}

% 原文中的 \= 重定义：\={\label} 表示“由该式给出”
\newcommand{\eqrefby}[1]{\stackrel{\eqref{#1}}{=}}

% 原文自定义宏的等价物（原文直接使用，必须提供，否则编译报错）
\newcommand{\rf}[1]{\eqref{#1}}            % 原文：\rf{label} 即公式引用
\newcommand{\glos}{\textsuperscript{g)}}   % 原文：术语表上标标记
\newcommand{\crd}{}                        % 原文：红色标记开关，排版时不显示
\newcommand{\hba}{\hat{\mathfrak{a}}}      % 原文第 6 节出现的记号
\newcommand{\wh}{\widehat}
\newcommand{\hbp}{\hat{\mathfrak{b}}}

% 原文用 \S 表示花体 S（class S theories）；此处重新定义以免与 \S 段落符冲突
\renewcommand{\S}{\mathcal{S}}

%---------------------------------------------------------------------
% 七、批注与要点环境（学术审校用）
%---------------------------------------------------------------------
\newcommand{\TutorNote}[1]{\par\smallskip\noindent
  \colorbox{yellow!18}{\parbox{0.97\textwidth}{\small\textbf{[审校批注]}\ #1}}\par\smallskip}

\newenvironment{keytakeaways}
  {\par\medskip\noindent
   \begin{minipage}{0.97\textwidth}
   \color{blue!45!black}\rule{\linewidth}{0.8pt}\par\smallskip
   \noindent\textbf{本节重点（Key Takeaways）}\par\smallskip
   \small}
  {\par\smallskip\noindent\color{blue!45!black}\rule{\linewidth}{0.8pt}\par\end{minipage}\medskip}

%=====================================================================
% 导言区结束
%=====================================================================


\begin{document}

%---------------------------------------------------------------------
% 标题页
%---------------------------------------------------------------------
\title{\bfseries 超对称场论与几何 Langlands 对应：\\[2pt] 硬币的另一面}

\author{Aswin Balasubramanian,\quad J\"org Teschner\\[4pt]
  \normalsize 汉堡大学数学系\\[2pt]
  \normalsize 及：DESY 理论组}

\date{arXiv:1702.06499v2\quad（2017 年 3 月 10 日）\\[4pt]
      \small 中文精译与学术审校版\\[2pt]
      \small 译自 String-Math 2016 会议文集}

\maketitle

%---------------------------------------------------------------------
% 摘要
%---------------------------------------------------------------------
\begin{center}
  {\bfseries 摘\quad 要}
\end{center}

\begin{quote}
\small
本短文报告了以下三方面工作之间关系的研究结果：Beilinson 与 Drinfeld 基于共形场论建立的几何 Langlands 对应、
Kapustin 与 Witten 基于 $N=4$ 超 Yang--Mills 理论给出的途径，以及 AGT 对应。本文把几何 Langlands 对应描述为：
在带面算符（surface operator）情形下，AGT 对应的一个推广在 Nekrasov--Shatashvili 极限下所取的形式。
沿着 Kapustin--Witten 与 Nekrasov--Witten 的进路，我们用在以 Hitchin 模空间为目标流形的二维 sigma 模型中的
有效描述，来诠释所得图像中的若干方面。
\end{quote}

\medskip
\noindent\rule{\textwidth}{0.6pt}

\begin{quote}
\small
\textbf{译者说明：}本译文以 arXiv:1702.06499v2 的 \LaTeX{} 源码与 PDF 版为准，逐段精译。
数学公式逐字保留，未改动任何符号、指标与指数；原文笔误与省略处均在 \verb|% [Tutor Note: ...]| 批注中指出，
不影响编译，阅读源码时可见。各节开头附有“本节重点”，为译者所加，非原文内容。
\end{quote}

\noindent\rule{\textwidth}{0.6pt}

\tableofcontents
\newpage

%=====================================================================
% 正文各节（由 build/sections/ 目录下的分节文件汇总而成）
%=====================================================================
%=====================================================================
% 第 1 节  引言
% 原文：Section 1. Introduction（arXiv:1702.06499v2, p. 1–2）
%=====================================================================
\section{引言}

\begin{keytakeaways}
\begin{itemize}
  \item \textbf{物理图像：}超对称规范场论与共形场论（CFT）之间存在一系列引人注目的显式联系。Alday--Gaiotto--Tachikawa 发现了其中最为显式的一类（AGT 对应），它与 Nekrasov 及其合作者命名为 BPS/CFT 对应的另一族结果相互交叠；更早地，Kapustin 与 Witten 揭示了 $N=4$ 超 Yang--Mills 理论与几何 Langlands 对应之间的联系，这一发展也自然地融入正在成形的图景。
  \item \textbf{数学工具：}通向几何 Langlands 对应的两条构造进路——Beilinson--Drinfeld 基于共形场论的进路与 Kapustin--Witten 基于四维超对称规范场论的进路；正文还将自由使用关于 Hitchin 模空间的标准定义与结果。
  \item \textbf{关键结论：}上述三条线索（Beilinson--Drinfeld 进路、Kapustin--Witten 进路、AGT 对应）之间的关系在许多方面仍不完整；本文所报告的结果为澄清这些关系带来若干启示，且限于底李代数为 $\mathfrak{sl}_2$ 的情形，更多细节与高秩情形留待后续系列论文。
  \item \textbf{阅读约定：}正文自由使用关于 Hitchin 模空间的标准定义与结果，其简要概述以术语表形式收于附录 \ref{Gloss}；术语首次出现时以上标标记，如 term$^{g)}$。
\end{itemize}
\end{keytakeaways}

\bigskip
\noindent\rule{\textwidth}{0.6pt}
\bigskip

过去几年中，人们发现了超对称规范场论与共形场论（conformal field theory, CFT）之间若干引人注目的联系。其中最为显式的一类由 Alday、Gaiotto 与 Tachikawa 发现\cite{AGT}，如今常被称为 AGT 对应。它与另一族结果相互交叠，Nekrasov 及其合作者将后者命名为 BPS/CFT 对应；关于这一主题的系列论文中的第一篇，以及这一方向先前工作的文献指引，见 \cite{N15}。一个更早的发展——Kapustin 与 Witten 所揭示的 $N=4$ 超 Yang--Mills 理论（SYM）与几何 Langlands 对应之间的关系\cite{KW}——则自然地融入这一正在成形的图景\cite{NW,NRS}。不过可以指出，Beilinson 与 Drinfeld 处理几何 Langlands 对应的进路与 CFT 有密切关联（综述及进一步文献见 \cite{Fr07}），而这些关联在 Kapustin 与 Witten 的进路中并不显然。上述主题之间的若干联系已在 \cite{NW,T10,NRS,GW} 中有所讨论；近期工作 \cite{KP,AFO} 表明，至少其中一部分关系容许进一步的形变，其动因来自超对称规范场论与弦论。然而，这一图景在许多方面似乎仍不完整。

本文是为 String-Math 2016 会议文集准备的，报告了一些结果，它们为澄清 Beilinson 与 Drinfeld 的进路、Kapustin 与 Witten 的进路以及 AGT 对应这三者之间的关系带来若干启示。为使本文篇幅保持在合理范围内，我们将主要限于讨论底李代数为 $\mathfrak{sl}_2$ 的情形。这里所述的结果是作者与 Ioana Coman-Lohi 共同推进的一个更大项目的一部分；一系列论文正在准备之中，将包含更多细节、对数学方面的扩展讨论以及对高秩情形的讨论。

在正文中，我们将自由使用若干关于 Hitchin 模空间的标准定义与结果。附录 \ref{Gloss} 以术语表的形式汇集了一份非常简短的概述。若某一术语在术语表中有对应条目，则该术语首次出现时将带上标，如 term$^{g)}$。

\bigskip
\noindent\rule{\textwidth}{0.6pt}

% [Tutor Note: 记号提醒：原文此处写 $N=4$（非 $\mathcal N=4$），而第 2.2 节又出现 $\mathcal N=(4,4)$，两种写法并存是原文自身的不统一；译文逐字照录、未作统一，引用时不要“顺手”改成 $\mathcal N=4$。另按术语表，“super-Yang--Mills theory (SYM)”统一作“超 Yang--Mills 理论”，连字符是 en-dash（--），不是单个“-”。]
% [Tutor Note: 含糊点拨：原文只说“certain $N=2$ supersymmetric gauge theories”，未指明具体理论。据 \cite{AGT} 的原始陈述（译者补充，非原文内容）：$SU(2)$、$N_f=4$ 的 $\mathcal N=2$ 超共形 QCD 在 $\Om$-背景（参数 $\ep_1,\ep_2$）下的瞬子配分函数，等于球面上插入四个 Liouville 顶点算符的 Virasoro 共形块；中心荷 $c=1+6(b+b^{-1})^2$，其中 $b^2=\ep_1/\ep_2$（依约定）。Nekrasov--Witten 对此的场论解释见本文第 2.3 节。]
% [Tutor Note: 跳步点拨：“Beilinson--Drinfeld 的进路与 CFT 有密切关联，而这些关联在 Kapustin--Witten 的进路中并不显然”——被省略的机制是：BD 构造把几何 Langlands 对应中的 $\CD$-模实现为临界水平 $k=-h^\vee$ 处仿射 Kac--Moody 代数 $\hat{\mathfrak g}_k$ 的共形块，Ward 恒等式赋予共形块层以 $\CD$-模结构；该构造之所以可能，关键在于 Feigin--Frenkel 定理给出 $\CU(\hat{\mathfrak g}_{-h^\vee})$ 的中心同构于形式圆盘上 ${}^L\mathfrak g$-opers 的空间，从而本征值 $E_i$ 可用来参数化 opers。详见本文第 2.1 节及 \cite{Fr07,FF92}。]
% [Tutor Note: 文献提示：\cite{KP,AFO} 所说的“further deformation”具体指：Kimura--Pestun 的 quiver $W$-代数（把 BPS/CFT 一侧的 $W$-代数结构推广到 $\Om$-形变情形）与 Aganagic--Frenkel--Okounkov 的量子 $q$-Langlands 对应（把几何 Langlands 对应沿 $q$-形变方向推广）。原文正文未展开，此处系译者据参考文献条目补出，供核对。]
% [Tutor Note: 本节小结——引言只交代动机与范围，技术内容自第 2 节开始：2.1 为 BD/CFT 进路，2.2 为 KW/规范场论进路，2.3 为 NW/AGT，2.4 指出前两条进路的共同根源（六维 $(2,0)$ 理论在 $\CM^6=\CM^4\times C$ 上的两种紧化，但所用的拓扑扭转不同），2.5 提及 Betti 几何 Langlands 与 Cordova--Jafferis 两条旁支。另：附录 \ref{Gloss} 的 \label{Gloss} 定义在 sec8.tex 中，单独编译本节时该 \ref 会显示为 ??，须整篇编译方可解析，属预期行为。]

%=====================================================================
% 第 2 节  综述
% 原文：Section 2. Review（arXiv:1702.06499v2；含 5 个子节，页码从略）
%=====================================================================
\section{综述}

\begin{keytakeaways}
\begin{itemize}
  \item \textbf{物理图像：}本节梳理了通向几何 Langlands 对应的三条进路。Beilinson--Drinfeld 的进路以共形场论为出发点，把对应实现为临界水平仿射李代数的共形块；Kapustin--Witten 的进路把 $N=4$ SYM 理论放在 $\Sigma\times C$ 上，沿 $C$ 紧化得到以 Hitchin 模空间为目标流形的二维 sigma 模型，并把 S-对偶实现为 SYZ 镜像对称；Nekrasov--Witten 的进路用 $\Om$-形变下的 $\CS$ 类理论，把 Virasoro 共形块实现为典范余迷向膜与 oper 膜之间的开弦态空间。
  \item \textbf{数学工具：}$\CD$-模与局部系（含 opers）、Hitchin 模空间及其环面纤维化、SYZ 镜像对称与膜的范畴、Hecke 函子与 Hecke 本征值性质、Virasoro 共形块与 Verlinde 线算符。
  \item \textbf{关键结论：}三条进路有共同的根源——六维 $(2,0)$ 理论在 $\CM^6=\CM^4\times C$ 上的紧化——但两种约化所用的拓扑扭转不同；这一差别既使结果的比较变得微妙，又恰好是把不同描述联系起来的杠杆（Gaiotto--Witten 对边界条件的分类、Betti 几何 Langlands、Cordova--Jafferis 的复 Chern--Simons 描述）。
  \item \textbf{记号提醒：}本节建立了贯穿全文的记号 $\Bun_G$、${}^{L}G$、$\CM_H(G)$、opers 与 $\Om$-形变参数 $\ep_1,\ep_2$（及 $\hbar=\ep_1/\ep_2$）；注意 Kapustin--Witten 与 Nekrasov--Witten 两处的 $\mathsf{B}_{\rm cc}$ 相似但不相同（见 §2.3 脚注），且镜像对称在交换 $G\leftrightarrow{}^{L}G$ 的同时把 $\hbar$ 换成 $1/\hbar$。
\end{itemize}
\end{keytakeaways}

\bigskip
\noindent\rule{\textwidth}{0.6pt}
\bigskip

\subsection{什么是几何 Langlands 对应？}

几何 Langlands 对应通常被宽泛地表述为如下对应：为 Riemann 面 $C$ 上的 ${}^{L}G$-局部系\glos 指派 $\Bun_G$ 上的 $\CD$-模。这里 ${}^{L}G$ 是单复李群 $G$ 的 Langlands 对偶群。这一对应中出现的 ${}^{L}G$-局部系，可以用由全纯 ${}^{L}G$-丛 $\mathcal{E}$ 与全纯联络 $\nabla'$ 组成的二元组 $(\mathcal{E},\nabla')$ 来表示；等价地，也可以用由 $(\mathcal{E},\nabla')$ 的和乐所确定的基本群 $\pi_1(C)$ 的表示 $\rho$ 来表示。我们将主要关注\emph{不可约}的 ${}^{L}G$-局部系的情形。相应地，$\Bun_G$ 上的 $\CD$-模可以更具体地描述为偏微分方程组，其形式为本征值方程 $D_i f= E_if$，其中 $D_i$ 是 $\Bun_G$ 上一族由 Hitchin 可积系统的 Hamilton 量量子化而得的微分算符。一个更具雄心的几何 Langlands 对应版本已在 \cite{AG} 中给出，在其中有必要把对应推广到若干类可约局部系。

几何 Langlands 对应的一些原始进路，从 ${}^{L}G$-局部系为\emph{oper}$^{g)}$ 的情形出发；所谓 oper，是二元组 $(\mathcal{E},\nabla')$，其中 $\nabla'$ 在规范等价意义下可取某种标准形式。全部 oper 构成的空间，在全部局部系的模空间中形成一个 Lagrange 子空间。Beilinson 与 Drinfeld 基于共形场论的进路，对每个 oper 构造 $\Bun_G$ 上 $\CD$-模范畴中的一个对象，即仿射李代数 $\hat{\mathfrak g}_k$ 在临界水平 $k=-h^\vee$ 处的共形块。刻画共形块的 Ward 恒等式，赋予共形块层以 $\CD$-模结构。在 $k=-h^\vee$ 处，泛包络代数 $\mathcal{U}(\hat{\mathfrak g}_{k})$ 具有一个很大的中心，它同构于形式圆盘上 ${}^{L}\mathfrak{g}$-oper 的空间\cite{FF92}。由此可以证明：来自 Ward 恒等式的 $\CD$-模结构，可以用量子化 Hitchin Hamilton 量的本征值方程组 $D_i f= E_if$ 来描述，而本征值 $E_i$ 参数化 opers 的选取\cite{Fr07}。

在 \cite[Section 9.6]{Fr07} 中描述了 Beilinson--Drinfeld 几何 Langlands 对应构造的一个扩充：由 oper 的情形推广到一般的不可约局部系。其依据是这样一个事实：这类局部系总在规范等价意义下化为带有某些额外奇点的 oper\cite{Ar16}。Beilinson 与 Drinfeld 的构造对这样的 oper 给出 $\hat{\mathfrak g}_{-h^{\vee}}$ 的共形块，其中在额外奇点处插入了由 $\mathfrak{g}$ 的有限维表示诱导出的某些退化表示。在这个意义上，我们可以把一般不可约局部系的几何 Langlands 对应，看作对通常（非奇异）oper 而言已存在的对应的一种扩充。从本文将要提出的视角来看，这一观点将被证明是自然的。还应指出，在 \cite{Ga15} 为证明 \cite{AG} 中表述的强化几何 Langlands 猜想所给出的纲要中，Beilinson--Drinfeld 构造起着重要作用。

为便于后文引用，我们指出：Beilinson--Drinfeld 构造把几何 Langlands 对应中出现的 $\CD$-模描述为共形块的空间；这些空间一方面自然地纤维化于 $\Bun_G$ 之上，另一方面又纤维化于不可约 ${}^{L}G$-局部系的空间之上。

\subsection{几何 Langlands 对应——Kapustin--Witten 的进路}

Kapustin 与 Witten \cite{KW} 提出的几何 Langlands 对应版本，基于在形如 $\Sigma\times C$ 的四维流形上考察规范群为 $G_c$（$G$ 的紧实形式）的 $N=4$ SYM 理论，其中 $C$ 是 Riemann 面。在 $C$ 上紧化，可将 $N=4$ SYM 的拓扑扭曲版本有效地表示为以 $\Sigma$ 上的 Hitchin 模空间$^{(g)}$ $M_H(G)$ 为目标流形的 $2d$ 拓扑扭曲 sigma 模型。Hitchin 模空间是完全可积的，这一点在把 $M_H(G)$ 描述为环面纤维化时表现得十分明显；正是这种完全可积性，使我们能把 $N=4$ SYM 理论 S-对偶的推论，描述为联系目标流形分别为 $M_H(G)$ 与 $M_H({}^{L} G)$ 的两个 $2d$ sigma 模型的 SYZ 镜像对称。

为把上述内容与几何 Langlands 对应联系起来，Kapustin 与 Witten 考虑 $\Sigma=R\times I$、$I=[0,\pi]$ 的情形。$N=4$ SYM 理论的超对称边界条件在 $C$ 上紧化之后，将定义以 $\CM_H(G)$ 为目标流形的 $\mathcal{N}=(4,4)$ sigma 模型在 $\Sigma$ 上的边界条件。设 $\mathfrak{B}$ 为如下范畴：其对象是称为膜（brane）、保持最大数目超对称的边界条件 $\mathsf{B}$；其态射是 $\BR\times I$ 带状区域上 sigma 模型的“开弦”态空间 ${\rm Hom}_{\CM_H(G)}^{}(\mathsf{B}_1,\mathsf{B}_2)$，其中边界 $\BR\times \{0\}$ 与 $\BR\times \{\pi\}$ 分别配以边界条件 $\mathsf{B}_1$ 与 $\mathsf{B}_2$。

所谓典范余迷向膜（canonical coisotropic brane）$\mathsf{B}_{\rm cc}$ 起着特殊的作用\cite{KW,NW}。向量空间 $\mathcal{A}_{\rm cc}={\rm Hom}_{\CM_H(G)}^{}(\mathsf{B}_{\rm cc},\mathsf{B}_{\rm cc})$ 具有自然的代数结构，其乘法由“拼接开弦”定义。空间 $\CH(\mathsf{B})={\rm Hom}_{\CM_H(G)}^{}(\mathsf{B}_{\rm cc},\mathsf{B})$ 是代数 $\mathcal{A}_{\rm cc}$ 上的左模，其作用通过在带状区域 $I$ 的左边界上由 $\mathcal{A}_{\rm cc}$“拼接开弦”来定义。Kapustin 与 Witten 论证：代数 $\mathcal{A}_{\rm cc}$ 包含 $\Bun_G$ 上整体微分算符构成的代数。由此可知，空间 $\CH(\mathsf{B})$ 表示 $\Bun_G$ 上的 $\CD$-模。

把支撑在 $\BR\times\{x\}\times P$（$x\in I$，$P\in C$）上的 Wilson 与 't Hooft 线算符约化到二维 TQFT，就在膜的范畴上定义了自然的函子，它们诱导出空间 $\CH(\mathsf{B})$ 的修改。这样定义的函子在 \cite{KW} 中被等同于几何 Langlands 对应中的 Hecke 函子。对于某些膜 $\mathsf{B}$，在每个固定的 $P\in C$ 处，可将所得的修改表示为 $\CH(\mathsf{B})$ 与 ${}^{L}G$ 的某个有限维表示 $V$ 的张量积。若让点 $P\in C$ 变化所得到的一族修改能粘接成一个局部系，则称膜 $\mathsf{B}$ 满足 Hecke 本征值性质。

\cite{KW} 中指出了一族具有该性质的膜 $\mathsf{F}_\mu$。这些膜 $\mathsf{F}_\mu$ 支撑在 Hitchin 环面纤维化的纤维上。$N=4$ SYM 的 S-对偶在以 $\CM_H(G)$ 为目标流形的 sigma 模型中，表现为 SYZ 镜像对称的一个变体：它把膜 $\mathsf{F}_\mu$ 与对偶 sigma 模型（目标流形为 $\CM_H({}^{L}G)$）中的膜联系起来，后者由在 $\mu\in\CM_H({}^{L}G)$ 处具有点状支撑的摩天楼层（skyscraper sheaf）$\check{\mathsf{F}}_\mu$ 表示。

\subsection{AGT 对应——Nekrasov--Witten 的进路}

Alday、Gaiotto 与 Tachikawa 发现了某些 $N=2$ 超对称规范场论的瞬子配分函数与 Virasoro 代数共形块之间的关系\cite{AGT}。这一发现激发了大量工作，尤其引出了这类关系的各种推广。为解释 \cite{AGT} 中发现的关系，Nekrasov 与 Witten 考虑了 $\mathcal{S}$ 类的四维 $N=2$ 超对称规范场论；这类理论由极大超对称的六维量子场论在形如 $\CM^4\times C$ 的流形上沿 Riemann 面 $C$ 紧化得到。在与李代数 $\mathfrak{g}=\mathfrak{sl}_2$ 相关联的情形，所得理论具有弱耦合的 Lagrangian 描述，这些描述与 $C$ 的一种裤子分解 $\si$ 的选取相联系。对于可局部表示为纤维积 $\mathbb{R}\times I\times S^1\times S^1$ 的四维流形 $\CM^4$，\cite{NW} 论证：(i) 可以定义参数为 $\ep_1,\ep_2$ 的 $\Omega$-形变；(ii) 通过在 $S^1\times S^1$ 上紧化，可以得到用 $R\times I$ 上以 $\CM_H(G)$ 为目标流形的 sigma 模型给出的有效描述。该 sigma 模型的耦合参数为 $\ep_1/\ep_2$。

在表示 $\CM^4\simeq\mathbb{R}\times I\times S^1\times S^1$ 中，区间 $I$ 的两个端点对应于 $\CM^4$ 完全正则的点。因此，在描述 $\mathcal{S}$ 类理论在 $\CM^4$ 上紧化的、定义在 $R\times I$ 上并以 $\CM_H(G)$ 为目标流形的 sigma 模型中，必须配备特殊的边界条件。当紧化给出以 $\CM_H(G)$ 为目标流形的 sigma 模型时，\cite{NW} 论证：相应的边界条件分别由位于 $\BR\times\{0\}$ 处的典范余迷向膜的一个变体 $\mathsf{B}_{\rm cc}$，以及一类称为“oper 膜”（brane of opers）的新膜（此处记作 $\mathsf{B}_{\rm op}$）描述。\footnote{此语境下记作 $\mathsf{B}_{\rm cc}$ 的膜，与 \cite{KW} 中所考虑的膜相似但并不相同。文献 \cite{NW} 用记号 $B_{N'}$ 表示此处记作 $\mathsf{B}_{\rm op}$ 的膜。}膜 $\mathsf{B}_{\rm op}$ 是 $\mathsf{B}_{\rm cc}$ 的镜像对偶；\cite{NW} 提出，$\mathsf{B}_{\rm op}$ 是一张 Lagrange 膜，支撑在 $M_H(G)$ 内的 oper 簇上。

\cite{NW} 进一步提出，空间 $\CH={\rm Hom}_{\CM_H(G)}^{}(\mathsf{B}_{\rm cc},\mathsf{B}_{\rm op})$ 可以等同于 Virasoro 共形块的空间。为说明这一等同的动机，Nekrasov 与 Witten 指出：当 $\hbar=\ep_1/\ep_2$ 时，代数 $\mathcal{A}_{\rm cc}^{\hbar}(G)={\rm Hom}_{\CM_H(G)}^{}(\mathsf{B}_{\rm cc},\mathsf{B}_{\rm cc})$ 同构于作用在 Virasoro 共形块空间上的 Verlinde 线算符代数。镜像对称给出 $\CH(G)\simeq\CH({}^{\rm L}G)$ 的对偶描述，即空间 ${\rm Hom}_{\CM_H({}^{L}G)}^{}({\mathsf{B}}_{\rm op}',{\mathsf{B}}_{\rm cc}')$，其中 ${\mathsf{B}}_{\rm op}'$ 与 ${\mathsf{B}}_{\rm cc}'$ 分别是 ${\mathsf{B}}_{\rm op}$ 与 ${\mathsf{B}}_{\rm cc}$ 的近亲，其超对称不变性质有所改变。在对偶描述中，代数 $\check{\mathcal{A}}_{\rm cc}^{1/\hbar}({}^{L}G)={\rm Hom}_{\CM_H({}^{L}G)}^{}(\mathsf{B}_{\rm cc}',\mathsf{B}_{\rm cc}')$ 有一个显然的右作用，其作用通过在带状区域 $I$ 的右边界上“拼接开弦”来定义。$\mathcal{A}_{\rm cc}^{\hbar}(G)$ 与 $\check{\mathcal{A}}_{\rm cc}^{1/\hbar}({}^{L}G)$ 之间存在（几乎）交换的作用，这是 Virasoro 共形块空间的一个特征性质。

\subsection{反方向的联系}

Nekrasov 与 Witten 的工作\cite{NW} 与 Kapustin 与 Witten 的进路\cite{KW} 有许多共同要素，这并非偶然。一个共同的根源在于：$\CS$ 类理论与 $N=4$ SYM\cite{Wi09} 都可以看作六维 $(2,0)$ 理论在六维流形 $\CM^6=\CM^4\times C$ 上紧化的结果，其中 $C$ 是 Riemann 面，而 $\CM^4$ 是一个可局部表示为形如 $\mathbb{R}\times I \times S^1\times S^1$ 的圆纤维化的四维流形。在 $C$ 上紧化给出 $\CS$ 类理论\cite{GMN2}；在 $S^1\times S^1$ 上紧化则给出 $\mathbb{R}\times I \times C$ 上的 $N=4$ SYM，这正是 \cite{KW} 所考察的框架，\cite{Wi09} 对此作了进一步讨论。

不过应当注意，两种紧化使用的是不同的拓扑扭转，这使得结果的比较有些微妙。尽管如此，这一事实仍可用来把由 $\CS$ 类理论紧化所得的、目标流形为 $\CM_H$ 的二维 sigma 模型中的超对称边界条件，与 $C$ 上 $N=4$ SYM 中的边界条件联系起来。这些边界条件已在 Gaiotto 与 Witten 的工作\cite{GW1} 中分类。由此可以建立典范余迷向膜与 $N=4$ SYM 中纯 Neumann 边界条件之间的关系\cite{KW,GW}。交换 $S^1\times S^1$ 中的两个圆，与 $N=4$ SYM 的 S-对偶相关，而 S-对偶蕴含了 \cite{GW2} 所研究的该理论边界条件之间的关系。这使 \cite{GW} 把膜 $\mathsf{B}_{\rm op}$（即 \cite{NW} 中典范余迷向膜的镜像对偶）与 $N=4$ SYM 中由所谓 Nahm 极点边界条件下降所得的边界条件联系起来，如下文第 \ref{S:Nahm} 小节将简要讨论的。

\subsection{其他进路}

如前所述，在上面考虑的两个从六维到四维的约化中，需要使用不同的扭转。例如，要得到 \cite{NW} 所研究的框架，需要把 $\CM^4\times C$ 上的 $(2,0)$ 理论扭转成在 $\CM^4$ 上拓扑的理论；而要由六维 $(2,0)$ 理论得到 \cite{KW} 所研究的 $\BR\times I\times C$ 上拓扑扭转的 $N=4$ SYM，显然需要另一种扭转。

为描述 $\BR\times I\times C$ 上拓扑扭转的 $N=4$ SYM 的维数约化，人们自然会考虑纯拓扑的边界条件，即不依赖于 $C$ 上复结构的边界条件。这一观点促使 Ben-Zvi 与 Nadler \cite{BN16} 提出 Betti 几何 Langlands 猜想；它是 \cite{AG} 中所表述的几何 Langlands 对应的一个纯拓扑变体，自然地适用于 \cite{KW} 所研究的四维 TQFT。

Cordova 与 Jafferis 在 \cite{CJ} 中提出了理解 AGT 对应的又一个进路。它从形如 $\CM^4\times C$ 的六维流形出发，其中 $\CM^4$ 是一个压扁的四维球（squashed four-sphere），可表示为区间 $I_0$ 上的三维球纤维化，局部地 $\CM^4\sim I_0 \times S^3_{\ep_1\ep_2}$，此时 $\ep_1,\ep_2$ 是压扁参数。在 $C$ 上作部分拓扑扭转，可以论证：带有 $C$ 上扭转的 $(2,0)$ 理论的 Weyl 不变性导致配分函数既不依赖于 $\CM^4$ 的整体半径，也不依赖于 $C$ 的体积。当 $C$ 的体积很小时，有效描述将由 Pestun \cite{Pe} 及后续大量文献所研究的（压扁）四维球上 $S$ 类理论的配分函数给出。可以论证：若先沿表示 $\CM^4\sim I_0 \times S^3_{\ep_1\ep_2}$ 中出现的三维球 $S^3_{\ep_1\ep_2}$ 紧化，就能得到用 $C$ 上 CFT 表示的等价描述。以五维 SYM 理论作为中间步骤，Cordova 与 Jafferis 找到了用 $I_0\times C$ 上复 Chern--Simons 理论给出的有效描述。$C$ 两端处的边界条件由 \cite{GW1,GW2,GW} 所研究的 Nahm 极点边界条件确定。于是可以援引 Chern--Simons 理论与二维共形场论之间通常对应的推广版本，其中 Nahm 极点边界条件转化为把 WZNW 模型约化为共形 Toda 理论的约束条件。

\bigskip
\noindent\rule{\textwidth}{0.6pt}

% [Tutor Note: 易错提醒（§2.1）：oper 是二元组 $(\mathcal{E},\nabla')$，其中 $\nabla'$ 在规范等价意义下取固定的标准形；当 ${}^{L}G=SL_2$ 时，局部坐标 $z$ 下 $\nabla'=d+\begin{pmatrix}0&1\\ q&0\end{pmatrix}dz$，$q$ 为二次微分。据此可核对原文“opers 构成 Lagrange 子空间”的断言：$\dim\mathrm{Op}_{SL_2}=\dim H^0(C,K_C^2)=3g-3$，而全部 $SL_2$ 局部系的模空间维数为 $6g-6$，恰为其一半；Lagrange 性还依赖局部系模空间上的 Atiyah--Bott 辛形式。原文未给证明，完整的 oper 理论见 Drinfeld--Simpson。]
% [Tutor Note: 跳步点拨（§2.1）：$k=-h^\vee$ 是临界水平，$h^\vee$ 为对偶 Coxeter 数（$\mathfrak{sl}_2$ 时 $h^\vee=2$）。此时 Sugawara 中心荷 $c=k\dim\mathfrak g/(k+h^\vee)$ 的分母为零、Virasoro 场本身不存在，故不能按通常 CFT 方式定义共形块；Beilinson--Drinfeld 的出路正是原文所说的“很大的中心”：Feigin--Frenkel 定理给出 $\mathcal U(\hat{\mathfrak g}_{-h^\vee})$ 的中心同构于形式圆盘上 ${}^L\mathfrak g$-opers 的函数环\cite{FF92}，由此把共形块组织成 $\CD$-模，并使本征值 $E_i$ 参数化 opers\cite{Fr07}。原文 “This can be used to show” 略去的正是这一步。]
% [Tutor Note: 易错提醒（§2.2）：SYZ 镜像对称要求纤维是 special Lagrangian 环面，而不只是复环面。Hitchin 纤维化的纤维关于一个复结构 $I$ 是全纯环面，关于另一个复结构（如超 Kähler 旋转后的 $J$）才是 special Lagrangian；正是 Hitchin 模空间的超 Kähler 结构（三个复结构 $I,J,K$）保证同一纤维化可以同时充当 SYZ 纤维化。忽略超 Kähler 旋转，就会误以为“环面纤维化 + 对偶环面”已足以给出镜像对称。]
% [Tutor Note: 记号易错（§2.3）：脚注已提醒，Nekrasov--Witten 语境下的 $\mathsf{B}_{cc}$ 与 Kapustin--Witten 的 $\mathsf{B}_{cc}$ 只是“相似但不相同”，NW 原文的 $B_{N'}$ 才是本文的 $\mathsf{B}_{op}$，切勿把两篇文章的 $\mathsf{B}_{cc}$ 当成同一张膜。另注意镜像对称在交换 $G\leftrightarrow{}^LG$ 的同时把形变参数 $\hbar=\ep_1/\ep_2$ 换成 $1/\hbar$：这正是对偶代数写作 $\check{\mathcal A}_{cc}^{1/\hbar}({}^LG)$ 的原因，也是两套（几乎）交换作用能同时出现的原因。]
% [Tutor Note: 含糊点拨（§2.4）：原文说两种紧化“用了不同的拓扑扭转，使比较有些微妙”，被略去的是：六维 $(2,0)$ 理论在不同紧化中选取不同的超荷子代数作 BRST 荷——一种使 $\CM^4$ 方向拓扑化，另一种在 $\BR\times I\times C$ 上给出 Kapustin--Witten 的 GL-扭转 $N=4$ SYM——两者不能同时实现。因此比较必须借道 Gaiotto--Witten 对 $N=4$ SYM 边界条件的分类\cite{GW1,GW2}，并用 $S^1\times S^1$ 两圆交换（对应 S-对偶）把两边的边界条件对应起来；原文的“This led \cite{GW} to…”正是这一链条的终点。]
% [Tutor Note: 跳步点拨（§2.5）：原文说“Nahm 极点边界条件转化为把 WZNW 模型约化为共形 Toda 理论的约束”，这实际上就是量子 Drinfeld--Sokolov 约化：取主幂零元（在边界条件上即 Nahm 极点），把 $\hat{\mathfrak g}$-WZNW 约化为 Toda。Ben-Zvi--Nadler 的 Betti 版则是把 de Rham 侧（$\CD$-模）换成 Betti 侧（可构造层），故称“纯拓扑变体”。本节小结：BD/CFT、KW/规范场论、NW/AGT 三条进路在 §2 中只是并列陈述，真正的桥接自 §2.4 的六维共同根源与“不同扭转”开始；本文的提议（带面算符情形下 AGT 对应的推广取 Nekrasov--Shatashvili 极限，即几何 Langlands 对应）建立在 §3 及以后。]

%=====================================================================
% 第 3 节  余二维缺陷与面算符
% 原文：Section 3. Defects of co-dimension two and surface operators
%       （arXiv:1702.06499v2）
%=====================================================================
\section{余二维缺陷与面算符}\label{codim2}

\begin{keytakeaways}
\begin{itemize}
  \item \textbf{物理图像：}把六维 $(2,0)$ 理论看作 M 理论中一叠 M5-膜上的有效理论之后，再加入一叠与原 M5-膜共享部分支撑的 M5'-膜，就在 $\CS$ 类理论中得到余二维面算符（面算符即余二维缺陷）。在二维 sigma 模型中，这类面算符由“零 Nahm 极点”边界条件——即固定 $C$ 上的全纯丛 $\CE$——来描述，对应的膜是 Hitchin 第二纤维化的纤维。
  \item \textbf{数学工具：}Gaiotto--Witten 对 1/2 BPS 边界条件的分类、Nahm 极点渐近行为、非阿贝尔 Hodge（NAH）对应、Hitchin 第二纤维化、仿射李代数共形块及其上的 $\mathcal{A}_{dR}^{\hbar}$/$\mathcal{A}_{B}^{\hbar}$ 模结构、$\bar\partial$-上同调，以及 Virasoro 单值化与 Narasimhan--Seshadri 定理的类比。
  \item \textbf{关键结论：}由固定 $\CE$ 定义的 Lagrangian 膜所给出的空间 $\CH^{\sst (2)}_x$，应与“缓增”（tempered）共形块 $\mathcal{CB}_{\rm te}(C,\hfg_k,\mathcal{E}_x)$ 同构，而非与通常代数定义的 $\mathcal{CB}_{\rm al}(C,\hfg_k,\mathcal{E}_x)$ 同构；后者太大。$\CH^{\sst (2)}_x$ 上存在两个模结构，是支持此猜想的第一条证据。
  \item \textbf{遗留问题：}$\mathcal{CB}_{\rm te}$ 的严格定义在本文中未能给出，作者声明留待另文；Virasoro 情形与 Kac--Moody 情形的讨论都只是“同构于共形块的子空间”的类比论证。
\end{itemize}
\end{keytakeaways}

\bigskip
\noindent\rule{\textwidth}{0.6pt}
\bigskip

M 理论提示了上述框架的一个自然推广。把六维 $(2,0)$ 理论视为 M 理论中一叠 M5-膜上的有效理论，则自然会引入一叠额外的 M5'-膜，使之与原 M5-膜共享一部分支撑（support），如下表所示：
\begin{center}
\begin{tabular}{l || c|c|c|c|c|c|c|c|c|c|c}
膜 & 0 & 1 & 2 & 3 & 4 & 5 & 6 & 7 & 8 & 9 & 10 \\ \hline\hline
M5  & x & x & x & x &   &   & x &   &   &   & x  \\
M5' & x & x &   &   &   &   & x &   & x & x & x \\
\end{tabular}
\end{center}
视所考虑的紧化类型不同，所得对象或是 $\CS$ 类理论中的面算符，或是 $N=4$ SYM 理论中的线算符。这将是第一种修正，而它后来被证明对本文的故事至关重要。

\subsection{$\CS$ 类理论中的面算符}

目前已积累了相当多的证据，表明把 M5-M5'-膜系统在 $C$ 上紧化所得的有效描述，正是被某一类面算符所修正的 $\CS$ 类理论。相关的面算符可通过规定规范场沿二维子流形 $\CM^2\subset\CM^4$ 的某种奇异性来刻画；当 $C$ 的亏格为 $g$、有 $n$ 个穿刺点时，这类奇异性由一组参数 $x=(x_1,\dots,x_{3g-3+n})$ 标定 \cite{AT,FGT}。在存在这类面算符时，可以定义瞬子配分函数的一个推广 $\CZ^{\rm\sst inst}_\si(a;x,\tau;\ep_1,\ep_2)$：除了依赖于标量零模的取值 $a$、与 $C$ 的复结构模 $\tau$ 相联系的紫外规范耦合、以及与 $\Om$-形变相联系的参数 $\ep_1,\ep_2$ 之外，它还额外地依赖于参数 $x$。下标 $\si$ 指决定所用 Lagrangian 表示的裤分解（pants decomposition）。

\cite{Br,AT,KPPW,Ne,Na} 中的显式计算为 AGT 对应的一个推广提供了证据：它把 $\CZ^{\rm\sst inst}_{\si}(a;x,\tau;\ep_1,\ep_2)$ 等同于与仿射李代数 $\hat{\mathfrak{g}}_k$ 的共形块相联系的“手征配分函数” $\CZ^{\rm\sst WZW}_{\si}(a;x,\tau;k)$，该代数的能级为 $k=-h^{\vee}-\ep_2/\ep_1$，而这些共形块由拼接构造（gluing construction）定义。在这一对应下，参数 $x$ 被等同于 $C$ 上全纯 $G$-丛的模空间 $\mathrm{Bun}_G$ 上的坐标。AGT 对应这一推广的证明，将可由 \cite{N15} 所宣布的结果推出，即 $\CZ^{\rm\sst inst}_\si$ 满足 KZB 方程。\footnote{
另见 {\cite{Ne}}，其中有这一方向上的一个较早结果。}

本节的目标是把 \cite{NW,GW} 的进路推广到带余二维面算符的 $\CS$ 类理论。我们将提出以 $\CM_H(G)$ 为目标流形的二维 sigma 模型作为其有效描述，其中 opers 膜 $\mathsf{B}_{\rm op}$ 被一族 Lagrangian 膜 $\mathsf{L}_{x}^{\sst (2)}$ 取代，后者支撑在 Hitchin 第二纤维化（Hitchin's second fibration）位于丛 $\mathcal{E}_x$ 之上的纤维上，而 $\mathcal{E}_x$ 由 $\mathrm{Bun}_G$ 的坐标 $x$ 标记。这些膜 $\mathsf{L}_{x}^{\sst (2)}$ 在此语境中的意义、以及它们与共形块的关系，最早由 E. Frenkel 在 \cite{Fr10} 中提出；这一点曾在 2012 年被指出给作者之一（J.T.）。下文的目标是为这一提议提供进一步的证据。

\subsection{Nahm 极点边界条件}\label{S:Nahm}

Gaiotto 与 Witten 在 \cite{GW1} 中用数据 $(\rho,H,T)$ 对 $N=4$ SYM 的 1/2 BPS 边界条件作了分类，其中 $\rho:\mathfrak{sl}_2\ra \mathfrak{g}$ 是 $\mathfrak{sl}_2$ 到规范群 $G_c$ 的李代数 $\mathfrak{g}$ 中的一个嵌入，$H$ 是 $\rho$ 的像在 $G_c$ 中的中心化子（commutant）的一个子群，而 $T$ 是一个具有 $N=4$ 超对称、且至少具有 $H$ 整体对称性的三维 SCFT。我们只需要其中最简单的两个边界条件。下面先简述 \cite{GW} 中所研究的所谓 Nahm 极点边界条件，它对应于三元组 $(\rho,\mathrm{Id},T_\emptyset)$，其中 $\rho$ 是一个\emph{主}（principal）$\mathfrak{sl}_2$-嵌入，$T_\emptyset$ 表示平凡的三维 SCFT。随后我们将讨论一个更简单的情形：把 $\rho$ 换成将 $\mathfrak{sl}_2$ 映到 $0\in\mathfrak{g}$ 的平凡嵌入；这一情形对我们来说特别重要。

对我们的目的而言，只需在刻画 $N=4$ SYM 中保持某些超对称的场构型的 BPS 方程 \cite{KW} 的解上描述 Nahm 极点边界条件。若把注意力限制在 $\mathbb{R}\times \mathbb{R}_+ \times C$ 上、沿 $\mathbb{R}$ 的平移不变的 BPS 方程解，则得到如下形式的微分方程组
\begin{subequations}\label{BPS}
\begin{align}
&[\,\CD_z\,,\,\CD_{\bar{z}}\,]=0\,,\quad [\,\CD_y\,,\,\CD_{{z}}\,]=0\,,
\quad [\,\CD_y\,,\,\CD_{\bar{z}}\,]=0\,,\label{Fterm}\\
&\sum_{i=1}^3[\,\CD_i\,,\,\CD_i^\dagger\,]=0\,. \label{Dterm}\,,
\end{align}
\end{subequations}
其中已使用记号 $z=x_2+ix_3$ 与 $y=x_1$，而微分算子 $\CD_i$ 取形式\footnote{我们的约定与 \cite{GW} 略有不同。}
\begin{equation}
\begin{aligned}
&\CD_z=\zeta\pa_z+\mathcal{A}_z\,,\\
&\CD_{\bar{z}}=\pa_{\bar{z}}+\mathcal{A}_{\bar{z}}\,,
\end{aligned}
\qquad
\begin{aligned}
&\mathcal{A}_z=\zeta A_z+\phi_z\,,\\
&\mathcal{A}_{\bar{z}}=A_{\bar{z}}+\zeta\phi_{\bar{z}}\,,
\end{aligned}
\qquad \CD_y=\pa_{y}+A_y-i\phi_y\,.
\end{equation}
参数 $\zeta$ 决定了哪些超对称被保持。\cite{GW} 提出：\rf{BPS} 的解空间模紧致规范变换，同构于“F-项”方程 \rf{Fterm} 的解空间模复规范变换所得的模空间。方程 $[\CD_z,\CD_{\bar{z}}]=0$ 在每个固定的 $y$ 处确定 $C$ 上的一个平坦复联络。\rf{Fterm} 中其余的方程则表明，这一平坦联络对 $y$ 的依赖由一个复规范变换来实现。

Nahm 极点型的边界条件在 \cite{GW} 中是这样定义的：要求 \rf{BPS} 的解具有如下形式的奇异性
\begin{equation}\label{Nahm}
\mathcal{A}_z\,\underset{y\ra 0}{\sim}\,\mathfrak{t}_- \,y^{-1}+\CO(y^0)\,,\qquad
\mathcal{A}_{\bar{z}}\,\underset{y\ra 0}{\sim}\,\CO(y^0)\,,\qquad
\mathcal{A}_1\,\underset{y\ra 0}{\sim}\,\mathfrak{t}_3 \,y^{-1}+\CO(y^0)\,,
\end{equation}
其中 $\mathfrak{t}_+=\mathfrak{t}_1+i\mathfrak{t}_2$，而 $\mathfrak{t}_i$（$i=1,2,3$）是 $\mathfrak{g}$ 的某个主 $\mathfrak{sl}_2$ 子代数的生成元。总可以通过一个规范变换把 $\mathcal{A}_{\bar{z}}$ 取为零，于是我们在每个 $y$ 处得到的 $C$ 上的平坦联络可表示为一个局部系统 $(\mathcal{E}_{y},\nabla_{y}')$，它由一个全纯丛和一个全纯联络 $\nabla_{y}'=dz(\pa_z+\mathcal{A}_z(z;y))$ 组成。在 $\mathfrak{g}=\mathfrak{sl}_2$ 的情形，我们可以把 \rf{Nahm} 的第一个条件改述为如下条件：存在一组截面基 $s=\{s_1,s_2\}$，使得 $\mathcal{A}$ 相对于它具有形式 $\mathcal{A}=g\tilde{\mathcal{A}}g^{-1}+gd g^{-1}$，其中
\begin{equation}\label{Nahm'}
\tilde{\mathcal{A}}_z\,\underset{y\ra 0}{\sim}\,\bigg(\begin{matrix} 0 & t\\ 1 & 0\end{matrix}\bigg)+\CO(y^1) \,,
\qquad g\underset{y\ra 0}{\sim} \bigg(\begin{matrix} y^{1/2}& 0 \\  0& y^{-1/2} \end{matrix}\bigg) +\CO(y^0).
\end{equation}
此时水平截面 $(d+\mathcal{A})s=0$ 的第一个分量 $s_1$ 将按 $y^{1/2}$ 的方式消失。正如 \cite{GW} 所解释的，这意味着 $C$ 上的局部系统 $\lim_{y\ra 0}^{}(\tilde{\mathcal{E}}_{y},\tilde{\nabla}_{y}')$ 是一个 oper。

Nahm 极点边界条件的特点是：由于主 $\mathfrak{sl}_2$-嵌入的中心化子是平凡的，它把 $G_c$ 极大地破缺。与之相反的另一个极端，则对应于平凡的 $\mathfrak{sl}_2$-嵌入：通过固定规范场 $\mathcal{A}_{\bar{z}}$ 的边界值，得到与三元组 $(0,\mathrm{Id},T_\emptyset)$ 相联系的类似边界条件。

在约化到二维时，在 $I$ 的边界上固定 $\mathcal{A}_{\bar{z}}$ 意味着在 $C$ 上固定了一个全纯丛，而复规范场的 $(1,0)$ 部分不受约束。偶对 $(\CE,\nabla_\zeta')$ 的模空间 $\CM_{dR}(G)$ 经由非阿贝尔 Hodge（NAH）对应$^{(g)}$ 同构于 Hitchin 模空间 $\CM_H(G)$。因此固定 $\CE$ 就在 $\CM_H(G)$ 中定义了一个子流形，它关于全纯辛形式 $\Omega_\zeta$ 是 Lagrangian 的，并且关于复结构 $I_\zeta$ 是全纯的。当 $\zeta=i$ 时有 $\Omega_\zeta=\Omega_J$、$I_\zeta=J$，于是由零 Nahm 极点边界条件约化所得的膜被认定为 A-模型中（辛结构取 $\omega_I$，即 \cite{NW} 中所用者）的一个 (A,B,A)-膜。

我们提出：零 Nahm 极点边界条件代表着余二维面算符的存在。的确，如 \cite{FGT} 所论证的，余二维面算符的存在会自然地引入额外的背景数据，而这些数据在几何上可以表示为 $C$ 上一个全纯丛的选取。

\subsection{与共形块的关系}

$\CM_H(G)$ 中由固定 $\CE$ 的偶对 $(\CE,\nabla')$ 所定义的 Lagrangian 子流形 $\mathsf{L}_{\CE,\zeta}^{\sst (2)}$，有时被称为 Hitchin 第二纤维化的纤维。取 $\mathrm{Bun}_G$ 的一组坐标 $x$，我们将用 $\mathsf{L}_{x}^{\sst (2)}$ 表示 $\mathsf{L}_{\CE_x,i}^{\sst (2)}$，其中 $\CE_x$ 是代表由坐标 $x$ 所指定的 $\mathrm{Bun}_G$ 中那一点的一个丛。

上面提到的、关于余二维面算符存在时 AGT 对应之推广的结果提示：空间 $\CH^{\sst (2)}_x={\rm Hom}_{\CM_H(G)}(\mathcal{B}_{\rm cc},\mathsf{L}_{x}^{\sst (2)})$ 可以等同于 $C$ 上能级为 $k=-h^{\vee}-\frac{\ep_2}{\ep_1}$ 的仿射李代数 $\hat{\mathfrak{g}}_k$ 的共形块空间。

然而必须指出：在眼下的无穷维情形中，要为 $\CH^{\sst (2)}_x$ 以及相应的共形块空间找到恰当的定义并非平凡。共形块通常的代数定义给出一个空间 $\mathcal{CB}_{\rm al}(C,\hfg_k,\mathcal{E}_x)$，它太大，对我们并不适用。我们需要考虑其中一个子空间，记作 $\mathcal{CB}_{\rm te}(C,\hfg_k,\mathcal{E}_x)$，即“缓增”（tempered）共形块。这两个空间之间的关系在某些方面类似于形式幂级数 $\sum_{n\in\BZ}a_nz^n$ 的空间 $\mathcal{K}_{\rm al}$ 与单位圆上缓增分布的空间 $\mathcal{K}_{\rm te}$ 之间的关系。后者可以在分布意义下用 Fourier 级数 $\sum_{n\in\BZ}a_ne^{in\si}$ 表示，从而给出 $\mathcal{K}_{\rm te}$ 到 $\mathcal{K}_{\rm al}$ 的一个嵌入。然而，“缓增”这一要求对系数 $a_n$ 施加了增长条件，使得 $\mathcal{K}_{\rm te}$ 严格小于 $\mathcal{K}_{\rm al}$。

E. Frenkel 在 \cite{Fr10} 中为 $\CH^{\sst (2)}_x$ 提出了一个代数对应物，此处记作 $\CH^{\sst (2)}_{{\rm al},x}$；他在那里证明了 $\CH^{\sst (2)}_{{\rm al},x}\simeq \mathcal{CB}_{\rm al}(C,\hfg_k,\mathcal{E}_x)$。我们将论证：在此语境下应当为 $\CH^{\sst (2)}_{x}$ 采用另一个不同的定义，从而得到同构 $\CH^{\sst (2)}_{x}\simeq \mathcal{CB}_{\rm te}(C,\hfg_k,\mathcal{E}_x)$，其中 $\mathcal{CB}_{\rm te}(C,\hfg_k,\mathcal{E}_x)$ 是 $\mathcal{CB}_{\rm al}(C,\hfg_k,\mathcal{E}_x)$ 中“缓增”共形块的子空间。在这篇短文里我们无法完整讨论 $\mathcal{CB}_{\rm te}(C,\hfg_k,\mathcal{E}_x)$ 的定义，拟另文再谈。下文将改为讨论来自 $2d$ sigma 模型的、支持 $\CH^{\sst (2)}_{x}\simeq \mathcal{CB}_{\rm te}(C,\hfg_k,\mathcal{E}_x)$ 的证据。

\subsubsection{模结构}

首先注意到，$\CH^{\sst (2)}_x$ 上带有\emph{两个}自然的模结构，它们来自与代数 $\mathcal{A}_{\rm cc}^\hbar={\rm Hom}_{\CM_G(C)}(\mathsf{B}_{\rm cc},\mathsf{B}_{cc})$ 的元素相联系的顶点算符。对于我们所关心的情形，应借助 NAH 对应把 $\CM_H(G)$ 表示为 $C$ 上平坦复联络的模空间；这样便可以考虑分别来自以下两种表示的两种代数结构：把 $\CM_H(G)$ 表示为特征簇$^{(g)}$ $\CM_{B}(G)$，以及表示为偶对 $(\CE,\nabla')$ 的模空间 $\CM_{dR}(G)$。不同的代数结构决定了在 $\mathcal{A}_{\rm cc}$ 中定义顶点算符时所用基本场变量的不同集合。相应地，可以区分
$\mathcal{A}_{B}^\hbar={\rm Hom}_{\CM_B(C)}(\mathsf{B}_{\rm cc},\mathsf{B}_{cc})$
与
$\mathcal{A}_{dR}^\hbar={\rm Hom}_{\CM_{dR}(C)}(\mathsf{B}_{\rm cc},\mathsf{B}_{cc})$，
其中 $\hbar=\ep_1/\ep_2$。

一方面，\cite{NW} 论证了代数 $\mathcal{A}_{B}^\hbar$ 可以等同于 $\CM_{B}(G)$ 上函数的量子化代数。另一方面，\cite{Ka08} 中的论证可以施用于 \cite{NW} 所研究的情形，从而得出结论：$\mathcal{A}_{dR}^\hbar\simeq \CD_{\hbar}$，即 $\CM_{dR}(G)$ 上函数的量子化代数；而后者又可以等同于 $\mathrm{Bun}_G$ 上某个线丛 $\CL^{-h^\vee-\ep_2/\ep_1}$ 上的微分算子代数。

众所周知，能级为 $k=-h^{\vee}-\frac{\ep_2}{\ep_1}$ 的仿射李代数 $\hat{\mathfrak{g}}_k$ 的共形块空间 $\mathcal{CB}_{\rm al}(C,\hat{\mathfrak{g}}_k,\mathcal{E}_x)$ 具有关于 $\mathcal{A}_{dR}^{\hbar}$ 的模结构。这个 $\mathcal{A}_{dR}^{\hbar}$-模结构是其定义性 Ward 恒等式的直接推论。

另一方面，$\mathcal{A}_{B}^{\hbar}$-模结构可以这样在共形块空间上定义：用 $\hat{\mathfrak{g}}_k$ 的退化表示来定义 \cite{AGGTV,DGOT} 中 Verlinde 线算符在此情形的类似物。我们将在别处证明：仿射李代数共形块上 Verlinde 线算符的代数可以等同于 $\mathcal{A}_{B}^{\hbar}$。现有的定义只对用拼接构造得到的共形块有效。这类共形块具有特别好的性质，其中之一是：典范联络可以在整个 $\mathrm{Bun}_G^{\rm vs}$ 上积分，而 $\mathrm{Bun}_G^{\rm vs}$ 是 $\mathrm{Bun}_G$ 中包含“非常稳定”（very stable）丛（即不承认幂零 Higgs 场的丛）的子集；对 $\mathcal{CB}^{\rm al}(C,\hat{\mathfrak{g}}_k,\mathcal{E}_x)$ 的一般元素而言则并非如此。这提示：$\mathcal{A}_{B}^{\hbar}$-模结构只能定义在“缓增”共形块的适当子空间 $\mathcal{CB}_{\rm te}(C,\hat{\mathfrak{g}}_k,\mathcal{E}_x)$ 上。

$\CH^{\rm (2)}_x$ 上存在两个模结构，可以看作支持所猜想同构
$\CH^{\rm (2)}_x\simeq\mathcal{CB}_{\rm te}(C,\hat{\mathfrak{g}}_k,\mathcal{E}_x)$
的第一条证据。下文将给出更多证据。

\subsubsection{与 $\bar\partial$-上同调的关系}

先回到 Nekrasov 与 Witten 在 \cite{NW} 中所研究的情形。首先注意，空间 ${\rm Hom}_{\CM_H(G)}(\mathsf{B}_{\rm cc}, \mathsf{B}_{\rm op})$ 可以在辛结构为 $\omega=\omega_I$ 的 A-模型中得到实现。长条（strip）上的边界条件取为典范余迷向 A-膜 $\mathsf{B}_{\rm cc}$ 与 Lagrangian A-膜 $\mathsf{B}_{\rm op}$。Chan--Paton 曲率 $F$ 由 Hitchin 模空间上的辛形式 $\omega_J$ 给出。于是由典范余迷向 A-膜 $\mathsf{B}_{\rm cc}$ 所确定的复结构为 $\omega^{-1}F=K$。\cite{NW} 论证：在此情形下，空间 ${\rm Hom}_{\CM_H(G)}(\mathsf{B}_{\rm cc}, \mathsf{B}_{\rm op})$ 可以等同于 $\CM_{dR}(G)$ 中 opers 子空间 $\mathrm{Op}_{\fsl_2}(C)$ 上、线丛 $K^{{1}/{2}}_{\rm op}$ 的全纯截面空间。事实上，${\rm Hom}_{\CM_H(G)}(\mathsf{B}_{\rm cc}, \mathsf{B}_{\rm op})$ 将可实现为长条上 A-模型零模的函数空间。对于代表 $\mathrm{Op}_{\fsl_2}(C)$ 上坐标的 sigma 模型场，我们在长条的两端都有 Neumann 型边界条件；而其余的场只在一端有 Dirichlet 型边界条件。因此 A-模型的零模可以用 $\mathrm{Op}_{\fsl_2}(C)$ 上的坐标表示。于是可以援引 \cite{NW} 中简述的标准论证，把长条上 sigma 模型的超对称基态空间等同于取值于 $K^{{1}/{2}}_{\rm op}$ 的 $\bar{\partial}$-上同调。

\subsubsection{与共形块的关系：Virasoro 情形}\label{HomVirConf}

注意 $\mathrm{Op}_{\fsl_2}(C)$（非典范地）同构于复向量空间 $H^0(C,K^2_C)$，后者在拓扑上是平凡的。而空间 $H^0(C,K^2_C)$ 典范同构于 Teichmüller 空间 $\mathcal{T}(C)$ 的余切纤维 $T^*\mathcal{T}(C)|_{C}$。对我们来说重要的是要看到：存在一个从余切空间 $T^*\mathcal{T}(C)|_{C}$ 到 Teichmüller 空间 $\mathcal{T}(C)$ 的映射，它把这两个空间上的自然复结构联系起来。

为了在 $\mathcal{T}(C)$ 上引入复结构，可以用调和 Beltrami 微分来描述 Teichmüller 变分：它们是满足 $\pa_z(\eta\mu)=0$ 的 $(-1,1)$-形式 $\mu$，其中 $\eta$ 是由 $C$ 上复结构唯一确定的双曲度量。于是调和 Beltrami 微分构成的向量空间上的复结构，就定义了 $\mathcal{T}(C)$ 上的复结构。对每个全纯 $(2,0)$-形式 $\theta$，可以配上相应的调和 Beltrami 微分 $\eta^{-1}\bar{\theta}$。这给出一个从 $H^0(C,K^2_C)$ 到调和 Beltrami 微分空间的复反线性映射，它表达了 $\mathcal{T}(C)$ 上与 $H^0(C,K^2_C)$ 上自然复结构之间的关系。由此可以把 $\mathrm{Op}_{\fsl_2}(C)$ 上的全纯函数空间，等同于 $\mathcal{T}(C)$ 上（反）全纯函数的空间 ${\rm Fun}_{\rm hol}(\mathcal{T}(C))$。

另一方面，已知 ${\rm Fun}_{\rm hol}(\mathcal{T}(C))$ 同构于 Virasoro 共形块的（子）空间。闭曲面\footnote{为简化叙述，我们暂时限于这一情形。} $C$ 上的共形块定义为 Virasoro 代数真空表示 $V_0$ 上的线性泛函 $f$，它满足共形 Ward 恒等式，后者刻画 $f$ 在一个点之外全纯的向量场代数之自然作用下的不变性。共形块构成的向量空间带有一个典范的射影平坦联络，它由 Virasoro 代数在 $V_0$ 上的作用定义。这一典范联络的曲率可以在复结构模空间 $\CM_{g}$ 上局部地平凡化。在所有代数定义的共形块空间的适当子空间上，可以对所得的平坦联络积分，从而在带局部坐标 $\tau$ 的 $\CM_{g}$ 的开子集上定义水平截面 $f_{\tau}$。共形块在最高权向量 $v_0$ 上的取值 $\CZ_f(\tau)=f_\tau(v_0)$ 称为手征配分函数。Virasoro 单值化定理与共形 Ward 恒等式把 $\CZ_f(\tau)$ 的导数与刻画共形块 $f$ 的数据联系起来。由此得到（局部定义的）函数 $\CZ(\tau)$ 与“可积”共形块 $f$ 之间的一一对应。这一对应本质上是典范的：Virasoro 单值化描述了 $\CM_{g}$ 的局部结构，而这一点已经编码在共形块的定义之中。

\subsubsection{与共形块的关系：Kac--Moody 情形}

可以论证 $\CH_x^{\sst (2)}$ 同构于 $\hat{\mathfrak{g}}_k$ 的、被丛 $\CE_x$ 扭（twisted）的共形块空间 $\mathcal{CB}(C,\hat{\mathfrak{g}}_k,\CE_x)$，论证方式与第 \ref{HomVirConf} 节类似。依照那里所述的同样推理，将会得到结论：$\CH_{\mathcal{E}}$ 同构于 $H^0(C,{\rm End}(\mathcal{E})\otimes K_C)$ 上的全纯函数空间。空间 $H^0(C,{\rm End}(\mathcal{E})\otimes K_C)$ 正是全纯余切空间 $T^*{\rm Bun}_G$。$T{\rm Bun}_G$ 上存在一个典范的 Hermite 度量，它类似于 Teichmüller 理论中的 Weil--Petersson 度量，其定义以 Narasimhan--Seshadri 定理作为单值化定理的类比。该度量把全纯余切空间等同于反全纯切空间，从而把 $H^0(C,{\rm End}(\mathcal{E})\otimes K_C)$ 与 $T{\rm Bun}_G$ 上的自然复结构联系起来。

$T{\rm Bun}_G$ 上的全纯函数空间自然地嵌入 $\on{Jet}_\mathcal{E}$，即 $\mathcal{E}$ 在 ${\rm Bun}_G$ 中形式邻域上的函数空间。剩下只需注意 $\on{Jet}_\mathcal{E}$ 同构于共形块空间，$\on{Jet}_\mathcal{E}\simeq\mathcal{CB}_{\rm al}(C,\hat{\mathfrak{g}}_k,\CE_x)$ \cite[Section 18.2]{FBZ}。E. Frenkel 正是利用这一观察在 \cite{Fr10} 中证明了我们论断的代数对应物 $\CH^{\sst (2)}_{{\rm al},x} \simeq \mathcal{CB}_{\rm al}(C,\hfg_k,\CE_x)$。由此可见，$T{\rm Bun}_G$ 上的全纯函数空间可以等同于 $\mathcal{CB}(C,\hat{\mathfrak{g}}_k,\CE_x)$ 中“性质良好”（well-behaved）的共形块的子空间，这进一步支持了上面提出的猜想 $\CH^{\sst (2)}_{x} \simeq \mathcal{CB}_{\rm te}(C,\hfg_k,\CE_x)$。

\bigskip
\noindent\rule{\textwidth}{0.6pt}

% [Tutor Note: 式 \eqref{Nahm} 与紧接其后的文字不自洽：正文定义的是 $\mathfrak{t}_+=\mathfrak{t}_1+i\mathfrak{t}_2$，而式中出现的是 $\mathfrak{t}_-$（极点系数）与 $\mathfrak{t}_3$（$A_y$ 的极点系数）。按通常约定 $\mathfrak{t}_+ = \mathfrak{t}_1+i\mathfrak{t}_2$、$\mathfrak{t}_- = \mathfrak{t}_1-i\mathfrak{t}_2$，则 $A_z$ 的极点系数应是 $\mathfrak{t}_-$、而文中只给出了 $\mathfrak{t}_+$ 的定义，读者若自行补 $\mathfrak{t}_-$ 须取复共轭（在 $\mathfrak{sl}_2$ 的实形式下 $\mathfrak{t}_-=\mathfrak{t}_+^\dagger$）。另注意第三条中的 $\mathcal{A}_1$ 与前文 $\CD_y=\pa_y+A_y-i\phi_y$ 中的 $A_y$ 是同一个分量的两种写法（$\mathcal{A}$ 带花体在别处专指 twisted connection），核对 Nahm 极点条件时不要把它误当成第四分量。]
% [Tutor Note: 式 \eqref{BPS} 到式 \eqref{Fterm} 之间被略去的一步是 Hitchin 方程（此处为三维约化后的 Nahm 方程）的超 Kähler 商论证：\eqref{Dterm} 的矩映射方程可通过一个复规范变换解出，而该复规范变换的选取不改变 \eqref{Fterm} 的复平坦性；也就是说，“BPS 解模紧致规范变换”与“F-项解模复规范变换”同构，靠的是 Kempf--Ness 型定理在无穷维的类比（Donaldson--Uhlenbeck--Yau），并且在 $y\ra0$ 处还需对渐近行为加正则性条件，否则模空间不具有限维性。作者引用 \cite{GW} 而未写出此论证。另注意算子定义中 $\zeta$ 的幂次并不对称：$A_z$ 带一个 $\zeta$、$\phi_{\bar z}$ 带一个 $\zeta$，而 $A_{\bar z}$、$\phi_z$ 不带；作者已加脚注声明与 \cite{GW} 约定不同；重标度 $\zeta$ 会同时改变 $\Om_\zeta$ 与 $I_\zeta$，从而改变膜的 A/B 类型判定，与 \cite{NW} 对照时不可混用两套归一化。]
% [Tutor Note: “固定 $\CE$ 的子流形关于 $\Omega_\zeta$ 是 Lagrangian 的”一句是本文后面所有膜判定的基础，但此处只作断言。实际上要用到两件事：(i) NAH 对应给出的是实解析（超 Kähler）同构，而非代数同构，它把 $\CM_{dR}$ 上的复结构 $I_\zeta$ 与全纯辛形式 $\Omega_\zeta$ 搬到 $\CM_H$ 一侧；(ii) Hitchin 第二纤维化 $\CM_H(G)\to\Bun_G$（丢掉 Higgs 场）的纤维维数为 $\dim\Bun_G$，恰为 $\CM_H$ 复维数的一半，且 $\Omega_\zeta$ 沿纤维为零——这一点来自 Hitchin 可积系统中 Poisson 对合函数的构造，需要单独验证，不能由“固定 $\CE$”直接看出。]
% [Tutor Note: $\mathcal{A}_{dR}^\hbar\simeq\CD_\hbar$ 并“等同于 $\CL^{-h^\vee-\ep_2/\ep_1}$ 上微分算子代数”这一步用的是 Beilinson--Drinfeld 关于临界能级中心（其谱即 opers）的定理之量子化。核对幂次：因 $k=-h^\vee-\ep_2/\ep_1$，指数 $-h^\vee-\ep_2/\ep_1$ 恰好就是 $k$，即线丛取 $\CL^{k}$；若按常见的 $\CL^{\otimes(k+h^\vee)}$ 写法，则指数是 $-\ep_2/\ep_1$，两者相差 $h^\vee$，极易写错。另外 $\hbar=\ep_1/\ep_2$ 与 $k+h^\vee=-\ep_2/\ep_1=-1/\hbar$ 互为倒数，后文比较 $W_{1/2\hbar}$ 时须一并核对。]
% [Tutor Note: 关于 $(-1,1)$-形式与条件 $\pa_z(\eta\mu)=0$ 的指标核算：$\mu$ 按 $\mu=\eta^{-1}\bar\theta$ 的类型计为 $K^{-1}\bar K$ 的截面（即 $(-1,1)$-型），与全纯二次微分 $\theta\in H^0(C,K^2_C)$ 配对后 $\mu\theta$ 是 $(1,1)$-形式方可积分；$\eta\mu=\bar\theta$ 的类型是 $(0,2)$，故 $\pa_z(\eta\mu)=0$ 应理解为“$\eta\mu$ 是 $C$ 上的全纯二次微分”（$\pa_z$ 作用在系数函数上），它同时给出 $\mu$ 的调和性。若把 $\mu$ 误当作 $(1,-1)$-型（即反变指标写反），$\pa_z$ 与 $\bar\pa_{\bar z}$ 的位置就会颠倒，$(-1,1)$ 与 $(1,-1)$ 也常被混用。]
% [Tutor Note: 两处易被带过的措辞：(i) “very stable 丛”的严格定义是“不承认\emph{非零}幂零 Higgs 场”，原文漏了“非零”二字（零 Higgs 场总是存在的），这一定义源于 Drinfeld 对 very stable 丛的引入（Laumon 讨论了其变体），且与“典范联络可在 $\mathrm{Bun}_G^{\rm vs}$ 上积分”直接相关；(ii) 缓增空间与形式幂级数空间之间，原文说嵌入方向是 $\mathcal{K}_{\rm te}\hookrightarrow\mathcal{K}_{\rm al}$，这是对的（多项式增长的 Fourier 系数必为形式幂级数），但作者未给出 $\mathcal{CB}_{\rm te}$ 的任何定义，读者不应把“tempered”误解为 $L^2$ 完备化——它对应的是 Fourier 系数（即共形块在穿刺点附近的模参数展开系数）的多项式增长条件。]

%=====================================================================
% 第 4 节  配分函数与共形块
% 原文：Section 4. Partition functions versus conformal blocks
%       （arXiv:1702.06499v2）
%=====================================================================
\section{配分函数与共形块}\label{partfct}

\begin{keytakeaways}
\begin{itemize}
  \item \textbf{物理图像：}AGT 对应及其推广以配分函数为核心，而 Kapustin--Witten 处理几何 Langlands 对应的进路中根本不出现配分函数。本节要在二维 sigma 模型中为二者找到统一的表述：四维半球上的路径积分给出拓扑向量空间中的一个向量，其波函数就是瞬子配分函数；而在二维一侧，同一个配分函数可以按两种不同的 Hamilton 解释读出。
  \item \textbf{数学工具：}带 $\Om$-形变的四维半球 $B^4_{\ep_1\ep_2}$ 上的拓扑扭曲、开/闭 TQFT 的态—边界态对应、扭 $\CD$-模与 KZB 方程，以及复 Fenchel--Nielsen 坐标下特征簇的可积结构。
  \item \textbf{关键结论：}配分函数 $\CZ(a;x,\tau;\ep_1,\ep_2)$ 既可以看成 $\CH_x^{\sst(2)}=\Hom(\mathsf{B}_{\rm cc},\mathsf{L}_x^{\sst(2)})$ 中元素的波函数，也可以看成 $\CH_a^{\sst(1)}=\Hom(\mathsf{B}_{\rm cc},\mathsf{L}_a^{\sst(1)})$ 中元素的波函数；这两种 Hamilton 解释由模不变性的开 TQFT 版本相联系，正是把 AGT 一侧与几何 Langlands 一侧缝合起来的关键。
  \item \textbf{关键构造：}边界条件 $\mathsf{L}^{\sst(1)}_{a}$ 实现为特征簇 $\CM_B(G)$ 中由长度函数 $a$ 生成的 Hamilton 流的轨道，再经和乐映射与 NAH 对应搬到 $\CM_H(G)$；这解释了记号中的“第一纤维化”。
\end{itemize}
\end{keytakeaways}

\bigskip
\noindent\rule{\textwidth}{0.6pt}
\bigskip

配分函数在 AGT 对应及其各种推广中处于核心地位，但在 Kapustin 与 Witten 处理几何 Langlands 对应的进路中却并不出现。为了建立一个统一的框架，下面我们讨论如何在二维 sigma 模型所提供的有效描述中定义相应的配分函数。

\subsection{四维与二维的 TQFT 框架}

利用 $\CS$ 类理论到二维的约化，下面我们要为配分函数在以 $\CM_H(G)$ 为目标流形的二维 sigma 模型中的描述提供动因。这一描述将基于又一类边界条件，记作 $\mathsf{L}^{\sst (1)}_{a}$。记号 $\mathsf{L}^{\sst (1)}_{a}$ 的选取源自它与 Hitchin \emph{第一}纤维化的联系，这一点后文将讨论。

仿照 \cite{NW}，我们考虑带 $\Om$-形变的、定义在半球 $B_{\ep_1\ep_2}^4$ 上的拓扑扭曲 $\CS$ 类理论。拓扑扭曲的 $\CS$ 类理论为 $M^3_{\ep_1\ep_2}=\pa B^4_{\ep_1\ep_2}$ 配上一个向量空间 $\CH_{\rm top}=Z(M^3_{\ep_1\ep_2})$，这里把它等同于超荷 $Q$ 的上同调，而 $Q$ 是能够在 $B_{\ep_1\ep_2}^4$ 上保持的那个超荷。可以用四维半球 $B_{\ep_1\ep_2}^4$ 上的路径积分定义一个向量 $\Psi \in\CH_{\rm top}$。向量 $\Psi$ 的波函数 $\Psi(a)$ 可以等同于配分函数 $Z(B^4_{\ep_1,\ep_2};\mathcal{B}_a)$，后者通过在 $M^3_{\ep_1\ep_2}$ 上施加适当的、由参数 $a$ 标记的 $Q$-不变边界条件 $\mathcal{B}_a$ 来定义。这类边界条件可以等同于定义 Nekrasov 配分函数时在 $\BR^4_{\ep_1\ep_2}$ 的无穷远处所加的边界条件，它们特别地把向量多重态中各个标量的零模固定为向量 $a=(a_1,\dots,a_{3g-3+n})$ 所收集的取值。边界条件 $\mathcal{B}_a$ 定义了一族边界态 $\beta_a$，从而可以把 $Z(B^4_{\ep_1,\ep_2};\mathcal{B}_a)$ 表示为交叠 $\langle\beta_a,\Psi\rangle$。

在把 $\CS$ 类理论约化到 $2d$ 拓扑 sigma 模型时，应当得到上面引入的 4d TQFT 数据的如下表示：
\begin{itemize}
\item 向量空间 $\CH_{\rm top}\simeq Z(S^3_{\ep_1,\ep_2})\ra Z(I)$。
\item 向量 $\Psi=Z(B^4_{\ep_1,\ep_2})\ra Z(T_{\ep_1,\ep_2})\in\CH_{\rm top}$，其中 $T_{\ep_1,\ep_2}$ 是挖去“上”边之后的那个开三角形，它在拓扑上等价于 $\BR_-\times I$，并通过在 $\BR_-$ 的无穷远处添加一个点而部分紧化。$T_{\ep_1,\ep_2}$ 的边界为 $\{0\}\times I$。
\item 配分函数 $Z(B^4_{\ep_1,\ep_2};\mathcal{B}_a)\ra Z(T_{\ep_1,\ep_2};\mathsf{B}_a)$ 则与三角形 $T_{\ep_1,\ep_2}$ 相联系，其上边 $\{0\}\times I$ 被赋予边界条件 $\mathsf{L}_{a}^{\sst (1)}$。$\mathsf{L}_{a}^{\sst (1)}$ 是由赋给 $M^3_{\ep_1\ep_2}$ 的边界条件 $\mathcal{B}_a$ 经约化到一维而定义的。
\end{itemize}
这意味着瞬子配分函数 $\CZ(a;x;\tau;\ep_1,\ep_2)$ 可以由三角形上 sigma 模型的配分函数表示，而该三角形的各边分别染上 $(\mathsf{B}_{\rm cc}, \mathsf{L}_a^{\sst (1)}, \mathsf{L}_x^{\sst (2)})$。

\subsection{$\CD$-模与配分函数}

量子理论中的态 $\psi$ 抽象地由 Hilbert 空间 $\CH$ 中的向量表示。要得到一个作为波函数空间 $\psi(x)$ 的具体表示，需要引入 $\CH$ 的（Hermite）对偶 $\CH^\dagger$ 中的一族元素 $\de_x$，使得 $\psi(x)$ 可以表示为 $\langle\de_x,\psi\rangle$，其中 $\langle . , .\rangle:\CH^{\dagger}\times\CH\ra \C$ 是自然的配对。

在目前的语境下，我们可以把配分函数 $\CZ=\CZ(a;x,\tau;\ep_1,\ep_2)$ 解释为元素 $\phi_x\in \CH_{x}^{\sst(2)}:=
{\rm Hom}_{\CM_H(G)}^{}(\mathsf{B}_{\rm cc}^{},\mathsf{L}_{x}^{\sst(2)})$ 的波函数 $\phi_x(a)=\langle\de_a^{\sst(1)},\phi_x^{}\rangle$，该元素由半无限长条 $\BR_-\times I$ 上的 sigma 模型路径积分产生。“边界态”$\de_a^{\sst(1)}$ 表示的正是加在 $\{0\}\times I$ 上的边界条件 $\mathsf{L}_{a}^{\sst(1)}$。

交叠 $\langle\de_a^{\sst(2)},\phi_x^{}\rangle$ 可以表示为拓扑 sigma 模型在三角形上的一个配分函数，该三角形三条边上分别赋以三种不同的边界条件 $(\mathsf{B}_{\rm cc},\mathsf{L}_{x}^{\sst(2)},\mathsf{L}_{a}^{\sst(1)})$。时间反演对称性把 $\CZ$ 的复共轭 $\CZ^*$ 等同于拓扑 sigma 模型在三角形上的一个配分函数，只是其边界条件 $(\mathsf{B}_{\rm cc}^{},\mathsf{L}_{a}^{\sst(1)},\mathsf{L}_{x}^{\sst(2)})$ 以不同的次序出现。$\CZ^*$ 具有一个 Hamilton 解释：它是元素 $\vf_a\in \CH_{x}^{\sst(1)}:=
{\rm Hom}_{\CM_H(G)}^{}(\mathsf{B}_{\rm cc}^{},\mathsf{L}_{a}^{\sst(1)})$ 的波函数 $\vf_a(x)=\langle\de_x^{\sst(2)},\vf_a^{}\rangle$，而该元素由半无限长条 $\BR_-\times I$ 上的 sigma 模型路径积分产生。

两个空间 $\CH_{a}^{\sst(1)}$ 与 $\CH_{x}^{\sst(2)}$ 都具有 $\mathcal{A}_{dR}^{\hbar}$-模结构，它来自代数 $\mathcal{A}_{dR}^{\hbar}$ 在 ${\rm Hom}_{\CM_H(G)}^{}(\mathsf{B}_{\rm cc}^{},\mathsf{B})$ 上的自然作用。上文已论证，这一模结构与由仿射代数 Ward 恒等式导出的、共形块空间上的模结构相重合。由此得到的扭 $\CD$-模结构，与手征配分函数（对应于仿射李代数的共形块）所满足的 KZB 方程所定义的结构相重合。配分函数 $\CZ$ 是 KZB 方程的解。$\CD$-模 $\CH_{a}^{\sst(1)}$ 与 $\CH_{x}^{\sst(2)}$ 可以由相应的波函数空间完全刻画，而这些波函数在物理上就表现为配分函数 $\CZ$。

在下面将要讨论的 $\ep_2\ra 0$ 极限下，我们将提出：$\CH_{{a}}^{\sst (1)}$ 出现在 Kapustin--Witten 处理几何 Langlands 对应的进路中，而 $\CH_x^{\sst(2)}$ 在 AGT 对应的各种推广的语境下似乎是更自然要考虑的空间。由模不变性的一个开 TQFT 版本所给出的、配分函数 $\CZ(x,{a},\tau;\ep_1,\ep_2)$ 之两种 Hamilton 解释之间的关系，正是这两个故事之间的关键纽带。

\subsection{膜 $\mathsf{L}_{a}^{\sst (1)}$ 的 sigma 模型描述}\label{branes1}

现在我们来给出边界条件 $\mathsf{L}_{a}^{\sst(1)}$ 作为 $\CM_H(G)$ 中 Lagrangian 子流形的一个描述。为后文之便，我们把指代用于定义膜 $\mathsf{L}_{a}^{\sst(1)}$ 的 $\CS$ 类理论之 Lagrangian 表示的标号 $\si$ 并入记号，改记为 $\mathsf{L}_{\si,a}^{\sst(1)}$。我们提出：膜 $\mathsf{L}_{\si,a}^{\sst(1)}$ 可以描述为特征簇 $\CM_B(G)$ 中由一族长度函数 $a=(a_1,\dots,a_{3g-3+n})$ 所生成的 Hamilton 流的轨道；这一族长度函数构成复 Fenchel--Nielsen 坐标系{$^{g)}$}中互相对易（Poisson 对易）的那一半，而该坐标系曾在 \cite{NRS} 的一个相关语境中被使用。这些轨道经和乐映射与 NAH 对应映到 Hitchin 模空间 $\CM_H(G)$。从特征簇 $\CM_B(G)$ 的可积结构来看，膜 $\mathsf{L}_{{\ep},a}^{\sst(1)\si}$ 是自然的对象。不同的复 Fenchel--Nielsen 坐标系由裤分解 $\si$ 标记。

我们的提议基于 {\crd \cite{DMO}} 中已知的关系：$\CS$ 类理论中的 Wilson 圈可观测量与六维 $(2,0)$ 理论中支撑在两个圆 $S^1_\CS$ 与 $S^1_C$ 之积上的余四维缺陷相关联，而这两个圆分别位于 $\CM^4$ 与 $C$ 之中。这些关系可以推广到存在面算符的情形。若把注意力集中在与 $\CM^4$ 中面算符不共享支撑的 Wilson 圈上，则可如前论证：Wilson 圈可观测量的本征值可以用标量零模 $a$ 来表示，而 $a$ 正是出现在瞬子配分函数 $\CZ({\crd a;x},\tau;\ep_1,\ep_2)$ 自变量中的量。

另一种做法是把 \cite{AGGTV} 中的论证推广到带有额外余二维面算符的情形，而这可以建立在第 \ref{Hecke1} 节所述结果的基础之上。

\bigskip
\noindent\rule{\textwidth}{0.6pt}

% [Tutor Note: 第 1 小节 itemize 第三条内部自相矛盾：先写 $Z(T_{\ep_1,\ep_2};\mathsf{B}_a)$，紧接着又说上边赋的是边界条件 $\mathsf{L}_{a}^{\sst (1)}$，并通篇只用 $\mathsf{L}_a^{\sst (1)}$。按上下文 $\mathsf{B}_a$ 应为 $\mathsf{L}_a^{\sst (1)}$ 的笔误（$\mathsf{B}$ 是膜/boundary condition 的旧记号，$\mathcal{B}_a$ 则是四维边界条件）；译文按原文保留，读者不要据此认为存在两个不同的边界条件。]
% [Tutor Note: 第二小节的指标有系统性的错位：由 $L_a^{\sst(1)}$ 定义的边界态本应记作 $\de_a^{\sst(1)}$，但正文写成交叠 $\langle\de_a^{\sst(2)},\phi_x\rangle$，上标 $(2)$ 与 $\phi_x$ 一侧重复；同样，$\vf_a$ 被声明属于 $\CH_{x}^{\sst(1)}:=\Hom(\mathsf{B}_{\rm cc},\mathsf{L}_{a}^{\sst(1)})$，下标 $x$ 应为 $a$（上标 $(1)$ 正确）。后文又出现正确的 $\CH_{a}^{\sst(1)}$，说明是排印疏漏。核对波函数的自变量与下标是否匹配（$\phi_x(a)$ 对 $\CH^{(2)}_x$、$\vf_a(x)$ 对 $\CH^{(1)}_a$）是理解本节的关键。]
% [Tutor Note: 配分函数自变量的分隔符在原文中不统一：$\CZ(a;x;\tau;\ep_1,\ep_2)$（分号分隔 $x$ 与 $\tau$）、$\CZ(a;x,\tau;\ep_1,\ep_2)$、$\CZ(x,{a},\tau;\ep_1,\ep_2)$（$x$ 与 $a$ 的次序颠倒且全用逗号）三种写法并存。按本文约定，分号用于区分四维数据 $(a;\tau,\ep)$ 与二维/共形块数据 $x$，因此第二种写法是规范的；第三种不仅是书写次序问题，其“两种 Hamilton 解释”论断恰恰要求把 $a$ 与 $x$ 视为分别属于 $\CH^{(1)}$ 与 $\CH^{(2)}$ 的两套标记，混写时须小心不要颠倒 $\langle\de_a,\phi_x\rangle$ 与 $\langle\de_x,\vf_a\rangle$ 的配对方向。]
% [Tutor Note: 膜 $\mathsf{L}_{{\ep},a}^{\sst(1)\si}$ 的记号中多出一个 $\ep$ 下标：同段前后文均只有 $\mathsf{L}_{\si,a}^{\sst(1)}$ 与 $\mathsf{L}_{a}^{\sst(1)}$，且 $\ep$ 在全文专指 $\Om$-形变参数 $\ep_1,\ep_2$，此处无任何作用。对照第 7 节对同一对象的引用（那里写作 $\mathsf{L}_{\si,a}^{\sst (1)}$）可以确认这是排印时残留的标记。]
% [Tutor Note: 4d TQFT 段落里“$\CH_{\rm top}=Z(M^3)$ 等同于 $Q$ 的上同调”这一等同需要扭曲使 $Q$ 幂零、并取 $Q$-上同调才是拓扑不变量（Witten 型论证）；而“向量 $\Psi$ 由半球上的路径积分给出、其波函数即 $\langle\beta_a,\Psi\rangle$”是开/闭 TQFT（Atiyah 公理）中态—边界态的通用构造，并不需要 $\CS$ 类理论的具体细节。真正被跳过的是末段“两种 Hamilton 解释由模不变性的开 TQFT 版本相联系”：这本质上是开弦—闭弦对偶/Cardy 条件在有边界与缺陷情形下的推广，其成立条件（边界态构成完备基、模变换的收敛性）在无穷维情形并未验证，作者把它当作出发点而非结论，引用时不宜说成已证明。]
% [Tutor Note: 记号 $S^1_\CS$ 与 $S^1_C$ 的下标在原文中并不自洽：全文的宏 \CS 定义为花体 $\mathcal{S}$（class S 的 S），故这里排出来是 $S^1_{\mathcal S}$，但从上下文（complex CS theory、Chern--Simons）看作者想写的是 Chern--Simons 的“CS”。译文按原文保留该记号。读者须按语境判断 $\CS$ 指 class S 还是 Chern--Simons：本节其余各处（$\CS$ 类理论）都是前者，只有 $S^1_\CS$ 这一处是后者，核对余四维缺陷所缠绕的两个圆时不要混淆。]

%=====================================================================
% 第 5 节  另一类面算符
% 原文：Section 5. Another type of surface operators
%       （arXiv:1702.06499v2）
%=====================================================================
\section{另一类面算符}

\begin{keytakeaways}
\begin{itemize}
  \item \textbf{物理图像：}除余二维面算符之外，六维 $(2,0)$ 理论还可被余四维缺陷修饰。弦论中把 M5-M5'-膜系统沿 $x_7$ 方向拉开，会在其间接上若干 M2-膜，由此得到一个红外等价的 M5-M2 构型；M2 的端点扫出 $C$ 上的除子 $\sum_k P_k$。于是“在 $P$ 处插入一个余四维面算符”就对应于“在 $P$ 处插入一个退化场”。
  \item \textbf{数学工具：}Hanany--Witten 膜产生效应、AGT 对应在（推广的）共形块层面的翻译、仿射李代数 $\widehat{\mathfrak{sl}}_{2,k}$ 的模 $W_{-(k+h^{\vee})\la}$ 与 Virasoro 退化表示 $V_{-1/2b}$、Sklyanin 分离变量（SOV）方法，以及二维单值表示。
  \item \textbf{关键结论：}本节的核心猜想（式 \rf{Heckeconj}）是：在 $P\in C$ 处添加余四维面算符，等价于在共形块中 $P$ 处插入 $\widehat{\mathfrak{sl}}_{2,k}$-模 $W_{1/2\hbar}$。这正好把 Beilinson--Drinfeld 用 $W_{-(k+h^{\vee})\la}$ 描述的 Hecke 函子翻译成物理语言，$\la=1/2$ 对应 ${}^LG=\mathrm{PSL}(2,\C)$ 的基本表示。
  \item \textbf{证据与缺口：}证据来自 M5-M5' 与 M5-M2 的红外对偶、$W_{1/2\hbar}$ 在临界能级极限下的因子化、以及它产生二维单值表示这一唯一性论证；关键的局域化检验尚未进行。
\end{itemize}
\end{keytakeaways}

\bigskip
\noindent\rule{\textwidth}{0.6pt}
\bigskip

几何 Langlands 对应的关键要素之一是 Hecke 本征值性质。为了给下一节关于 Hecke 本征值性质的讨论作准备，我们在此讨论另一类面算符，它为我们提供了 Hecke 函子（Hecke functor）的物理实现。

\subsection{弦论动机}

从六维 $(2,0)$ 理论的角度看，研究用余四维缺陷修改该理论会产生什么效应，似乎是很自然的。这类缺陷构成对六维 $(2,0)$ 理论及其维数约化的自然探针，相关研究在 AGT 对应的语境中已经相当深入 \cite{AGGTV}。下面我们简要讨论：若同时考虑余二维与余四维的面算符，图像将如何改变。这类系统的存在性由弦论提示。特别地，依照 \cite{FGT} 的讨论，可以考虑 M 理论中支撑如下表所示的 M5-M5'-M2-膜构型：
\begin{center}
\begin{tabular}{l || c|c|c|c|c|c|c|c|c|c|c}
膜 & 0 & 1 & 2 & 3 & 4 & 5 & 6 & 7 & 8 & 9 & 10 \\ \hline\hline
M5  & x & x & x & x &   &   & x &   &   &   & x  \\
M5' & x & x &   &   &   &   & x &   & x & x & x \\
M2  & x & x &   &   &   &   &  &  x &   &   &   \\
\end{tabular}
\end{center}
具有这类支撑意味着：$\CM^4\times C$ 上的余四维面算符（其中 $\CM^4$ 局部地是一个如下形式的纤维积\footnote{记号 
$S_i^1$ 中的下标指的是用来定义参量为 $\ep_i$ 的 $\Om$-形变的那一个旋转对称性。} $\BR\times I\times
S^1_1\times S^1_2$）将表现为
\begin{itemize}
\item
$\CM^4$ 上支撑在 $\CM^2\simeq \BR\times S_{1}^1$ 上的面算符，
\item
$\BR\times I\times C$ 上 $N=4$ SYM 中的线算符，支撑在 $\BR\times \{\pi\}\times \{P\}$，
\item
或 $\BR\times C$ 上复 {\crd CS}（Chern--Simons）理论中的线算符，支撑在 $\BR\times \{P\}$。
\end{itemize}
由这些观察很容易看出，这类余四维面算符的存在会修改描述约化到 $\BR\times I$ 的二维 sigma 模型中的边界条件。下面我们简要描述：在作为二维 CFT 的有效描述中，余四维面算符的存在会导致哪些修改。

\subsection{对（推广的）AGT 对应带来的修改}

\cite{AGGTV} 论证：六维 $(2,0)$ 理论中余四维面算符的存在，在 $C$ 上紧化之后被描述为这样一个系统——一个 $4d$ 的 $\CS$ 类理论与 $\CM^4$ 中某个二维曲面 $\CM^2$ 上的二维规范线性 sigma 模型（GLSM）相耦合。相关 GLSM 的特征在于：其红外描述是一个以 Grassmann 流形为目标空间的非线性 sigma 模型。在最简单的 $\mathfrak{g}= \mathfrak{sl}_2$ 情形，例如得到的就是流向 $\mathbb{CP}^1$ sigma 模型的 GLSM。

关于在 CFT 内部、在余二维面算符存在时如何实现余四维面算符，我们的主要猜想如下：
\begin{equation}\label{Heckeconj}
\boxed{ 
\quad
\begin{aligned}
&\text{在 $P\in C$ 处添加一个余四维面算符}\\
& \text{由在 $P$ 处插入 $\widehat{\mathfrak{sl}}_{2,k}$-模 $W_{1/2\hbar}$ 来表示。}
\end{aligned}
\quad
}
\end{equation}
为了列举这一猜想的现有证据，先注意到它是由 \cite{FGT} 所述图像的一个轻微推广所提示的。该文献观察到，由 M5-M5'-膜组成的系统呈现出一种有趣的红外对偶，把它与一个 M5-M2-膜系统联系起来，其中若干 M2-膜位于 $C$ 上的各个点。这一红外对偶来自 Hanany--Witten 膜产生效应（brane creation effect）的一个变体，后者描述了一族相互红外等价的、通过把上述构型中已有的 M5-M5'-膜沿 $x_7$ 轴分开而得到的有效理论。为了得到红外等价的有效描述，需要引入额外的、悬挂在正在被分开的膜之间并沿 $x_7$ 轴延伸的 $M2$-膜。这些 $M2$-膜的端点定义了 $C$ 上的一个除子 $\sum_{k}P_k$。\cite{FGT} 论证：这一红外对偶表现为超对称配分函数之间的关系；借助 AGT 对应的各种推广，这些关系可翻译为仿射代数 $\widehat{\mathfrak{sl}}_{2,k}$ 与 $\mathfrak{Vir}_c$ 的共形块之间已知的关系。$\mathfrak{Vir}_c$ 的相关共形块是由在点 $P_1,\dots,P_h$（$h=4g-4+n$）处插入额外的、满足 $b^2=\hbar$ 的退化表示 $V_{-1/2b}$ 来定义的。这些关系来自 Sklyanin 分离变量（SOV）方法的一个推广。

\cite{FGT} 所考虑框架的一个简单变体是：在背景中再放入一个额外的“原初”（primordial）M2-膜，把它定域在 $P_0\in C$ 处并沿 $x_7$ 轴延伸，以此修改初始的 M5-M5'-膜系统。此时在红外对偶描述中，除 $P_1,\dots,P_h$ 处的插入之外，还会在 $P_0$ 处得到额外插入的一个退化表示 $V_{-1/2b}$。如 \cite{T10} 所解释的，这无非就是当在 $P$ 处额外插入 $\widehat{\mathfrak{sl}}_{2,k}$-模 $W_{1/2\hbar}$ 时、由 SOV 得到的 $\widehat{\mathfrak{sl}}_{2,k}$ 共形块的那个表示。

\subsection{更多证据}

另一种论证可以通过对 \cite{AGGTV} 中论证的适当推广得到：那里用退化场的插入来描述额外 M2-膜的作用。加入这样一个顶点算符，会使瞬子配分函数的空间再张量上 $\mathbb{CP}^1$ sigma 模型二维的真空向量空间。余四维面算符位置的变动，会在这个二维向量空间上生成一个单值表示。取 $\ep_2\ra 0 $ 的极限，会导致面算符上的自由度与体（bulk）中的自由度发生部分退耦，从而使瞬子配分函数分解为一个四维的奇异部分与一个二维的正则部分。$\widehat{\mathfrak{sl}}_{2,k}$-模 $W_{1/2\hbar}$ 是仿射李代数共形块中满足下述两条性质的唯一候选插入：(i) 生成一个二维单值表示；(ii) 在临界能级极限下，$W_{1/2\hbar}$ 的插入会从共形块中因子化出去。下一节我们将更详细地描述这一论证路线的若干要素。

若能通过局域化计算来检验我们的猜想 \rf{Heckeconj}，那将是很好的；这需要在两方面作推广：一方面是上面所引的关于余二维面算符的结果，另一方面是关于余四维面算符的已知结果（见 \cite{GLPP} 及其中的文献）。

在约化到 $N=4$ SYM 时，人们会用 $N=4$ SYM 中的 Wilson 与 't Hooft 线算符来表示余四维面算符 \cite{Wi12}。这些线算符与几何 Langlands 对应的 Hecke 函子之间的关系，已在 \cite{KW} 中作了广泛讨论。在 Cordova 与 Jafferis 于 \cite{CJ} 中所研究的约化方案中，将得到复 Chern--Simons 理论的 Wilson 线，这应能为 \rf{Heckeconj} 中所表述的猜想提供另一种视角。

最后可以指出，Beilinson 与 Drinfeld 工作中的 Hecke 函子，是由把 $\mathfrak{g}_k$ 的模 $W_{-(k+h^{\vee})\la}$ 插入共形块时所取的临界能级极限来描述的，其中 $\la$ 是 ${}^LG$ 的一个有限维表示的权，综述见 \cite[Section 20.5]{FBZ}。\rf{Heckeconj} 中出现的那些模则对应于 $\la=1/2$ 的情形，即 ${}^LG=\mathrm{PSL}(2,\C)$ 的基本表示。

\bigskip
\noindent\rule{\textwidth}{0.6pt}

% [Tutor Note: 式 \rf{Heckeconj} 的模与结尾 Beilinson--Drinfeld 公式的对应必须逐项核对：由 $k=-h^{\vee}-\ep_2/\ep_1$ 得 $k+h^{\vee}=-\ep_2/\ep_1=-1/\hbar$，即 $-(k+h^{\vee})=1/\hbar$，故 $\la=1/2$ 时 $-(k+h^{\vee})\la=\dfrac{1}{2\hbar}$，与 $W_{1/2\hbar}$ 相符——其中只用到本文的约定 $\hbar=\ep_1/\ep_2$。若换用文献中同样流行的 $\hbar=\ep_2/\ep_1$（并把能级写成 $k=-h^{\vee}-1/\hbar$），则同一个模会被写成 $W_{\hbar/2}$，于是 $1/(2\hbar)$ 与 $\hbar/2$ 这两套写法在两种约定下指的是同一个对象，直接照抄别人的公式时极易把指数写反。凡比较形如 $W_{\pm}$ 的下标，先固定 $\hbar$ 的定义再代换。]
% [Tutor Note: 退化表示 $V_{-1/2b}$ 的约定同样易错。取 Liouville/Virasoro 的标准记法 $\al_{m,n}=\frac{1-m}{2}b+\frac{1-n}{2}\frac{1}{b}$、$c=1+6(b+1/b)^2$，则 $V_{-1/2b}$ 对应 $\al_{1,2}=-1/(2b)$，而 $\al_{2,1}=-b/2$ 给出的是 $V_{-b/2}$；二者在 $b\ra1/b$（即 $\ep_1\leftrightarrow\ep_2$）下互换。正文给出 $b^2=\hbar$，故 $b\ra\infty$ 对应 $\ep_2\ra0$、$\hbar\ra\infty$、$k\ra-\infty$（临界能级），此时 $\al_{1,2}\ra0$，退化场趋于恒等场——这正是下面性质 (ii)“在临界能级极限下因子化”的另一种说法。]
% [Tutor Note: 表与正文的对应关系有两个易混点。其一，M5 与 M5' 在方向 7 上都没有支撑，因此沿 $x_7$ 把它们分开之后才可以悬挂 M2；M2 行给出的正是 0、1、7 三个方向，其中 0、1 是它在 $\CM^4$ 中留下的二维支撑，7 是被拉开的方向。其二，脚注中的 $S^1_1,S^1_2$ 是定义 $\Om$-形变（参数 $\ep_1,\ep_2$）的两个旋转对称，而第 4 节里余四维缺陷所缠绕的是 $S^1_\CS\subset\CM^4$ 与 $S^1_C\subset C$；把 $S^1_1$ 直接等同于 $S^1_\CS$ 缺乏依据，脚注的下标只在 $\CM^4$ 内部有意义。]
% [Tutor Note: “$W_{1/2\hbar}$ 是满足 (i)(ii) 的唯一候选”这一唯一性论断只是启发式的：性质 (i) 里的“二维单值表示”来自 $\mathbb{CP}^1$ GLSM 的两个真空，其单值群为 $\mathrm{GL}(2)$；性质 (ii) 则要求该模在临界能级极限下从共形块中分裂出来。二者合起来只说明“该模必须是二重退化且在临界能级退耦”，并不构成范畴层面的唯一性定理。真正的检验应如作者所言用局域化计算完成，而这一步在本文中并未进行，引用该猜想时不宜表述为已证。]
% [Tutor Note: 插入点数目 $h=4g-4+n$ 不要与 $\dim_{\C}\Bun_{SL_2}=3g-3+n$（即参数 $x$ 的个数）混为一谈；原文没有解释 $h$ 的来历，也没有给出 $\mathfrak{Vir}_c$ 的中心荷 $c$ 与 $k$、$\hbar$ 的关系式。读者若要把此处的共形块与第 3 节的 $\mathcal{CB}_{\rm te}(C,\hfg_k,\CE_x)$ 对齐，需要先补上 AGT/Liouville 字典中 $c=1+6(\sqrt{\hbar}+1/\sqrt{\hbar})^2$ 这一步，注意本文的 $b^2=\hbar$ 约定下 $c$ 随 $\hbar\ra\infty$ 发散，与 $\ep_2\ra0$ 的临界能级极限一致。]

%=====================================================================
% 第 6 节  重新得到几何 Langlands 对应
% 原文：Section 6. Recovering the geometric Langlands correspondence
%       （arXiv:1702.06499v2）
%=====================================================================
\section{重新得到几何 Langlands 对应}\label{Hecke1}

\begin{keytakeaways}
\begin{itemize}
  \item \textbf{物理图像：}带面算符的 $\CS$ 类理论在 Nekrasov--Shatashvili 极限 $\ep_2\ra0$ 下，
        其配分函数的指数按 $1/\ep_2$ 发散；发散系数 $\mathcal{Y}_\si$ 是 opers 的 Lagrangian
        子流形（在复 Fenchel--Nielsen 坐标下）的生成函数，而指数中的有限部分给出 Hitchin
        本征方程的解 $\Psi$。这正是几何 Langlands 对应在物理上的具体体现。
  \item \textbf{数学工具：}仿射李代数共形块在临界水平极限下的渐近展开、opers 与复
        Fenchel--Nielsen 坐标、Hitchin 纤维化与 Hitchin Hamiltonians，以及 Cecotti--Vafa 的
        $tt^*$ 方程（其 $tt^*$-联络是平坦的 $SL(2,\mathbb{C})$-联络）。
  \item \textbf{关键结论：}余四维面算符的加入等价于在 $t$ 处插入退化场；$\ep_2\ra0$ 时配分函数
        按 $\psi_v(t)$ 因子化，这正是 Hecke 本征值性质。$\psi_v(t)$ 的单值性定义一个 oper
        局部系 $\chi$，$tt^*$-联络与之规范等价。
  \item \textbf{技术要点：}$\ep_2\ra0$ 与 CFT 临界水平极限的对应由 $k=-h^{\vee}-\ep_2/\ep_1$
        实现；$tt^*$-联络要成为 oper 联络，须其在基 $\Pi_i^v$ 下的反全纯部分幂零，且由水平性
        条件给出的二阶方程在 $C_\tau$ 上无奇点。于是 Beilinson--Drinfeld 的 $\CD$-模由
        $D_r-E_r$ 生成的理想之商给出。
\end{itemize}
\end{keytakeaways}

\bigskip
\noindent\rule{\textwidth}{0.6pt}
\bigskip

本节的目标是详细阐述：在面算符存在的情形下，AGT 对应的各种推广形式与几何 Langlands
对应之间的关系。我们将论证，在 Nekrasov--Shatashvili 极限 $\ep_2\ra 0$ 下会出现一种直接的
关联。为此我们先指出：配分函数 $\CZ_\si(a;x,\tau;\ep_1,\ep_2)$ 的渐近行为中含有若干解，
它们满足描述几何 Langlands 对应中 $\CD$-模结构的 Hitchin 本征方程。然后我们讨论：
如何从带面算符的 $\CS$ 类理论的 Nekrasov--Shatashvili 极限这一角度，来理解相应
$\CD$-模的关键性质——Hecke 本征值性质。这些现象在以 $\CM_H(G)$ 为目标流形的 $2d$ sigma
模型中的可能解释，将在下一节讨论。

\subsection{Nekrasov--Shatashvili 极限}

极限 $\ep_2\ra 0$ 依 \cite{NS} 而称为 Nekrasov--Shatashvili 极限。关系式
$k=-h^{\vee}-\ep_2/\ep_1$ 把它等同于 CFT 中的临界水平极限。在 CFT 中已知 \cite{T10,FGT}：
仿射李代数共形块的手征配分函数在临界水平极限下表现为
\begin{equation}\label{NSlim}
\mathcal{Z}_\si(a;x,\tau;\ep_1,\ep_2)\sim
e^{-\frac{1}{\ep_2}\mathcal{Y}_\si(a;\tau;\ep_1)}
\Psi(a;x,\tau;\ep_1)\,\big(1+\CO(\ep_2)\big)\,,
\end{equation}
其中
\begin{itemize}
\item $\mathcal{Y}_\si(a;\tau;\ep_1)$
是 opers 的 Lagrangian 子流形 ${\rm Op}_{\fsl_2}(C_\tau)$ 的生成函数，该对象已在 \cite{NRS}
的一个相关语境中考虑过；而 $a$ 是 ${\rm Op}_{\fsl_2}(C_\tau)$ 的复 Fenchel--Nielsen
坐标，由沿 $C$ 上决定裤分解 $\si$ 的一组曲线的和乐之迹定义。
\item $\Psi(a;x,\tau;\ep_1)$ 是 Hitchin Hamiltonians 的本征函数，$D_r\Psi=E_r\Psi$，其本征值
$E_r=E_r(a;\tau;\ep_1)$ 由 $\mathcal{Y}_\si(a;\tau;\ep_1)$ 按
$E_r=\partial_{\tau_r}\mathcal{Y}_\si(a;\tau;\ep_1)$ 给出。
\end{itemize}
式\eqref{NSlim} 的行为可看作几何 Langlands 对应的一种具体体现：特征标簇中复
Fenchel--Nielsen 坐标为 $(a,\pa_a\mathcal{Y}_\si(a;\tau;\ep_1))$ 的点，可以表示为
$\ep_1$-opers $(\CE_{o},\nabla'_{\ep_1,u})=
 \big(\CE_{o},\big(\ep_1\pa_z+ \big(
 \begin{smallmatrix} 0 & u \\ 1 & 0\end{smallmatrix}\big)\big)dz\big)$ 的和乐。
对一个满足 $u=u^{\sharp}+\sum_{r=1}^{3g-3+n}E_r\vartheta_r(x)$ 的 oper（这里 $u^\sharp$ 定义某个
固定的参考 oper，而 $\vartheta_r(x)(dz)^2$，$r=1,\dots,3g-3+n$ 构成 $H^0(C,K^2)$ 的一组基），
我们可以联系起方程组 $D_r\Psi=E_r\Psi$，$r=1,\dots,3g-3+n$，它定义了几何 Langlands
对应中与 oper $(\CE_{o},\nabla'_{\ep_1,u})$ 相联系的 $\CD$-模。

形如式\eqref{NSlim} 的行为，从带面算符的 $\CS$ 类理论的角度是容易理解的。
在极限 $\ep_2\ra 0$ 下，人们会预期：沿一个在 $\ep_2\ra 0$ 时仍保持有效紧致的子流形
$\CM^2$ 修改规范场的边界条件，不会影响该极限下的领头奇异行为。

\subsection{CFT 中的 Hecke 本征值性质}\label{Hecke-CFT}

当两类面算符同时存在时，可以定义配分函数的进一步推广
$\CZ_v(a;x,\tau;t;\ep_1,\ep_2)$：除此前考虑过的全部变量之外，它额外依赖于一个变量 $t$，
而 $t$ 可以等同于 $C$ 上一点的坐标，即一条“原初”（primordial）M2-膜的端点。上文已论证：
这类配分函数在 CFT 中表示为共形块的手征配分函数，并在 $t$ 处多出一个退化场的插入。
额外的标号 $v$ 指涉某个基的元素，该基张成给出 $4d+2d$ 拉格朗日描述的 GLSM 的真空空间。
在固定一种裤分解的选取之后，我们暂时略去标号 $\si$。

由 CFT \cite{T10,FGT} 可知，配分函数 $\CZ_v(a;x,\tau;t;\ep_1,\ep_2)$ 在极限 $\ep_2\ra 0$ 下表现为
\begin{equation}\label{NSlim+}
\CZ(a;x,\tau;t;\ep_1,\ep_2)\sim
e^{-\frac{1}{\ep_2}\mathcal{Y}(a;\tau;\ep_1)}
\Psi(x;a,\tau;\ep_1)\,\psi_v(t;a,\tau;\ep_1)\big(1+\CO(\ep_2)\big)\,,
\end{equation}
其中 $\mathcal{Y}(a;\tau;\ep_1)$ 与 $\Psi(x;a,\tau;\ep_1)$ 就是式\eqref{NSlim} 中出现的那两个
函数，而 $\psi_v(t)=\psi_v(t;a,\tau;\ep_1)$ 张成微分方程
$(\ep_1^2\partial_t^2-u(t))\psi_v(t)=0$ 之解空间的一组基，这里 $u=u(t)$ 把该 oper 联络表示为
$\ep_1\partial_t+\big(\begin{smallmatrix} 0& u \\ 1 & 0\end{smallmatrix}\big)$。
指标 $v=+,-$ 标记解空间的一组基。

$\psi_v(t)$ 之所以能作为一个因子被分离出来，与 Hecke 本征值性质有关，这一点已在
\cite{T10} 中指出。与前面一样，我们可以把 $\mathcal{Z}_v^*$ 解释为重叠
$\langle \de_x^{\sst (2)},\xi_a^{}\rangle$，其中 $\xi_a\in\CH^{\sst (2)}_a$。若式\eqref{NSlim+}
对所有 $x$ 成立，则对一切 $\xi_a\in\CH^{\sst (2)}_a$ 都有类似的因子化。此外，因子
$\psi_v(t)$ 并不依赖于为把配分函数表为 $\langle \de_x^{\sst (2)},\xi_a^{}\rangle$ 所选用的
向量 $\xi_a\in\CH^{\sst (2)}_a$。把 $\psi_v(t)$ 解析延拓所生成的单值性定义一个 oper
局部系 $\chi$。由于它不依赖于 $\xi_a\in\CH^{\sst (2)}_a$，人们可以把这一改动归因于空间
$\CH^{\sst (2)}_a$ 自身的改动；于是该空间满足 Hecke 本征值性质。

\subsection{由 $tt^*$ 方程导出 Hecke 本征值性质}\label{Hecke-4d}

本节的目标是详细展开 Hecke 函子的 $(4d+2d)$ 描述。由于 $\Om$-形变把涨落局域化到圆作用的
固定点，人们可以把极限 $\ep_2\ra 0$ 看作把半球 $B^4_{\ep_1\ep_2}$ 退紧化为
$B^2_{\ep_1}\times\BR^2$，其中 $B^2_{\ep_1}$ 是一个二维半球。我们所关心的面算符将支撑在
$B^2_{\ep_1}$ 上。

目标是给式\eqref{NSlim+} 的行为一个四维的解释。回想：配分函数
$\mathcal{Z}_v(a;x,\tau,t;\ep_1,\ep_2)$ 可以解释为带 $\Om$-形变的规范场论在四维半球
$B^4_{\ep_1\ep_2}$ 上的配分函数，其定义用到了这样的边界条件：它把规范场路径积分在
$B^4_{\ep_1\ep_2}$ 上生成的状态投影为 Wilson 圈的本征态 \cite{NW,TV13}。变量 $a$ 参数化
Wilson 圈的本征值，而该 $4d$ 超对称规范场论的紫外规范耦合的集合记作 $\tau$。变量 $x$
参数化定义余二维面算符的边界条件 \cite{FGT}，而变量 $t$ 表示 GLSM 的 FI 参数 $t$，
该 GLSM 实现额外的余四维面算符 \cite{AGGTV}。

\subsubsection{因子化}

当 $\ep_2\ra 0$ 时，面算符上的自由度不可能对自由能的发散有贡献，因为此极限下出现的
二维椭球最终具有有限体积。因此我们预期 $\CF=\log\CZ$ 表现为
\begin{equation}
\mathcal{F}\sim-\frac{1}{\ep_2}\mathcal{Y}^{\rm 4d}(a;\tau;\ep_1)+\mathcal{W}^{\rm 2d}(a,\tau;x,t;\ep_1)+\dots\,,
\end{equation}
其中 $\mathcal{W}^{\rm 4d}(a;\tau;\ep_1)$ 与二维耦合 $x$、$t$ 无关。与式\eqref{NSlim} 比较给出
\[
\mathcal{W}^{\rm 2d}(a,\tau;x,t;\ep_1)=\mathcal{W}^{\rm\sst M2}(a,\tau;t;\ep_1)
+\mathcal{W}^{\rm\sst M5'}(a,\tau;x;\ep_1)\,,
\]
其中 $\mathcal{W}^{\rm\sst M2}(a,\tau;t;\ep_1)=\log\psi_v(t;a,\tau;\ep_1)$，
$\mathcal{W}^{\rm\sst M5'}(a,\tau;x;\ep_1)=\log\Psi(x;a,\tau;\ep_1)$。这种解耦大概可以这样
理解：定义余二维面算符的边界条件只涉及体场，而不涉及用来描述余四维面算符的自由度。

于是很自然地把 $\mathcal{Z}^{\rm 2d}_v(a,\tau;t;\ep_1):=\psi_v(t)$ 解释为：GLSM 位于下方的
二维半球 $B^2_{\ep_1}$ 上、带 $\Om$-形变、耦合到体内的 $\CS$ 类理论、并服从记作
$\mathfrak{b}_v$ 的边界条件时的配分函数。这类配分函数的例子近来在 \cite{FKNO} 中研究过，
那里发现它与 \cite{HO,HR} 所研究的半球配分函数有密切关系。人们观察到
$\mathcal{Z}^{\rm 2d}_v(a,\tau;t;\ep_1)$ 对 FI 参数 $t$ 的依赖是全纯的。此外，该配分函数与
GLSM 的拓扑扭曲版本中的一个重叠有关 \cite{HO,HR}：
\begin{equation}\label{Zoverlap}
\mathcal{Z}^{\rm 2d}_v(a,\tau;t;\ep_1)\,=\,\langle\,\mathfrak{b}_v\,|\,0\,\rangle\,,
\end{equation}
其中 $|\,0\,\rangle$ 是半球上无算符插入时的路径积分所生成的状态，而
$\langle\,\mathfrak{b}_v\,|$ 是与边界条件 $\mathfrak{b}_v$ 相联系的边界态。

\subsubsection{$tt^*$ 方程}

重叠 $\langle\,\mathfrak{b}_v\,|\,0\,\rangle$ 有一个自然的推广：把 $|\,0\,\rangle$ 换成另一个
超对称基态 $|\,1\,\rangle$，后者由手征环生成元 $v$ 从 $|\,0\,\rangle$ 得到，
$|\,1\,\rangle= v|\,0\,\rangle$。\cite{HIV} 中已指出，$\Pi_i^v:=\langle\,\mathfrak{b}_v\,|\,i\,\rangle$
表示 Cecotti 与 Vafa \cite{CV} 的 $tt^*$-联络的一个水平截面。在我们的情形中，只有一个耦合
$t$ 和一个二维的超对称真空空间。因此 $tt^*$-联络具有平坦 $SL(2,\mathbb{C})$-联络的形式
$\nabla'=(\partial_t+\mathcal{A}_t)dt$，$\nabla''=(\bar\partial_{\bar{t}}+\mathcal{A}_{\bar{t}})d{\bar{t}}$
\cite{Ga09}。由 \cite{FKNO,HO,HR} 中观察到的全纯性可得：在所讨论的情形中，
$\mathcal{A}_{\bar{t}}$ 在基 $\Pi_i^v$ 下必须幂零。这是存在一个保持 $\Pi_0^v$、并把联络
$\mathcal{A}$ 与某个 oper 相联系的规范变换的必要条件。水平性条件
$\nabla' \Pi^v=0=\nabla''\Pi^v$ 于是给出 $\Pi_0^v$ 的一个二阶微分方程。以上观察为
\cite{DFKMMN} 中关于 opers 的描述提供了一个物理实现。

由 \cite{AGGTV} 中给出的论证可得，耦合 $t$ 可以等同于 Riemann 面 $C_\tau$ 上的局部坐标。
从物理上应当预期：由水平性条件得到的这个常微分方程在 $C_\tau$ 上应当没有奇点。这将意味着
$tt^*$-联络必须规范等价于一个 oper 联络。\cite{FKNO,HO,HR} 的结果支持这一论断。

\subsubsection{固定单值性}

当 $C_\tau$ 的复结构处在 Riemann 面模空间的边界附近——在那里 $C_\tau$ 可表示为一组裤
（pairs of pants）用长管 $T_1,\dots,T_{3g-3+n}$ 连接起来——$\CS$ 类四维理论存在弱耦合
的拉格朗日描述。若 $t$ 位于某条特定的管 $T_r$ 上，则耦合的 $4d$–$2d$ 系统是弱耦合的。
该拉格朗日描述中含有 GLSM 与四维体理论之间 $ta_r/\ep_1$ 形式的耦合，其中 $a_r$ 是与 $T_r$
相联系的多重态中标量场在余四维面算符支撑上的限制 \cite{AGGTV}。因此，与 $t\ra t+2\pi$
相应的配分函数 $\mathcal{Z}^{\rm 2d}_v(a;\tau;t;\ep_1)$ 的单值性是对角的，并可由乘以
$e^{2\pi va_r/\ep_1}$ 表示。这固定了 $tt^*$-联络单值性的一半；而由于该联络规范等价于一个
oper 联络，这一半信息已足以把它完全固定下来。

于是 $tt^*$-联络被等同于几何 Langlands 对应一侧出现的 oper 局部系 $\rho$。把 ${\rm Bun}_G$
上整体微分算子代数对由 $D_r-E_r$（$r=1,\dots,3g-3+n$）生成的理想取商，便定义了 Langlands
对应另一侧的相应 $\CD$-模 $\De_{\rho}$ \cite{Fr07}。因子化式\eqref{NSlim+} 可以看作 Hecke
本征值性质在与 $\CD$-模 $\De_{\rho}$ 相联系的微分方程之解的层面上的体现。

\bigskip
\noindent\rule{\textwidth}{0.6pt}
% [Tutor Note: 临界水平的比例约定（易错）。关系式 $k=-h^{\vee}-\ep_2/\ep_1$ 必须与第 5 节的 $b^2=\hbar$ 以及第 7 节的 $\hbar=\ep_1/\ep_2$ 这一套约定配套理解：于是 $k=-h^{\vee}-1/b^2$，$\ep_2\ra0$ 即 $\hbar\ra\infty$、$b\ra\infty$，$k\ra-h^{\vee}$，确实是临界水平。若有人取相反约定 $b^2=\ep_2/\ep_1$，就会把临界水平错配到 $\ep_1\ra0$ 的极限，并连带把式\eqref{NSlim} 指数中的 $1/\ep_2$ 误写成 $1/\ep_1$。另须注意：临界水平的发散因子是 $\frac{1}{k+h^{\vee}}=-\hbar=-\ep_1/\ep_2$，因此 $\mathcal{Y}$ 必须按 $\mathcal{Y}=\ep_1\cdot(\text{经典 Yang--Yang 函数})+\CO(\ep_1^2)$ 归一化，才能使 $\mathcal{Y}/\ep_2$ 在同时标度 $\ep_i\ra\lambda\ep_i$ 下不变；核对办法是取 $\ep_1\ra0$ 应回到 \cite{NRS} 的经典 Yang--Yang 函数。]
% [Tutor Note: $E_r=\pa_{\tau_r}\mathcal{Y}_\si$ 极易被误读为 $E_r=\pa_{a_r}\mathcal{Y}_\si$。两点提示：(i) 本节后文明确以 $\tau$ 记“该 $4d$ 超对称规范场论的紫外规范耦合的集合”，故 $\tau_r$ 是第 $r$ 个裤分解分量的耦合，而 $\pa_a\mathcal{Y}$ 给出的是与长度坐标 $a$ 共轭的复 Fenchel--Nielsen 坐标（即上文 $(a,\pa_a\mathcal{Y})$ 中的第二个分量）；(ii) 该式是经典 Matone 关系 $u=\pa\mathcal{F}/\pa\tau$ 在 NS 极限下的类比，其成立依赖 $\mathcal{Y}$ 的归一化（即式\eqref{NSlim} 中 $1/\ep_2$ 的约定）。若误取 $E_r=\pa_{a_r}\mathcal{Y}$，则 $u=u^\sharp+\sum_r E_r\vartheta_r$ 不再是 Hitchin 基上的坐标，式\eqref{NSlim} 的“几何 Langlands 具体体现”一说也就随之失效。]
% [Tutor Note: 原文在 $tt^*$ 一节跳过了关键一步。$\mathcal{A}_{\bar{t}}$ 在基 $\Pi_i^v$ 中幂零，正是“存在保持 $\Pi_0^v$ 的规范变换把 $tt^*$-联络化为 oper 形式”的条件：oper 联络的 $(0,1)$ 部分必须是严格上三角的幂零矩阵（在 $\Pi_0^v$ 为最高权向量的规范下），这样水平性条件给出的二阶方程才不含 $\bar\partial_{\bar t}$ 项、成为 $C_\tau$ 上真正的二阶常微分方程 $(\ep_1^2\pa_t^2-u)\Pi_0^v=0$。而“该方程在 $C_\tau$ 上无奇点”这一步，原文只以“从物理上应当预期”带过；数学上它对应 \cite{DFKMMN} 中 oper 与 Hitchin 方程（$SL(2)$ 自对偶方程在共形极限下）之间的等价，需另行验证，不能视为已证。]
% [Tutor Note: 单值性因子 $e^{2\pi va_r/\ep_1}$ 需核对三处。(1) 没有 $i$：$t$ 是复化（全纯）的 FI 参数，耦合 $ta_r/\ep_1$ 对 $t$ 全纯，故 $t\ra t+2\pi$ 给出实指数的乘性因子，而非 $e^{2\pi i(\cdots)}$；若把 $t$ 写成 $t=\theta+i\xi$ 型实/虚分解，这个 $2\pi$ 对应 $\theta$ 方向的周期。(2) 分母是 $\ep_1$ 而非 $\ep_2$：出现在耦合里的 $\ep_1$ 来自沿面算符支撑方向的 $\Om$-形变参数，与式\eqref{NSlim} 指数中的 $\ep_2$ 是两个不同的参数，不可互换。(3) 原文把解空间基的指标 $v=\pm$ 当作整数荷 $v=\pm1$ 写进指数，这是记号混用（见本小节前文“指标 $v=+,-$ 标记解空间的一组基”）；若要严格，应说明该荷来自 $W_{1/2\hbar}$ 插入的权。此外“固定一半单值性即足以完全固定 $tt^*$-联络”隐含了 $SL(2,\mathbb{C})$ 平坦联络由单值表示（模共轭）决定这一 Riemann--Hilbert 事实，并须先假设该 oper 联络确由这些单值性实现。]
% [Tutor Note: 本节原文有三处笔误/记号不一致，译文一律保留原样，请读者注意。(a) 式(6.3) 写的是 $\mathcal{Y}^{\rm 4d}$，紧接着的一句却称“$\mathcal{W}^{\rm 4d}(a;\tau;\ep_1)$ 与二维耦合无关”——按上下文该处应为 $\mathcal{Y}^{\rm 4d}$，$\mathcal{W}$ 与 $\mathcal{Y}$ 混用了。(b) 式\eqref{NSlim} 中写作 $\Psi(a;x,\tau;\ep_1)$，而式\eqref{NSlim+} 中写作 $\Psi(x;a,\tau;\ep_1)$，自变量次序不一致（正文又说二者“就是”同一函数）。(c) 式\eqref{NSlim+} 左端漏掉了下标 $v$，而右端 $\psi_v$ 带 $v$。另，“面算符上的自由度不可能贡献自由能发散，因为二维椭球体积有限”这一论证只是启发式的：$1/\ep_2$ 极点实际来自四维体理论在 $\Om$-背景下的局域化（每个瞬子贡献在 $\ep_2=0$ 处的极点），而面算符部分给出的是 $\CO(\ep_2^0)$ 的正规项。]

%=====================================================================
% 第 7 节  对 sigma 模型的诠释？
% 原文：Section 7. Sigma model interpretation?（arXiv:1702.06499v2）
%=====================================================================
\section{对 sigma 模型的诠释？}

\begin{keytakeaways}
\begin{itemize}
  \item \textbf{物理图像：}把 AGT 与几何 Langlands 的图像放回到以 $\CM_H(G)$ 为目标流形的
        二维拓扑 sigma 模型中。Kapustin--Witten 用的是以 $\omega_K$ 为辛结构的 A-模型、
        其 cc-brane 的 Chan--Paton 曲率为 $\omega_J$；Nekrasov--Witten 用的是以
        $\frac{\ep_2}{\ep_1}\omega_I$ 为辛结构的 A-模型。二者之间没有明显参数可以连续移动，
        这正是本节要化解的障碍。
  \item \textbf{数学工具：}Hitchin 模空间的超 Kähler 旋转（Hitchin 的 $\BC^*$ 作用把
        $(I_\zeta,\Omega_\zeta)$ 族中除 $\zeta=0,\infty$ 外的复结构与辛结构互相联系起来）、
        SYZ 镜对称与 Hitchin 纤维化、非阿贝尔 Hodge（NAH）对应及其共形极限、以及
        $\ep_2\ra0$ 下的 WKB 近似。
  \item \textbf{关键结论：}若 sigma 模型 TQFT 在 $\BC^*$ 作用下具有适当形式的不变性，则取
        $\zeta=\ep_2\xi$ 即可把 NW 与 KW 的两个 A-模型插值起来；在 NS 极限下 brane
        $\mathsf{L}_{\si,a}^{\sst(1)\hbar}$ 退化为支撑在 Hitchin 第一纤维化纤维
        $\mathsf{F}_{(u_{\mathsf{a}},0)}$ 上、Chan--Paton 丛平凡的 (B,A,A)-brane，于是
        $\ep_2\ra0$ 时共形块与 Kapustin--Witten 的 Hecke 本征 brane 之间的对应可写成显式同构。
  \item \textbf{技术要点：}该同构的自然形式是把 Beilinson--Drinfeld 进路中的 oper
        $\rho_u$ 经 NAH 对应的共形极限（$R\ra0$、$\zeta\ra0$、$\zeta/R$ 有限）送到 Hitchin
        模空间的 $(u,0)$ 点；对一般局部系 $\chi$ 则需假设 $\CM_{dR}({}^LG)$ 与
        $\CM_H({}^LG)$ 之间存在由该共形极限相联的分层。
\end{itemize}
\end{keytakeaways}

\bigskip
\noindent\rule{\textwidth}{0.6pt}
\bigskip

上一节的讨论揭示了（面算符存在时）AGT 对应的各种推广与几何 Langlands 对应的
Beilinson--Drinfeld 进路之间的密切联系。我们已经把几何 Langlands 对应识别为：
用“量子化的特征标簇”描述共形块空间这一描述在 $\ep_2\ra 0$ 时的极限。
在这最后一节中，我们要回到这样一个问题：由此浮现的图像与 Kapustin--Witten 进路下的
几何 Langlands 对应之间是什么关系。在 \ref{comp-I} 小节中我们提出如何在 $2d$ sigma 模型中
描述极限 $\ep_2\ra 0$ 的问题；\ref{Cstaraction} 与 \ref{SecNSlim} 两小节给出这一问题的一个
推测性回答。本节其余部分将基于类似的论证，对 Beilinson--Drinfeld 与 Kapustin--Witten
两种进路作更明确的比较。

\subsection{比较 I：AGT 对应与几何 Langlands 对应之间的关系}\label{comp-I}

我们曾论证：此语境中出现的 $\CD$-模可以用相应微分方程的解空间来刻画，而这些解在物理上
实现为配分函数。在带面算符的 AGT 对应语境中，我们曾论证空间
$\CH_{x}^{\sst(2)}:=
{\rm Hom}_{\CM_H(G)}^{}(\mathsf{B}_{\rm cc}^{},\mathsf{L}_{x}^{\sst(2)})$
可以等同于仿射代数共形块在丛 $\CE_x$ 上的纤维。在 CFT 中，我们可以把解 KZB 方程的
手征配分函数 $\CZ(a;x,\tau;k)$ 看作共形块空间 $\mathcal{CB}(C,\hfg_k,\CE_x)$ 的子空间
$\mathcal{CB}^\si_a(C,\hfg_k,\CE_x)$ 中元素 $\vf_a$ 的代表元，而这些子空间由这样的共形块
组成：它们在 Verlinde 线算符的某个极大对易子集下具有由 $a$ 参数化的固定本征值。

开 sigma 模型 TQFT 的模不变性提示我们可以等同
$\mathcal{CB}^\si_a(C,\hfg_k)\simeq 
{\rm Hom}_{\CM_H(G)}^{}(\mathsf{B}_{\rm cc}^{},\mathsf{L}_{\si,a}^{\sst(1)})$。
这一等同可以由下述观察支持：边界条件 $\mathsf{L}_{\si,a}^{\sst(1)}$ 固定了条带上
sigma 模型中某些零模的取值，而这应当导致相对于 $\mathcal{A}_{B}^{\hbar}$ 的某个子代数的
本征值性质，该子代数由相应的量子化迹函数生成，而这些迹函数则与决定裤分解 $\si$ 的
那些曲线相联系。

在上一节中我们讨论了子空间 $\mathcal{CB}_a^\si(C,\hfg_k,\CE_x)$ 在极限 $\ep_2\ra 0$ 下如何
与几何 Langlands 对应中出现的 Hecke 本征 $\CD$-模相联系，其中代表“本征值”的局部系由
以 $a$ 为参数的 opers 族表示。然而不清楚的是：我们关于 Hecke 本征值性质在极限
$\ep_2\ra 0$ 下出现的观察，如何从 sigma 模型的角度来理解。似乎存在一个直接的障碍：
\cite{KW} 在几何 Langlands 对应语境中考虑的 sigma 模型，是以辛结构 $\omega_K$ 的 A-模型，
而 cc-brane 的 Chan--Paton 曲率为 $\omega_J$。在 AGT 对应的 Nekrasov--Witten 进路中则做了
不同的选择，那里用的是以辛结构 $\frac{\ep_2}{\ep_1}\omega_I$ 的 A-模型。没有任何明显的参数
能使我们在这两种情形之间连续移动。

鉴于在 CFT 中可以在临界水平极限下得到几何 Langlands 对应，很自然地要问：这能否在以
$\CM_H(G)$ 为目标流形的 $2d$ sigma 模型内理解。从 \ref{Cstaraction} 小节起，我们将推测
一条可能的途径。

\subsection{超 Kähler 旋转}\label{Cstaraction}

一个诱人的做法是：用来自 Hitchin 模空间上圆作用 $\vf\ra e^{i\theta}\vf$ 的超
Kähler$^{(g)}$ 旋转，来修改 Nekrasov 与 Witten 的二维 TQFT 构架。\cite[Section 9]{Hi87}
中证明：这一圆作用可以复化为一个 $\BC^*$-作用，它把族 $(I_\zeta,\Omega_\zeta)$ 中除
$\zeta=0,\infty$ 之外的所有复结构与辛结构互相联系起来。超 Kähler 旋转使我们能够在
$\zeta\neq 0,\infty$ 时识别 $\CM_H(G)$ 上的关键几何结构，其含义是这些旋转可以由
$\CM_H(G)$ 的微分同胚实现 \cite[Proposition (9.1)]{Hi87}。这相当于对 sigma 模型场作
重新参数化，而后者也许并不影响物理内容。

sigma 模型路径积分在场重定义下的不变性可能是个微妙的问题。圆作用
$\vf\ra e^{i\theta}\vf$ 应当容易理解；但复化为 $\BC^*$-作用之后，它是否在适当意义上
容易理解为 sigma 模型 TQFT 的对称性，我们并不清楚。我们想指出的是：若以 $\CM_H(G)$
为目标流形的拓扑 sigma 模型在 $\BC^*$-作用下具有某种适当形式的不变性，我们就能简单地
在 \cite{KW} 与 \cite{NW} 中分别考虑的那两个 sigma 模型之间插值。这样就能对上一节的观察
给出一个有吸引力的解释。

我们关心的是带边曲面上的开 A-模型，其边界条件或属典范余迷向（canonical coisotropic）型，
或属拉格朗日型。我们可以用全纯辛形式族 $\Omega_{\zeta}$ 来定义一族辛结构为
$\mathrm{Im}(\Omega_\zeta)$ 的 A-模型，其中余迷向 brane 的 Chan--Paton 曲率为
$\mathrm{Re}(\Omega_\zeta)$。$\CM_H(G)$ 中对 $\Omega_\zeta$ 为拉格朗日的子流形，在这族
A-模型中定义 brane。若 sigma 模型 TQFT 具有适当形式的 $\BC^*$-不变性，则可以预期：
在这族 A-模型中定义的配分函数 $\CZ_\zeta$ 实际上与 $\zeta$ 无关。

然而应当注意：我们用来定义边界条件的所有拉格朗日子流形都必须协调地变动，才能保持
$\CZ_{\zeta}$ 不变。我们所关心的子流形是特征标簇中复 Fenchel--Nielsen$^{(g)}$ 扭转流的
轨道，经由 NAH 对应与和乐映射视作 $\CM_H(G)$ 的子流形。固定特征标簇中的一条轨道，
随着超 Kähler 参数 $\zeta$ 的变动，就得到 $\CM_H(G)$ 的一族单参数子流形。

若已有 $\BC^*$-作用下的不变性，我们就可以把耦合参数 $\ep_1/\ep_2$ 的变动与适当的超 Kähler
旋转结合起来，从而得到一族以 $\CM_H(G)$ 为目标流形、并在 \cite{NW} 与 \cite{KW} 所考虑的
两类模型之间插值的拓扑 A-模型。这可以通过取 $\zeta=\ep_2\xi$ 来实现。表达式
$\hat{\omega}_\zeta=\frac{\ep_2}{\ep_1}\mathrm{Im}(\Omega_{\zeta})$ 在极限
$\ep_2\ra 0$ 下约化为 $\hat{\omega}_\zeta=
\frac{1}{2\ep_1|\xi|^2}
(\mathrm{Re}(\zeta)\omega_K-\mathrm{Im}(\zeta)\omega_J)$，
对实 $\xi$ 恰好重现用来定义 \cite{KW} 所研究的 A-模型的辛形式。

\subsection{brane $\mathsf{L}_{\si,a}^{\sst (1)}$ 的 Nekrasov--Shatashvili 极限} \label{SecNSlim}

为看清由此得到的方案是否真的能解决上述困惑，先注意：\ref{branes1} 小节所考虑的 (A,B,A)
brane $\mathsf{L}_{\si,a}^{\sst (1)}$ 可以很容易地推广为 $(B,A,A)_\zeta$-brane
$\mathsf{L}_{\si,a}^{\sst (1)\zeta}$，只要在它们的定义中使用依赖于 $\zeta$ 的 NAH 对应。
如上面建议的那样以依赖于 $\hbar$ 的方式选取 $\zeta$，就会把它们变成超 Kähler 参数为
$\hbar=\ep_1/\ep_2$ 的那族 A-模型中的 brane $\mathsf{L}_{\si,a}^{\sst (1)\hbar}$。

我们会看到，在 brane 族 $\mathsf{L}_{\si,a}^{\sst (1)\hbar}$ 的 Nekrasov--Shatashvili 极限中
发生了有趣的事情：若在此极限中把参数 $a$ 按
${a}_r=\frac{\ep_1}{\ep_2}\,\check{a}_r$ 标度（$\check{a}_r$ 有限），则 brane
$\mathsf{L}_{\si,a}^{\sst(1)\hbar}$ 有良定义的极限 $\hbar\ra \infty$。为理解极限
$\ep_2\ra 0$，主要需要研究 $\pa_u+\mathcal{A}_u$ 的和乐的 WKB 近似，其中
$\mathcal{A}_u=A_u+\frac{1}{\ep_2}\vf$。通过规范变换总可以局部达到形式
$\mathcal{A}_u=\big(\begin{smallmatrix} 0 & u_{\ep_2} \\ 1 & 0\end{smallmatrix}\big)$。
出现在 $\mathcal{A}_u$ 右上矩阵元中的函数 $u_{\ep_2}(t)$ 一般只是 $u$ 的亚纯函数，
但其留量为 $\ep_2$ 阶。我们注意到 $A_u$ 只出现在展开的次领头阶。由这些事实立即得到
$\vartheta_{\ep_2}(t)=\frac{1}{2}{\rm tr}(\vf^2)+\CO(\ep_2)$。因此在 $\ep_2$ 的领头阶，
可以把 $(\pa_u+\mathcal{A}_u)s_{\ep_2}=0$ 的解 $s$ 写成
\begin{equation}
s_{\ep_2}(t)\,\sim\, e^{\pm\frac{1}{\ep_2}\int_{\CC_t} \la}\,,
\end{equation}
其中 $\CC_t$ 是谱曲线 $\Sigma$ 上终止于 $t\in \C$ 的某个提升 $\hat{t}\in\Sigma$ 的路径，
而 $\la$ 是 $\Sigma$ 上的典范微分。结论是：标度后的 Fenchel--Nielsen 长度坐标
$\check{a}_r$ 在极限 $\ep_2\ra 0$ 下表现为
$
\check{a}_r=\frac{1}{\ep_1}\mathsf{a}_r+\CO(\ep_2),
$
其中 $\mathsf{a}_r$ 是微分 $\la$ 沿一组由 $\si$ 确定的典范同调基的某个拉格朗日子空间所取的
周期。

在这个意义上，人们可以把“量子周期”$\check{a}_r$ 看作 Hitchin 可积系统的角变量
$\mathsf{a}_r$ 的形变。有趣的是：对 $\si$ 的依赖在极限 ${\ep}\ra 0$ 下消失——即使坐标
$\mathsf{a}_r$ 的定义带有对 $\si$ 的残余依赖，由 $\mathsf{a}_r$ 的取值所确定的 Hitchin
纤维化的纤维却并非如此。

这样我们便被引向如下结论：brane $\mathsf{L}_{\si,a}^{\sst(1)\hbar}$ 的标度极限由
(B,A,A)-brane 表示，后者支撑在 Hitchin（第一）纤维化的纤维
$\mathsf{F}_{\sst (u_{\mathsf{a}},0)}$ 上，其上的 Chan--Paton 丛是平凡丛，而
$u_\mathsf{a}$ 是 Hitchin 基 $\CB$ 上由坐标 $\mathsf{a}_r$ 确定的点。

最后要注意：也可以用类似方式讨论 brane $\mathsf{L}_{x}^{\sst(2)\hbar}$。容易看出，
当 ${\hbar}\ra 0$ 时得到支撑在 Hitchin 第二纤维化的纤维上的 brane
$\mathsf{L}_{x}^{\sst(2)}$。

\subsection{比较 II：Kapustin--Witten 进路与 CFT 的关系？}\label{compII}

任何在 Beilinson--Drinfeld 与 Kapustin--Witten 两种进路之间的比较，都必须回应下面这一点。
在 \cite{KW} 中，Hecke 本征 brane 被描述为 $\CM_H(G)$ 上的摩天楼层（skyscraper sheaves）。
然而，为了在 Hitchin 纤维化上使用 SYZ 镜对称，Kapustin 与 Witten 采用的是适配复结构 $I$ 的
$\CM_H(G)$ 与 $\CM_H({}^L G)$ 之表示，而 Beilinson--Drinfeld 进路则在对应的一侧考虑
${}^LG$-局部系 $(\CE,\nabla')$。为把二者联系起来，需要用到 NAH 对应。若 $\mu$ 表示
$\CM_H(G)$ 的一个点（该空间表为环面纤维化），令 $\check{\mathsf{F}}_{\mu}$ 为 B-模型中
相应的摩天楼层，$\mathsf{F}_{\mu}$ 为 SYZ-对偶的 A-模型中的 Hitchin 纤维。为表述两种进路
之间的一个猜想性关系，我们需要用 NAH 对应（的逆）找出与局部系 $\chi$ 相联系的点
$\mu(\chi)\in\CM_H(G)$。

于是，对 Beilinson--Drinfeld 与 Kapustin--Witten 两种进路之间可能关系的一个自然猜测，
就是下述 $\CD$-模同构成立：
\begin{equation} \label{BD-KW}
\begin{aligned}
{\rm Hom}_{\CM_H(G)}^{}(\mathsf{B}_{\rm cc},\mathsf{F}_{\mu(\chi)}) 
& \;\simeq\; \mathcal{CB}(C,\hfg_{-h^{\vee}},\CE)\;\;\text{for}\;\;\CE\in\mathrm{Bun}_G^{\rm vs}\\
& \;\simeq\; \big[\,
\text{Fiber of}\;\, \De_\chi\;\,\text{over}\;\,\CE\in\mathrm{Bun}_G^{\rm vs}\,\big]
\end{aligned}
\end{equation}
其中 $\mathrm{Bun}_G^{\rm vs}$ 是“非常稳定”（very stable）丛 $\CE$ 的空间，即不容许幂零
Higgs 场的丛；$\De_\chi$ 是由 $\mathcal{CB}(C,\hfg_{-h^{\vee}},\CE)$ 表示的 $\CD$-模。
由于存在一个典范平坦联络把与不同 $\CE\in\mathrm{Bun}_G^{\rm vs}$ 相联系的纤维等同起来，
右端不依赖于 $\CE$ 的选取。$\mathcal{CB}(C,\hfg_{-h^{\vee}},\CE)$ 的维数在离开
$\mathrm{Bun}_G^{\rm vs}$ 时可能发生跳跃，讨论见 \cite[Section 9.5]{Fr07}。

在此我们还可以指出 Beilinson--Drinfeld 与 Kapustin--Witten 进路比较中出现的另一个困惑。
如上文所指出的、并由 \cite{T10} 所研究的例子所说明的：在 Beilinson--Drinfeld 进路中出现的
$\CD$-模，在离开 $\CM_{dR}({}^LG)$ 内的 opers 子簇时表现出某种不连续的行为，其表现是：
在这些 $\CD$-模以共形块形式给出的表示中，会出现若干个额外的退化表示。而在 Kapustin 与
Witten 的进路中看不到这种不连续行为。

\subsection{共形极限}

我们还愿意指出一条也许能为 \ref{compII} 小节提出的问题提供答案的途径。它基于
\cite{Ga14,DFKMMN} 中的观察：在共形极限下 NAH 对应可能大大简化；该极限是把通过
$\vf\ra R\vf$ 引入 NAH 对应的参数 $R$ 与超 Kähler 参数 $\zeta$ 一起趋于零，同时保持
$\zeta/R$ 有限。作为准备，这里先讨论该极限可能具有的相关性。所涉及的问题与
\ref{Cstaraction} 小节中讨论的超 Kähler 旋转的相关性问题类似，但并不相同，因此值得
单独讨论。

作替换 $\vf\ra R\vf$ 定义了 NAH 对应的一族单参数形变，看起来会修改 $\CM_H(G)$ 上的
超 Kähler 度量。然而，参数 $R$ 对 Hitchin 模空间的几何而言是非本质的，其含义是：
不同 $R$ 值对应的超 Kähler 度量之间由微分同胚相联系。人们可以用由修改后的 NAH 对应
所得的 $\CM_H(G)$ 上的超 Kähler 度量，来定义一族单参数 sigma 模型作用量 $S_R$。
但是，既然这些作用量是由以微分同胚相联的超 Kähler 度量构造的，人们或会倾向于把仅在
参数 $R$ 的选取上不同的两个作用量 $S_{R_1}$ 与 $S_{R_2}$，认同为同一个 sigma 模型量子场论
之物理等价的拉格朗日表示。在诸如 $u=\frac{1}{2}\mathrm{tr}(\vf^2)$ 这样的合适坐标中，
可以发现相应的微分同胚实现为简单的标度变换 $u\ra ur^{2}$。相应的场重定义确实只应以
非本质的整体因子来重新标度配分函数。

若用作用量 $S_R$ 定义的配分函数 $Z_R(u)$ 依赖于一个边界参数 $u$，而该参数对应于
$\CM_H(G)$ 的一个坐标、在 $\vf\ra R\vf$ 下按 $u\ra ur^{2}$ 标度，则我们发现：sigma 模型
度量在微分同胚意义下的 $R$-无关性表现为 $Z_{r R}(u)\propto Z_R(ur^{2})$，同样可能只差
非本质的整体因子。

边界 B-模型中由 $\CM_{dR}({}^LG)$ 上的点 $\chi$ 定义的边界条件所对应的配分函数
$Z'(\chi)$，可以用依赖于点 $\mu_{R,\zeta}(\chi) \in \CM_{H}({}^LG)$ 的配分函数
$Z_R''(\mu_{R,\zeta}(\chi))$ 来表示，因为由 NAH 对应之修改所产生的 $R$-依赖，被 sigma 模型
作用量的相应修改所补偿。若存在使 $\mu_{R,\zeta}(\chi)$ 大大简化的极限，这一观察可能有用。
我们将看到，这样的极限正是所谓的共形极限 $R\ra 0$、$\zeta\ra 0$ 且保持 $\zeta/R$ 有限。
如下文所述，它将给出所提议的关系式\eqref{BD-KW} 的一个有趣的重述。

\subsection{更明确地表述 Kapustin--Witten 进路与 CFT 的关系}

我们将主要限于这样的情形：局部系 $\chi$ 为 opers
$\rho_u=\big(\CE_o,\big(\pa_t+\big(\begin{smallmatrix} 0& u\\ 1 & 0\end{smallmatrix}\big)\big)dt\big)$，
它们在 Beilinson 与 Drinfeld 的进路中起基本作用。有趣的是：在此情形下 NAH 对应的共形极限
变得特别简单 \cite{DFKMMN}，它把 opers $\rho_u$ 与一类形式非常特殊的 Higgs 对联系起来。
对 ${}^L G=\mathrm{PSL}(2,\BC)$，相应的 Higgs 对 $(\mathcal{E},\vf)$ 形如
$(K^{1/2}\oplus K^{-1/2},(\begin{smallmatrix} 0& u \\ 1 &0\end{smallmatrix})dz)$，其中 $u$
是在 $C$ 的 Fuchs 单值化中表示 $\rho_u$ 的那个二次微分。若把 $\CM_H({}^L G)$ 描述为由
$(\Sigma_u,\CL)$ 这样的对表示的环面纤维化（$\Sigma_u$ 是与 Hitchin 基上一点 $u\in\CB$ 相联系的谱曲线，
$\CL$ 是 $\Sigma_u$ 上的线丛），则得到线丛 $\CL_0=\pi^*(K^{1/2})$，其中 $\pi: \Sigma_u\ra C$
是覆盖投影。该线丛代表参数化 $\CL$ 各种选择的 Jacobian/Prym 的一个典范“原点” \cite{W15}。

令 $\check{\mathsf{F}}_{\sst(u,0)}$ 为 $\CM_{H}({}^LG)$ 上支撑在点 $(\Sigma_u,\CL_0)$ 的
摩天楼层，并令 $\mathsf{F}_{\sst (u,0)}$ 为与 $\check{\mathsf{F}}_{\sst(u,0)}$ SYZ 对偶的
Hitchin 纤维化的纤维。于是，Beilinson--Drinfeld 与 Kapustin--Witten 两种进路之间的可能
关系式\eqref{BD-KW} 可以更明确地写成
\begin{equation} \label{BD-KW-conf}
{\rm Hom}_{\CM_H(G)}^{}(\mathsf{B}_{\rm cc},\mathsf{F}_{\sst (u,0)})
\;\simeq\;
\big[\,
\text{Fiber of}\;\, \De_{\rho_u}\;\,\text{over}\;\,\CE\in\mathrm{Bun}_G^{\rm vs}\,\big]\,.
\end{equation}
我们指出：$\CM_{H}({}^LG)$ 中一般点 $(u,0)$ 在 NAH 对应下的像，当且仅当从
$\CM_{H}({}^LG)$ 到 $\CM_{dR}({}^LG)$ 的映射是用非阿贝尔 Hodge 对应的共形极限定义时，
才由某个 oper 表示。\footnote{“if”这一方向已在 \cite{DFKMMN} 中证明。
$(u,0)\in\CM_{H}({}^LG)$ 中的点在 NAH 下的像由具有实和乐的联络组成，这些联络只离散地与
opers 簇相交。我们感谢 A. Neitzke 的这一说明。} 这一点，以及该极限在 \ref{Hecke-4d}
小节中的相关性，表明这一极限很适合用来表述 Beilinson--Drinfeld 与 Kapustin--Witten
两种进路之间的关系。

关于把 \eqref{BD-KW-conf} 推广到更一般的局部系 $\chi$，我们猜想：存在
$\CM_{dR}({}^LG)$ 与 $\CM_{H}({}^LG)$ 的自然分层，其各层之间由 NAH 对应的共形极限
相联系。这将使我们能把关系式\eqref{BD-KW-conf} 推广到一般的不可约局部系，从而把几何
Langlands 对应中 $\CD$-模的不连续行为与从一个层过渡到另一个层联系起来。我们计划在后续
文章中回到这一点。

\subsection{展望}

我们将在别处讨论形如
\begin{equation}\label{spectral}
\CH_{x}^{\sst(2)}\simeq \int^{\oplus} d\mu_\si(a)\;\CH_{\si,a}^{\sst(1)}\,,
\end{equation}
的关系之存在的现有证据，以及关于与不同裤分解 $\si$ 相联系的各空间
$\CH_{\si,a}^{\sst(1)}$ 之间存在线性关系的证据。这就把较弱版本的 $\si$-无关性恢复到了
与非零 $\epsilon_2$ 相联系的那套图景之中。

几何 Langlands 对应有时被表述为经典 Langlands 纲领中出现的自守形式空间之谱分解的类比。
我们把本短文的内容视为一些提示：梦想一个由超越（transcendental）与分析（analytic）两方面
扩展的几何 Langlands 纲领的变体，或许并不算放肆。超越方面涉及配分函数，它们表示由
$\CD$-模所定义的微分方程组的解；分析方面则涉及如 \eqref{spectral} 所提议的谱分解。
把配分函数认同为自守形式的类比物，会进一步加强对经典 Langlands 纲领的类比。配分函数
构成了 $\CM_H(G)$ 的代数结构之间的桥梁，而这些代数结构分别与 $\CM_H(G)$ 作为局部系
模空间 $\CM_{dR}(G)$ 的表示、以及作为特征标簇 $\CM_B(G)$ 的表示相联系。由此可以期望得到
一个更大的图像，把几何 Langlands 纲领中拓扑的方面与依赖复结构的方面统一起来。

我们计划在后续文章中讨论这些内容、“量子几何 Langlands 对偶”这一解释，以及与 \cite{T10}
中称为模对偶（modular duality）的另一种 Langlands 对偶模式的化身之间的关系。

最后注意，从规范场论方面对几何 Langlands 纲领的近期进展已在 \cite{Ga16a,Ga16b} 中描述。
分析它们与我们工作的关系应当是有意义的。

{\bf 致谢：}
A.B 感谢 D. Ben-Zvi、M. Mulase、A. Neitzke 与 R. Wentworth 的讨论，并感谢 String-Math 2016
的组织者筹办了一场富有启发的会议。A.B 还感谢 L'Institut Henri Poincar\'e（巴黎）与
ICTS-TIFR（班加罗尔）在工作进行期间的接待。

J.T. 感谢 String-Math 2016 的组织者筹办了一场鼓舞人心的会议并提供报告本文工作的机会，
感谢 M. Mulase、A. Neitzke 的讨论。特别感谢 E. Frenkel 的多方讨论、告知其未发表工作
\cite{Fr10} 的内容，以及对草稿的批评性意见。

本工作由德国科学基金会（DFG）通过协作研究中心 SFB 676 “Particles, Strings and the Early
Universe” 项目 A10 资助。

\bigskip
\noindent\rule{\textwidth}{0.6pt}
% [Tutor Note: 超 Kähler 旋转这一步最易被误读。Hitchin 的 $\BC^*$-作用对实 $\theta$ 是 $U(1)$ 等距，但复化之后**不是等距**：\cite[Proposition (9.1)]{Hi87} 只断言 $\zeta\neq0,\infty$ 时超 Kähler 旋转由 $\CM_H(G)$ 的微分同胚实现，而该微分同胚会改变度量（它是“旋转复结构”而非“保持度量”）。因此“这相当于对 sigma 模型场作重新参数化，也许并不影响物理内容”是一个需要证明的断言，而非平凡事实：sigma 模型的作用量依赖目标度量，一般的目标空间微分同胚（非等距）会改变作用量；此外场重定义还会移动 $B$ 场与路径积分测度，除非同时作相应补偿且无反常，物理才真正不变。原文自己也承认“may be a subtle issue”。若要在 KW 与 NW 两个 A-模型之间插值，需要的正是这一不变性；而 A-模型只依赖辛形式，故真正要验证的是 $\mathrm{Im}(\Omega_\zeta)$（连同 $\mathrm{Re}(\Omega_\zeta)$ 作为 Chan--Paton 曲率）在 $\BC^*$-作用下的行为。]
% [Tutor Note: 第 7.2 小节末的 $\hat\omega_\zeta$ 极限式（原文未编号）请用附录中的公式复核。附录给出 $\Omega_\zeta=\frac{1}{2\zeta}(\omega_J+i\omega_K)+i\omega_I+\frac{\zeta}{2}(\omega_J-i\omega_K)$。代入 $\zeta=\ep_2\xi$ 得 $\frac{1}{\zeta}=\bar\xi/(\ep_2|\xi|^2)$，故 $\Omega_\zeta\simeq\frac{1}{2\ep_2|\xi|^2}\big[(\mathrm{Re}\xi\,\omega_J+\mathrm{Im}\xi\,\omega_K)+i(\mathrm{Re}\xi\,\omega_K-\mathrm{Im}\xi\,\omega_J)\big]+i\omega_I$，于是 $\mathrm{Im}\Omega_\zeta=\frac{1}{2\ep_2|\xi|^2}(\mathrm{Re}\zeta\,\omega_K-\mathrm{Im}\zeta\,\omega_J)+\omega_I$。乘上 $\ep_2/\ep_1$ 后，$\omega_I$ 项给出 $\frac{\ep_2}{\ep_1}\omega_I$，它在 $\ep_2\ra0$ 时被丢掉——原文未曾明言这一丢弃，而这一项正是“NW 的 $\omega_I$ 模型”与“KW 的 $\omega_K$ 模型”之间差别的残留。另注意对实 $\xi$ 只得到正比于 $\omega_K$ 的辛形式，整体因子为 $1/(2\ep_1\xi)$，故与 \cite{KW} 的 $\omega_K$ 只相差一个常数标度；常数标度在 A-模型中只改变整体因子，不影响 brane 的拉格朗日性。]
% [Tutor Note: WKB 与标度这一步有三处易错。(1) 指数里是 $\la$ 而不是 $u$ 或 $\sqrt{u}$：$\la$ 是谱曲线 $\Sigma\subset T^*C$ 上的典范微分，$\pm$ 对应两叶，$\CC_t$ 是到提升点 $\hat t$ 的路径；$1/\ep_2$ 来自 $\mathcal{A}_u$ 中 $\vf$ 前的 $1/\ep_2$，故 WKB 的大参数是 $\hbar=\ep_1/\ep_2$（在 $\ep_2\ra0$、$\ep_1$ 固定时趋于无穷），不要误写成 $1/\ep_1$。(2) 这是奇异极限，$e^{\pm\frac{1}{\ep_2}\int\la}$ 只在各 Stokes 区域内渐近有效，绕行转折点（$\la=0$ 处）时按 Stokes 矩阵跳跃，真正的解须经 Borel 求和才确定；原文直接写“解 $s$ 的形式”跳过了 Stokes 现象，而 \cite{GMN2} 表明这一跳跃正对应谱曲线的墙穿越（WKB 曲线的 BPS 态计数），它会影响下文 $\check a_r$ 的 $\CO(\ep_2)$ 修正。(3) 标度关系可疑：把 $a_r=\frac{\ep_1}{\ep_2}\check a_r$ 与 $\check a_r=\frac{1}{\ep_1}\mathsf a_r+\CO(\ep_2)$ 联立得 $a_r=\frac{1}{\ep_2}\mathsf a_r+\CO(\ep_1)$，令 $\ep_2\ra0$ 会发散；而若 $\mathsf a_r$ 确如正文所说是 $\la$ 的周期、且经典长度坐标即该周期，则更自然的写法应为 $\check a_r=\frac{\ep_2}{\ep_1}\mathsf a_r+\CO(\ep_2)$。原文此处 $\ep_1,\ep_2$ 的归属值得核对（译文保留原样）。此外下一句把 $\mathsf a_r$ 称为 Hitchin 可积系统的“角变量”，但“$\la$ 的周期”给出的是作用量（action）变量，与 NRS 及本附录中作为 Darboux 坐标的长度坐标 $a_r$ 同一角色；角变量应是其共轭的扭转（twist）坐标。]
% [Tutor Note: 共形极限可直接用附录的关系式核对，这是本节最容易被跳过的一步。附录给出 $\nabla_{\zeta,R}=\zeta^{-1}R\vf+D_{\CE,h}+R\zeta\vf^{\dagger_h}$，以及 $\ep=\zeta/R$、$\nabla'_\ep=\ep\pa_{\CE,h}+\vf$。取 $R\ra0,\ \zeta\ra0$、$\zeta/R=\ep$ 固定，则 $R\zeta=\zeta^2/\ep\ra0$ 而 $\zeta^{-1}R=R/\zeta=1/\ep$，故 $\nabla_{\zeta,R}$ 的 $(1,0)$ 部分趋于 $\pa_{\CE,h}+\frac{1}{\ep}\vf$，即 $\ep\nabla'\ra\ep\pa_{\CE,h}+\vf$ 恰为 $\ep$-联络（oper 型）。同时自对偶方程 $F_{\CE,h}+R^2[\vf,\vf^{\dagger_h}]=0$ 退化为 $F_{\CE,h}=0$，即 $h$-酉联络平坦、和乐为实。这恰好解释了脚注中“$(u,0)$ 在 NAH 下的像有实和乐、只离散地与 opers 相交”：不用共形极限时 $(u,0)$ 并非 oper；只有在共形极限下 $\vf$ 的 Higgs 部分才“存活”为 oper 联络的 $u$，$(u,0)$ 的像才落到 opers 上。换言之，式\eqref{BD-KW-conf} 里的 $\rho_u$ 与 $(u,0)$ 之间必须插入“共形极限的 NAH”这一环节，不能直接用通常的 NAH。注意别把 $\ep=\zeta/R$ 与 $R/\zeta$ 弄混：$\ep$ 小时 $R/\zeta=1/\ep$ 是大的，这正是 $\frac{1}{\ep}\vf$ 这一项。]
% [Tutor Note: 关于式\eqref{BD-KW} 的三点提醒。(1) 必须限制到 $\mathrm{Bun}_G^{\rm vs}$（不含幂零 Higgs 场的“非常稳定”丛）：那里 $\De_\chi$ 的纤维是有限维的，左端 $\mathrm{Hom}_{\CM_H(G)}(\mathsf B_{\rm cc},\mathsf F_{\mu(\chi)})$ 作为两个拉格朗日 brane 的交（一般为横截的有限集）也有限维，同构才可能成立；原文自承“维数可能跳跃”，见 \cite[Section 9.5]{Fr07}——在非非常稳定丛处纤维会变成无穷维（额外的退化表示出现），故不要把式\eqref{BD-KW} 读成 $\mathrm{Bun}_G$ 上向量丛的同构。(2) “右端不依赖 $\CE$”用的是共形块丛在 $\mathrm{Bun}_G$ 上的典范平坦联络（KZ/BGG 联络），沿它平移即可等同不同 $\CE$ 处的纤维；这一步与“非常稳定”条件是两回事，勿混。(3) 第 7.6 小节中 Higgs 对的扭变可逐项核对：取 $E=K^{1/2}\oplus K^{-1/2}$，则 $\mathrm{End}(E)\ot K$ 的右上元属于 $\mathrm{Hom}(K^{-1/2},K^{1/2})\ot K=K^2$，故 $u\in H^0(C,K^2)$ 正确；左下元属于 $\mathrm{Hom}(K^{1/2},K^{-1/2})\ot K=\CO$，故常数 $1$ 正确；对角元必须为零——这正是 oper 条件（$\vf$ 关于 filtration 严格下三角）。注意 $K^{\pm1/2}$ 的存在依赖一个自旋结构（$\Theta$-特征）的选取，故“$\CL_0=\pi^*K^{1/2}$ 是典范原点”只在固定该自旋结构后成立；$SL(2)$ 与 $GL(2)$ 的差别（Prym 对 Jacobian）正是这一 $\mathbb{Z}_2$ 数据造成的，附录中“Jacobian 若 $GL(2)$、Prym 若 $SL(2)$”即指此。]
% [Tutor Note: 展望段的式\eqref{spectral} 是个直积分，其中测度 $d\mu_\si(a)$ 原文未给出；按 \cite{T10} 的谱分解，它应是复 Fenchel--Nielsen 长度坐标上的 Plancherel 型测度（含 Liouville 理论的反射振幅/DOZZ 密度），而非平凡的 Lebesgue 测度，否则“$\sigma$-无关性”无法成立。另注意：右端对每个 $\si$ 给出一个分解，而左端与 $\si$ 无关，因此“不同裤分解之间的线性关系”正是使该分解自洽的关键；其经典极限是生成函数为 \cite{NRS} 的 Yang--Yang 函数 $\mathcal Y$ 的辛变换（即 Fenchel--Nielsen 坐标变换），量子化后应是量子膨胀/量子二对数（quantum dilogarithm）型的变换。原文只说“将在别处讨论”，此处不宜过度解读为已证的定理。]


%---------------------------------------------------------------------
% 附录
%---------------------------------------------------------------------
\appendix
%=====================================================================
% 附录  Hitchin 模空间
% 原文：Appendix. Hitchin's moduli spaces（arXiv:1702.06499v2, 附录）
%=====================================================================
% 原文宏 \BC（即复数域 \mathbb{C}）：若共享 preamble 已提供则此处自动跳过，
% 未提供时就地补定义，以保证公式逐字无损且整篇可编译。
\providecommand{\BC}{\mathbb{C}}

\section{Hitchin 模空间}\label{Gloss}

\begin{keytakeaways}
\begin{itemize}
  \item \textbf{物理图像：} Hitchin 模空间 $\CM_H(G)$ 是全文三条进路共同的目标流形兼相空间。Kapustin--Witten 的进路把它取作 $N=4$ 超 Yang--Mills 理论沿 $C$ 约化后所得二维 sigma 模型的目标流形；AGT 对应及其 Nekrasov--Shatashvili（NS）极限把它看作被量子化的 Hitchin 可积系统的相空间，共形块给出 Hitchin 哈密顿量的本征函数；Beilinson--Drinfeld 的 CFT 进路则经由 $\CM_{dR}$ 中的 oper 子空间与它对接。
  \item \textbf{数学工具：} 谱曲线 $\Sigma\subset T^*C$ 连同 Hitchin 纤维化把 $\CM_H(G)$ 实现为底空间 $\CB\simeq H^0(C,K^2)$ 上的代数完全可积系统；非阿贝尔 Hodge（NAH）对应给出 $\CM_H(G)$、$\CM_{dR}(G)$ 与 $\CM_B(G)$ 三种描述的同一性；以 $\mathbb{P}^1$ 参数化的一族复结构 $I_\ze$ 与全纯辛形式 $\Om_\ze$ 构成超 K\"ahler 结构，是统一上述描述的语言。
  \item \textbf{关键结论：} 同一流形上的复结构 $I$、$J$、$K$ 分别对应 Higgs 丛侧、平坦联络（de Rham）侧与基本群表示（Betti，即特征簇）侧。三条进路的差别恰好体现为取用了族中不同的复结构——Kapustin--Witten 的 Hecke 本征膜按复结构 $I$ 下的 Hitchin 纤维化描述，Beilinson--Drinfeld 一侧取的是 $\CM_{dR}$ 中的局部系统——而由 $\vf\ra e^{i\theta}\vf$ 生成的超 K\"ahler 旋转（在 $\ze\neq0,\infty$ 处可由 $\CM_H(G)$ 的微分同胚实现）把三者联系起来。
  \item \textbf{附录的作用：} 正因如此，Hitchin 模空间才是三条进路可以逐项比较的公共几何框架；本附录以术语表形式固定全文记号，条目之间的衔接（谱曲线—纤维化、NAH—三种描述、$\Om_\ze$—辛结构）正是正文论证的几何骨架。
\end{itemize}
\end{keytakeaways}

\bigskip
\noindent\rule{\textwidth}{0.6pt}
\bigskip

本文假定 $G=SL(2)$，且 $C$ 是亏格为 $g$、带 $n$ 个穿刺点的 Riemann 面。

{\bf Hitchin 模空间 $\CM_H(G)$} \cite{Hi87}。
这是偶 $(\CE,\vf)$ 的模空间，其中 $\CE=(E,\bar{\partial}_{\CE})$ 是光滑向量丛 $E$ 上的全纯结构（holomorphic structure），而 $\vf\in H^0(C,\mathrm{End}(\CE)\ot K)$。这类偶模去自然规范变换后所得的模空间记作 $\CM_H(G)$。

{\bf Hitchin 可积系统} \cite{Hi87}。
给定 $(\mathcal{E},\vf)$，可构造谱曲线
$\Sigma=\{(u,v);v^2=\frac{1}{2}\mathrm{tr}(\vf^2)\}\subset T^*C$，
以及代表 $\vf-v$ 的余核的线丛 $\CL$。
由 $(\Sigma, \mathcal{L})$ 可以重构 $(\mathcal{E},\vf)$：
$\mathcal{E}=\pi_*(\mathcal{L})$，$\vf=\pi_*(v)$。
这就把 $\CM_{H}(G,C)$ 描述为底空间
$\CB\simeq H^0(C,K^2)$
上的环面纤维化，其中纤维——即 $\CL$ 的选取——当 $G=GL(2)$ 时等同于 $\Sigma$ 的 Jacobi 簇（Jacobian），当 $G=SL(2)$ 时等同于 Prym 簇（Prym variety）。底空间 $\CB$ 的自然坐标由 Hitchin 哈密顿量给出，其定义为将
$\frac{1}{2}\mathrm{tr}(\vf^2)=\sum_{r=1}^{3g-3+n}\vartheta_rH_r$
展开，其中 $\{\vartheta_r, r=1,\dots,3g-3+n\}$ 构成 $H^0(C,K^2)$ 的一组基。

{\bf 局部系统}。即偶 $(\CE,\nabla_\ep')$，其中 $\CE$ 是如上全纯向量丛，$\nabla_\ep'$ 是全纯 $\ep$-联络，对函数 $f$ 与 $E$ 的光滑截面 $s$ 满足
$\nabla_\ep' (fs)=\ep (\pa f)s+f\nabla_\ep's$。这类偶的模空间记作 $\CM_{dR}(G)$。
此处常把局部系统等同于相应的\emph{平坦丛}（flat bundle），即转移函数为常值的局部平凡化系统；或等同于由 $(\CF,\nabla_\ep')$ 的和乐得到的基本群表示（模去共轭）。由此得到 $\CM_{dR}(G)$ 与

{\bf 特征簇（character variety）$\CM_B(G)$}：$\pi_1(C)$ 到 $G$ 的表示模去整体共轭所成之空间；作为代数簇，它被描述为一个 GIT 商
$\BC[\mathrm{Hom}(\pi_1(C),G]^G$。

{\bf oper}。这是一类特殊的局部系统：其中 $\CE=\CE_{\rm op}$，即唯一可使全纯联络 $\nabla_\ep'$ 取形式
$\nabla_\ep'=dz\big(\ep\pa_z+\big(
\begin{smallmatrix} 0 & u \\ 1 & 0\end{smallmatrix}\big)\big)$
的扩张
$0\ra K^{1/2}\ra\CE_{\rm op} \ra K^{-1/2}\ra 0$。

{\bf 非阿贝尔 Hodge（NAH）对应} \cite{Hi87,Si}。给定 Higgs 偶 $(\CE,\vf)$，存在
唯一的调和度量 $h$ 满足
%\begin{align}\label{selfdual}
$F_{\CE,h}+R^2[\vf,\vf^{\dagger_h}]=0$
%\end{align}
其中 $F_{\CE,h}$ 是唯一以
$(0,1)$-部分为 $\bar{\partial}_{\CE}$ 的 $h$-酉联络 $D_{\CE,h}$ 的曲率。
% In the holomorphic gauge where $\bar{\partial}_{\CE}=\bar{\partial}$,
%we have $D_{\CE,h}=\partial+ h(\partial h^{-1}) +\bar\partial$, and \rf{selfdual} becomes
%\begin{align}\label{selfdual'}
%\bar\partial(h(\partial h^{-1}))+R^2\big[\vf,\bar{h}^{-1}\vf^{\dagger}\bar{h}\big]=0\,,
%\end{align}
%Equation \rf{selfdual} determines $h$ uniquely for given $(\CE,\vf)$.  
于是可以构造相应的双参数平坦联络族
%\begin{align}
$\nabla_{\zeta,R}=\zeta^{-1}{R}\,\vf+D_{\CE,h}+R\zeta\,\vf^{\dagger_h}$。
%\end{align}
把 $\nabla_{\zeta,R}$ 分解为 $(1,0)$ 与 $(0,1)$-两部分 %$\nabla'$ and $\nabla''$ 
即给出一个偶 $(\CF,\nabla_\ep')$，它由
$\CF=(E,\bar{\partial}_{\CF})$ 与 $\ep$-联络
$\nabla_\ep'=\ep\nabla'=\ep\pa_{\CE,h}+\vf$ 组成，其中 $\ep=\zeta/R$，且 $\nabla_\ep'$ 在
$\bar{\partial}_{\CF}$ 所定义的复结构下全纯。

{\bf 超 K\"ahler 结构} \cite{Hi87}。存在由 $\mathbb{P}^1$ 参数化的一族复结构 $I_\zeta$ 与全纯
辛结构 $\Omega_\zeta$。后者定义为
$
\Omega_\zeta =\frac{1}{2}\int_C\;{\rm tr}(\de\mathcal{A}_\zeta\wedge\de\mathcal{A}_\zeta)
$.
三重辛形式 $(\omega_I,\omega_J,\omega_K)$ 可由 $\Omega_\zeta$ 的展开定义：
$
\Omega_\zeta=\frac{1}{2\zeta}(\omega_J+i\omega_K)+i\omega_I+\frac{1}{2}{\zeta}(\omega_J-i\omega_K).
$
相应的复结构为
$
I_{\zeta}=\frac{1}{1+|\zeta|^2}((1-|\zeta|^2)I-i(\zeta-\bar{\zeta})J-(\zeta+\bar{\zeta})K).
$

{\bf 复 Fenchel--Nielsen 坐标} \cite{NRS}。这是与 $C$ 的裤子分解（pants decomposition）$\si$ 相对应的 $\CM_{B}(G)$ 上的 Darboux 坐标，而该分解由沿闭曲线 $\ga_i$（$i=1,\dots,3g-3+n$）切割而得。
复长度坐标把迹函数 $L_i=\mathrm{tr}(\rho(\ga_i))$ 参数化为 $L_r=2\cosh(a_r/2)$。可以定义正则共轭坐标
$\kappa_r$，使得自然的 Poisson 结构表示为
$\{a_r,\kappa_s\}=\de_{r,s}$、$\{a_r,a_s\}=0=\{\kappa_r,\kappa_s\}$。

\bigskip
\noindent\rule{\textwidth}{0.6pt}
% [Tutor Note: 被注释掉的三个公式环境（原文写作 %\begin{align}…%\end{align}）内容如下：其一是 Hitchin 1987 的核心方程 $F_{\CE,h}+R^2[\vf,\vf^{\dagger_h}]=0$，即 Hitchin 方程，亦称自对偶方程（它是 4d 自对偶 Yang–Mills 方程在 $C$ 上的 $G$-不变约化，故原文径称 self-dual）；其二给出全纯规范 $\bar\partial_{\CE}=\bar\partial$ 下解 $h$ 的等价非线性 PDE $\bar\partial(h(\partial h^{-1}))+R^2[\vf,\bar h^{-1}\vf^\dagger\bar h]=0$；其三是 $\nabla_{\zeta,R}$ 的 align 环境。它们被注释的原因应是排版遗留而非数学内容删改：同一段里被引用的 \label{selfdual}、\label{selfdual'} 也一并被注释，注释中的 \rf{selfdual}（=\eqref）因此指向不存在的 label，作者最终改以行内无编号公式呈现。译者严格保持注释状态、不改回可见公式，并按“公式逐字保留”原则保留其行内可见版本。]
% [Tutor Note: $R^2$ 项的来源：$R$ 是 4d$\to$2d 约化中紧化曲面 $C$ 的尺度／耦合参数（Hitchin 原文即以 $R$ 记之）；本文正文 Conformal limit 一节明确它由 NAH 对应中的标度 $\vf\ra R\vf$ 引入，并把保持 $\ze/R$ 有限的极限 $R\ra0$、$\ze\ra0$ 作为几何 Langlands 的关键。核对量纲：$F_{\CE,h}$ 与 $[\vf,\vf^{\dagger_h}]$ 都是 $(1,1)$-形式（$\vf$ 为伴随值 $(1,0)$-形式、取值于 $K$，$\vf^{\dagger_h}$ 为 $(0,1)$-形式），故 $R^2$ 是使两项同阶的无量纲因子；按 $\vf\ra R\vf$ 改写时第二项整体乘 $R^2$，与正文标度约定一致。注意 $\vf^{\dagger_h}$ 显含 Hermitian 度量 $h$，这正是该方程作为 $h$ 的方程、而非代数方程的原因。]
% [Tutor Note: 复结构 $I,J,K$ 与 $\Om_\ze$ 的约定核对：由附录的 $\Om_\ze$ 展开可得 $\ze=0$ 时 $I_\ze=I$、$\ze=i$ 时 $I_\ze=J$、$\ze=-1$ 时 $I_\ze=K$（代入 $1+|\ze|^2$：$\ze=i$ 给出 $I_i=\frac12(-i\cdot2i)J=J$；$\ze=-1$ 给出 $I_{-1}=\frac12\cdot2K=K$）。注意 $\Om_\ze$ 在 $\ze=0$ 有极点，故“$\ze=0$ 处的全纯辛形式”应理解为 $\ze\Om_\ze|_{\ze=0}=\frac12(\omega_J+i\omega_K)$，这才是 $I$-全纯的 $(2,0)$-形式；文献对 $\ze$ 是否吸收这个 $1/\ze$ 归一化不同，比较公式前须核对。另：$\de$ 是模空间（场空间）上的外微分，$\mathcal{A}_\ze$ 为 $\CM_H$ 上对应的联络形式，$\wedge$ 同时含 $C$ 上的形式与场空间微分，$\Om_\ze$ 才是 $\CM_H$ 上的辛形式。]
% [Tutor Note: 稳定性条件与“模空间”的确切含义：原文只说“模去自然规范变换”（modulo natural gauge transformations），但 $\CM_H(G)$ 的严格定义须取 $(E,\vf)$ 的 S-等价类，即 polystable 对象。Hitchin 方程可解 $\Leftrightarrow$ $(E,\vf)$ polystable，此即 Hitchin (1987) 与 Simpson (1992) 的定理；附录中 NAH 对应断言的 $\CM_H(G)\cong\CM_{dR}(G)\cong\CM_B(G)$ 实为 polystable 模空间之间的同构，稳定性条件不可省略。另注意 $\vf$ 是伴随值 $(1,0)$-形式（见 $\mathrm{End}(\CE)\ot K$），$\mathrm{tr}(\vf^2)\in H^0(C,K^2)$ 是二次微分而非函数；对 $G=SL(2)$ 还应取无迹部分 $\mathrm{End}_0$，否则底空间会多出 $H^0(C,K)$ 方向（详见交付报告“原文可疑之处”第 2 条）。]
% [Tutor Note: 谱曲线与 Hitchin 纤维化：$\Sigma=\{(u,v);v^2=\frac12\mathrm{tr}(\vf^2)\}\subset T^*C$ 是 $\vf$ 的特征方程定义的二重覆盖（$v$ 为 $T^*C$ 上的 tautological 1-形式），$\CL=\mathrm{coker}(\vf-v)$ 是 $\Sigma$ 上的线丛；$\CE=\pi_*\CL$、$\vf=\pi_*(v)$ 即“Higgs 丛 = 谱曲线上线丛的直像”。$G=SL(2)$ 时 $\det\CE$ 平凡，故 $\CL$ 须满足 $\mathrm{Nm}(\CL)=\CO_C$，其纤维是 Prym 簇，即 $\ker(\mathrm{Nm}:\mathrm{Pic}(\Sigma)\to\mathrm{Pic}(C))$ 的单位连通分支，维数 $g_\Sigma-g=3g-3+n$，恰等于底空间 $\CB$ 的维数，故 $\dim_\BC\CM_H=2(3g-3+n)$。Hitchin 哈密顿量 $H_r$ 是 $\frac12\mathrm{tr}(\vf^2)$ 在基 $\vartheta_r$ 下的系数；正文 NS 极限中 oper 的 $u=u^\sharp+\sum_rE_r\vartheta_r$ 用的正是同一组 $\vartheta_r$，$E_r$ 为 $D_r\Psi=E_r\Psi$ 的本征值——附录与正文的 $\vartheta_r$ 记号须一并理解。]


%---------------------------------------------------------------------
% 参考文献（由原 .bbl 文件改写为内嵌 thebibliography 环境）
%---------------------------------------------------------------------
%=====================================================================
% 参考文献：由原文 AGT-GL.bbl 改写（内嵌 thebibliography）
%=====================================================================
\renewcommand{\refname}{参考文献}
\begin{thebibliography}{99}

\bibitem{AGT}
L.~F. Alday, D.~Gaiotto, and Y.~Tachikawa, {\it {Liouville Correlation
  Functions from Four-dimensional Gauge Theories}},  {\em Lett.Math.Phys.} {\bf
  91} (2010) 167--197, [\href{http://xxx.lanl.gov/abs/0906.3219}{{\tt
  arXiv:0906.3219}}].

\bibitem{N15}
N.~Nekrasov, {\it {BPS/CFT correspondence: non-perturbative Dyson-Schwinger
  equations and qq-characters}},  {\em JHEP} {\bf 03} (2016) 181,
  [\href{http://xxx.lanl.gov/abs/1512.0538}{{\tt arXiv:1512.0538}}].

\bibitem{KW}
A.~Kapustin and E.~Witten, {\it {Electric-Magnetic Duality And The Geometric
  Langlands Program}},  {\em Commun. Num. Theor. Phys.} {\bf 1} (2007) 1--236,
  [\href{http://xxx.lanl.gov/abs/hep-th/0604151}{{\tt hep-th/0604151}}].

\bibitem{NW}
N.~Nekrasov and E.~Witten, {\it {The Omega Deformation, Branes, Integrability,
  and Liouville Theory}},  {\em JHEP} {\bf 1009} (2010) 092,
  [\href{http://xxx.lanl.gov/abs/1002.0888}{{\tt arXiv:1002.0888}}].

\bibitem{NRS}
N.~Nekrasov, A.~Rosly, and S.~Shatashvili, {\it {Darboux coordinates, Yang-Yang
  functional, and gauge theory}},  {\em Nucl. Phys. Proc. Suppl.} {\bf 216}
  (2011) 69--93, [\href{http://xxx.lanl.gov/abs/1103.3919}{{\tt
  arXiv:1103.3919}}].

\bibitem{Fr07}
E.~Frenkel, {\it Lectures on the Langlands Program and Conformal Field Theory},
   in {\em Frontiers in Number Theory, Physics, and Geometry II} (P.~Cartier,
  P.~Moussa, B.~Julia, and P.~Vanhove, eds.), pp.~387--533.
\newblock Springer Berlin Heidelberg, Berlin, Heidelberg, 2007.

\bibitem{T10}
J.~Teschner, {\it {Quantization of the Hitchin moduli spaces, Liouville theory,
  and the geometric Langlands correspondence I}},  {\em Adv. Theor. Math.
  Phys.} {\bf 15} (2011) 471--564,
  [\href{http://xxx.lanl.gov/abs/1005.2846}{{\tt arXiv:1005.2846}}].

\bibitem{GW}
D.~Gaiotto and E.~Witten, {\it {Knot Invariants from Four-Dimensional Gauge
  Theory}},  {\em Adv. Theor. Math. Phys.} {\bf 16} (2012), no.~3, 935--1086,
  [\href{http://xxx.lanl.gov/abs/1106.4789}{{\tt arXiv:1106.4789}}].

\bibitem{KP}
T.~Kimura and V.~Pestun, {\it {Quiver W-algebras}},
  \href{http://xxx.lanl.gov/abs/1512.0853}{{\tt arXiv:1512.0853}}.

\bibitem{AFO}
M.~Aganagic, E.~Frenkel, and A.~Okounkov, {\it {Quantum q-Langlands
  Correspondence}},  \href{http://xxx.lanl.gov/abs/1701.0314}{{\tt
  arXiv:1701.0314}}.

\bibitem{AG}
D.~Arinkin and D.~Gaitsgory, {\it Singular support of coherent sheaves and the
  geometric Langlands conjecture},  {\em Selecta Mathematica} {\bf 21} (2015),
  no.~1, 1--199.

\bibitem{FF92}
B.~Feigin and E.~Frenkel, {\it {Affine Kac-Moody algebras at the critical level
  and Gelfand-Dikii algebras}},  {\em Int. J. Mod. Phys.} {\bf A7, \rm Suppl.
  1A} (1992) 197--215.

\bibitem{Ar16}
D.~{Arinkin}, {\it {Irreducible connections admit generic oper structures}},
  {\em ArXiv e-prints} (Feb., 2016)
  [\href{http://xxx.lanl.gov/abs/1602.0898}{{\tt arXiv:1602.0898}}].

\bibitem{Ga15}
D.~Gaitsgory, {\it Outline of the proof of the geometric {L}anglands conjecture
  for {$GL_2$}},  {\em Ast\'erisque} (2015), no.~370, 1--112.

\bibitem{Wi09}
E.~Witten, {\it {Geometric Langlands From Six Dimensions}},
  \href{http://xxx.lanl.gov/abs/0905.2720}{{\tt arXiv:0905.2720}}.

\bibitem{GMN2}
D.~Gaiotto, G.~W. Moore, and A.~Neitzke, {\it Wall-crossing, {H}itchin systems,
  and the {WKB} approximation},  {\em Adv. Math.} {\bf 234} (2013) 239--403.

\bibitem{GW1}
D.~Gaiotto and E.~Witten, {\it {Supersymmetric Boundary Conditions in N=4 Super
  Yang-Mills Theory}},  {\em J. Statist. Phys.} {\bf 135} (2009) 789--855,
  [\href{http://xxx.lanl.gov/abs/0804.2902}{{\tt arXiv:0804.2902}}].

\bibitem{GW2}
D.~Gaiotto and E.~Witten, {\it {S-Duality of Boundary Conditions In N=4 Super
  Yang-Mills Theory}},  {\em Adv. Theor. Math. Phys.} {\bf 13} (2009), no.~3,
  721--896, [\href{http://xxx.lanl.gov/abs/0807.3720}{{\tt arXiv:0807.3720}}].

\bibitem{BN16}
D.~{Ben-Zvi} and D.~{Nadler}, {\it {Betti Geometric Langlands}},  {\em ArXiv
  e-prints} (June, 2016) [\href{http://xxx.lanl.gov/abs/1606.0852}{{\tt
  arXiv:1606.0852}}].

\bibitem{CJ}
C.~Cordova and D.~L. Jafferis, {\it {Toda Theory From Six Dimensions}},
  \href{http://xxx.lanl.gov/abs/1605.0399}{{\tt arXiv:1605.0399}}.

\bibitem{Pe}
V.~Pestun, {\it {Localization of gauge theory on a four-sphere and
  supersymmetric Wilson loops}},  {\em Commun. Math. Phys.} {\bf 313} (2012)
  71--129, [\href{http://xxx.lanl.gov/abs/0712.2824}{{\tt arXiv:0712.2824}}].

\bibitem{AT}
L.~F. Alday and Y.~Tachikawa, {\it {Affine SL(2) conformal blocks from 4d gauge
  theories}},  {\em Lett.Math.Phys.} {\bf 94} (2010) 87--114,
  [\href{http://xxx.lanl.gov/abs/1005.4469}{{\tt arXiv:1005.4469}}].

\bibitem{FGT}
E.~Frenkel, S.~Gukov, and J.~Teschner, {\it {Surface Operators and Separation
  of Variables}},  {\em JHEP} {\bf 01} (2016) 179,
  [\href{http://xxx.lanl.gov/abs/1506.0750}{{\tt arXiv:1506.0750}}].

\bibitem{Br}
A.~Braverman, {\it {Instanton counting via affine Lie algebras. 1. Equivariant
  J functions of (affine) flag manifolds and Whittaker vectors}},  in {\em
  {Algebraic Structures and Moduli Spaces, vol. 38 of CRM Proc. Lecture Notes,
  pp. 113--132, Amer. Math. Soc., Providence}}, 2004.

\bibitem{KPPW}
C.~Kozcaz, S.~Pasquetti, F.~Passerini, and N.~Wyllard, {\it {Affine sl(N)
  conformal blocks from N=2 SU(N) gauge theories}},  {\em JHEP} {\bf 1101}
  (2011) 045, [\href{http://xxx.lanl.gov/abs/1008.1412}{{\tt
  arXiv:1008.1412}}].

\bibitem{Ne}
A.~Negut, {\it {Affine Laumon Spaces and the Calogero-Moser Integrable
  System}},  \href{http://xxx.lanl.gov/abs/1112.1756}{{\tt arXiv:1112.1756}}.

\bibitem{Na}
S.~Nawata, {\it {Givental J-functions, Quantum integrable systems, AGT relation
  with surface operator}},  {\em Adv. Theor. Math. Phys.} {\bf 19} (2015)
  1277--1338, [\href{http://xxx.lanl.gov/abs/1408.4132}{{\tt
  arXiv:1408.4132}}].

\bibitem{Fr10}
E.~Frenkel, {\it On the brane interpretation of conformal blocks},  {\em {\rm
  available at http://edwardfrenkel.com/confblocks.pdf}} (2010).

\bibitem{Ka08}
A.~Kapustin, {\it {A Note on Quantum Geometric Langlands Duality, Gauge Theory,
  and Quantization of the Moduli Space of Flat Connections}},
  \href{http://xxx.lanl.gov/abs/0811.3264}{{\tt arXiv:0811.3264}}.

\bibitem{AGGTV}
L.~F. Alday, D.~Gaiotto, S.~Gukov, Y.~Tachikawa, and H.~Verlinde, {\it {Loop
  and surface operators in N=2 gauge theory and Liouville modular geometry}},
  {\em JHEP} {\bf 1001} (2010) 113,
  [\href{http://xxx.lanl.gov/abs/0909.0945}{{\tt arXiv:0909.0945}}].

\bibitem{DGOT}
N.~Drukker, J.~Gomis, T.~Okuda, and J.~Teschner, {\it {Gauge Theory Loop
  Operators and Liouville Theory}},  {\em JHEP} {\bf 1002} (2010) 057,
  [\href{http://xxx.lanl.gov/abs/0909.1105}{{\tt arXiv:0909.1105}}].

\bibitem{FBZ}
E.~Frenkel and D.~Ben-Zvi, {\em {Vertex algebras and algebraic curves}}.
\newblock Providence, USA: Am. Math. Soc. (2004) 400 p, 2004.

\bibitem{DMO}
N.~Drukker, D.~R. Morrison, and T.~Okuda, {\it {Loop operators and S-duality
  from curves on Riemann surfaces}},  {\em JHEP} {\bf 09} (2009) 031,
  [\href{http://xxx.lanl.gov/abs/0907.2593}{{\tt arXiv:0907.2593}}].

\bibitem{GLPP}
J.~Gomis, B.~Le~Floch, Y.~Pan, and W.~Peelaers, {\it {Intersecting Surface
  Defects and Two-Dimensional CFT}},
  \href{http://xxx.lanl.gov/abs/1610.0350}{{\tt arXiv:1610.0350}}.

\bibitem{Wi12}
E.~Witten, {\it {Fivebranes and knots}},  {\em Quantum Topol.} {\bf 3} (2012)
  1--137, [\href{http://xxx.lanl.gov/abs/1101.3216}{{\tt 1101.3216}}].

\bibitem{NS}
N.~A. Nekrasov and S.~L. Shatashvili, {\it {Quantization of Integrable Systems
  and Four Dimensional Gauge Theories}},  in {\em {Proceedings, 16th
  International Congress on Mathematical Physics (ICMP09): Prague, Czech
  Republic, August 3-8, 2009}}, pp.~265--289, 2009.
\newblock \href{http://xxx.lanl.gov/abs/0908.4052}{{\tt arXiv:0908.4052}}.

\bibitem{TV13}
G.~Vartanov and J.~Teschner, {\it {Supersymmetric gauge theories, quantization
  of moduli spaces of flat connections, and conformal field theory}},  {\em
  Adv. Theor. Math. Phys.} {\bf 19} (2015) 1--135,
  [\href{http://xxx.lanl.gov/abs/1302.3778}{{\tt arXiv:1302.3778}}].

\bibitem{FKNO}
T.~Fujimori, T.~Kimura, M.~Nitta, and K.~Ohashi, {\it {2d partition function in
  $\Omega$-background and vortex/instanton correspondence}},  {\em JHEP} {\bf
  12} (2015) 110, [\href{http://xxx.lanl.gov/abs/1509.0863}{{\tt
  arXiv:1509.0863}}].

\bibitem{HO}
D.~Honda and T.~Okuda, {\it {Exact results for boundaries and domain walls in
  2d supersymmetric theories}},  {\em JHEP} {\bf 09} (2015) 140,
  [\href{http://xxx.lanl.gov/abs/1308.2217}{{\tt arXiv:1308.2217}}].

\bibitem{HR}
K.~Hori and M.~Romo, {\it {Exact Results In Two-Dimensional (2,2)
  Supersymmetric Gauge Theories With Boundary}},
  \href{http://xxx.lanl.gov/abs/1308.2438}{{\tt arXiv:1308.2438}}.

\bibitem{HIV}
K.~Hori, A.~Iqbal, and C.~Vafa, {\it {D-branes and mirror symmetry}},
  \href{http://xxx.lanl.gov/abs/hep-th/0005247}{{\tt hep-th/0005247}}.

\bibitem{CV}
S.~Cecotti and C.~Vafa, {\it {Topological antitopological fusion}},  {\em Nucl.
  Phys.} {\bf B367} (1991) 359--461.

\bibitem{Ga09}
D.~Gaiotto, {\it {Surface Operators in N = 2 4d Gauge Theories}},  {\em JHEP}
  {\bf 11} (2012) 090, [\href{http://xxx.lanl.gov/abs/0911.1316}{{\tt
  arXiv:0911.1316}}].

\bibitem{DFKMMN}
O.~{Dumitrescu}, L.~{Fredrickson}, G.~{Kydonakis}, R.~{Mazzeo}, M.~{Mulase},
  and A.~{Neitzke}, {\it {Opers versus nonabelian Hodge}},  {\em ArXiv
  e-prints} (July, 2016) [\href{http://xxx.lanl.gov/abs/1607.0217}{{\tt
  arXiv:1607.0217}}].

\bibitem{Hi87}
N.~J. Hitchin, {\it The self-duality equations on a {R}iemann surface},  {\em
  Proc. London Math. Soc. (3)} {\bf 55} (1987), no.~1 59--126.

\bibitem{Ga14}
D.~Gaiotto, {\it {Opers and TBA}},
  \href{http://xxx.lanl.gov/abs/1403.6137}{{\tt arXiv:1403.6137}}.

\bibitem{W15}
E.~Witten, {\it {More On Gauge Theory And Geometric Langlands}},
  \href{http://xxx.lanl.gov/abs/1506.0429}{{\tt arXiv:1506.0429}}.

\bibitem{Ga16a}
D.~Gaiotto, {\it {S-duality of boundary conditions and the Geometric Langlands
  program}},  \href{http://xxx.lanl.gov/abs/1609.0903}{{\tt arXiv:1609.0903}}.

\bibitem{Ga16b}
D.~Gaiotto, {\it {Twisted compactifications of 3d N = 4 theories and conformal
  blocks}},  \href{http://xxx.lanl.gov/abs/1611.0152}{{\tt arXiv:1611.0152}}.

\bibitem{Si}
C.~Simpson, {\it Higgs bundles and local systems},  {\em Publ. Math. IHES} {\bf
  75} (1992) 5--95.

\end{thebibliography}


%---------------------------------------------------------------------
% 译校说明（译者所加）
%---------------------------------------------------------------------
%=====================================================================
%  译校说明（Translation Notes）
%  本文件为译者所加，非原文内容，置于参考文献之后。
%=====================================================================
\section*{译校说明}
\addcontentsline{toc}{section}{译校说明}

本译文由英文学术论文 \emph{Supersymmetric field theories and geometric Langlands: The other side of the coin}
（A.~Balasubramanian, J.~Teschner, arXiv:1702.06499v2，提交于 2017 年 3 月 10 日，String-Math 2016 会议文集）
精译而成。为便于读者核对与引用，此处说明翻译所依据的底本、处理原则与已知问题。

\subsection*{一、底本与工作方式}

\begin{enumerate}[label=（\arabic*）,leftmargin=*]
  \item \textbf{底本}：以 arXiv 公开的 \LaTeX{} 源码（\texttt{AGT-GL.tex}, 2320 行，及 \texttt{AGT-GL.bbl}）为主要依据，
        并与随任务提供的 PDF（35 页，经逐页核对与 arXiv 版一致）对照。以源码而非 OCR 为准，
        可保证公式、指标、上下标、希腊字母在转换中绝对无损。
  \item \textbf{公式}：所有行内公式与显示公式均逐字保留源码原貌，包括 \verb|\label|、\verb|\left/\right|、\verb|\,|、
        矩阵环境与 \verb|\underset| 等排版细节；未重排、未化简、未“顺手”改符号。原文中被作者注释掉的公式
        （如附录中 Hitchin 方程的自对偶形式）保持注释状态，仅在批注中说明其内容。
  \item \textbf{正文}：按中文学术论文的语序与行文习惯重写，不作逐词硬译；术语以文末《术语译名对照表》为准，
        人名一律保留拉丁字母。
  \item \textbf{批注}：正文源码中以 \verb|% [Tutor Note: ...]| 形式插入译者批注，共 44 条，分“易错提醒”与
        “跳步点拨”两类。批注只出现在源码注释中，不影响排版；在 .tex 源码中可直接看到。
        各节开头另有“本节重点（Key Takeaways）”，为译者所加，非原文内容。
  \item \textbf{编译}：本文档使用 XeLaTeX 编译（中文支持依赖 \texttt{ctex} 宏包）。\texttt{pdfLaTeX} 不能直接编译含中文的版本。
\end{enumerate}

\subsection*{二、图表清单}

\begin{enumerate}[label=（\arabic*）,leftmargin=*]
  \item 经核查，原文 PDF 共 35 页，\textbf{不含任何 Figure/Diagram 图像对象}
        （PDF 内部图像对象计数为 0），故无图可裁、无图可移。
  \item 原文含两张表格，均已原样保留于正文并作了表头翻译：
        第 3 节的 M5/M5$'$ 膜支撑方向表，第 5 节的 M5/M5$'$/M2 膜支撑方向表。
        原文未给表题，故未添加 \verb|\caption|。
\end{enumerate}

\subsection*{三、原文可疑之处与已知笔误}

以下各条均为译者在逐式核对中发现的原文问题。译文一律\textbf{保留原样}，另以批注说明；此处集中列出，供教授复核。

\begin{enumerate}[label=\textbf{[\arabic*]},leftmargin=*]
  \item \textbf{记号不统一（多节）}：原文 $N=4$ 与 $\mathcal{N}=4$、$\mathcal{N}=(4,4)$ 混用；
        ${}^{\rm L}G$ 与 ${}^{L}G$ 混用；术语表上标有 \verb|\glos|、$^{g)}$、$^{(g)}$ 三种写法。
  \item \textbf{§2.5 印刷错误}：geometric Lang\textbf{s}lands（应为 Langlands）；
        “branes \textbf{the} dual sigma model” 缺介词 of。
  \item \textbf{§2.5 逻辑疑点}：Cordova--Jafferis 一段称边界条件位于“$C$ 的两端”，但 $C$ 为紧 Riemann 面、无端点；
        按上下文（$I_0\times C$ 上的复 Chern--Simons）疑应为区间 $I_0$ 的两端。
  \item \textbf{§3 式 (3.3)}：正文只定义 $\mathfrak{t}_+=\mathfrak{t}_1+i\mathfrak{t}_2$，式中却用 $\mathfrak{t}_-$；
        第三条的 $\mathcal{A}_1$ 与前文 $A_y$ 记法不统一。
  \item \textbf{§3}：$\mathcal{A}_{\rm cc}^\hbar=\mathrm{Hom}_{\CM_G(C)}(\cdots)$ 中 $\CM_G(C)$ 全文未定义，
        按上下文应为 $\CM_H(G)$；$T^*\mathcal{T}(C)|_{C}$ 的下标应为 Teichmüller 空间中的点 $[C]$。
  \item \textbf{§3}：“very stable 丛 = 不承认幂零 Higgs 场”漏掉“非零”二字（零 Higgs 场总是存在）。
  \item \textbf{§4 指标错位（系统性）}：$\langle\de_a^{\sst(2)},\phi_x\rangle$ 应为 $\de_a^{\sst(1)}$；
        $\vf_a\in\CH_x^{\sst(1)}$ 应为 $\CH_a^{\sst(1)}$；$\mathsf{L}_{\ep,a}^{\sst(1)\si}$ 多一个 $\ep$ 下标；
        配分函数自变量 $(a;x;\tau)$、$(a;x,\tau)$、$(x,a,\tau)$ 三种写法并存。
  \item \textbf{§4 宏误用}：原文 $S^1_{\CS}$ 处的 \verb|\CS| 在导言区定义为花体 $\mathcal{S}$（class $\mathcal{S}$ 的 S），
        但此处语义为 Chern--Simons；译文按“复 CS（Chern--Simons）理论”处理。
  \item \textbf{§6}：式 (6.3) 写 $\mathcal{Y}^{\rm 4d}$，紧接一句却称 $\mathcal{W}^{\rm 4d}$；属笔误。
  \item \textbf{附录 括号不配对}：$\BC[\mathrm{Hom}(\pi_1(C),G]^G$ 内层圆括号未闭合；
        且不变量环 $\BC[\mathrm{Hom}]^G$ 与 GIT 商 $\mathrm{Hom}/\!/G$ 混写。
  \item \textbf{附录 $\mathrm{End}$ 应为 $\mathrm{End}_0$}：$G=SL(2)$ 时 Higgs 场取值应为无迹自同构丛，
        否则与 $\CB\simeq H^0(C,K^2)$ 及 $GL(2)$ 情形不自洽。
  \item \textbf{附录 谱曲线方程与 $GL(2)$ 讨论错配}：$v^2=\frac12\mathrm{tr}(\vf^2)$ 只在 $\mathrm{tr}\,\vf=0$ 时成立；
        紧随其后的 $G=GL(2)$ 讨论中正确方程应为 $\det(v-\vf)=0$，即 $v^2-(\mathrm{tr}\vf)v+\det\vf=0$。
  \item \textbf{附录 退化未注明}：“由 $(\CF,\nabla_\ep')$ 的和乐得到基本群表示”只在 $\ep\neq0$ 时成立；
        $\ep=0$ 时 $\nabla_0'=\vf$ 退化为 Higgs 场，“和乐”无意义。
  \item \textbf{附录 下标混用}：$L_i=\mathrm{tr}(\rho(\ga_i))$ 用 $i$，随后写作 $L_r=2\cosh(a_r/2)$，$i$ 与 $r$ 混用。
\end{enumerate}

\subsection*{四、阅读建议}

\begin{enumerate}[label=（\arabic*）,leftmargin=*]
  \item 建议先读 §2.1（Beilinson--Drinfeld 进路）、§2.2（Kapustin--Witten 进路）、§2.3（Nekrasov--Witten/AGT），
        再读 §2.4 指出的共同根源（六维 $(2,0)$ 理论在 $\CM^6=\CM^4\times C$ 上的两种紧化）。
        本文的核心主张在 §6 与 §7：带面算符的 AGT 对应之 Nekrasov--Shatashvili 极限给出几何 Langlands 对应。
  \item 附录 “Hitchin's moduli spaces” 是全文术语的汇总，建议与术语译名对照表配合使用。
  \item 源码中 44 条 \verb|% [Tutor Note: ...]| 批注按“易错/跳步”标注，可直接用于组会报告或备课。
\end{enumerate}


%---------------------------------------------------------------------
% 术语对照表
%---------------------------------------------------------------------
%=====================================================================
%  术语对照表（Glossary of Terms）
%  统一译名，保证全文一致。凡表中条目，译文中必须采用对应写法。
%=====================================================================
\section*{术语译名对照表}
\addcontentsline{toc}{section}{术语译名对照表}

下表给出本文使用的统一译名。凡同一术语在中文文献中有多种译法者，本文择定一种并在右栏给出取舍说明；原文首次出现处保留英文原名，以便与文献对照。

\medskip
\noindent
\begin{longtable}{@{}p{0.30\textwidth}p{0.26\textwidth}p{0.36\textwidth}@{}}
\toprule
\textbf{英文原文} & \textbf{本译名} & \textbf{说明} \\
\midrule
\endfirsthead
\toprule
\textbf{英文原文} & \textbf{本译名} & \textbf{说明} \\
\midrule
\endhead
\bottomrule
\endfoot

%----- 场论与规范理论 -----
gauge theory                & 规范场论            & 不作“规范理论” \\
gauge field                 & 规范场              & \\
gauge group                 & 规范群              & \\
gauge transformation        & 规范变换            & \\
Yang--Mills theory          & Yang--Mills 理论    & 人名不译 \\
super-Yang--Mills (SYM)     & 超 Yang--Mills 理论 & 缩写 SYM 保留 \\
supersymmetry (SUSY)        & 超对称              & \\
supersymmetric              & 超对称的            & \\
topological quantum field theory (TQFT) & 拓扑量子场论 & 缩写 TQFT 保留 \\
twisted / topological twist & 拓扑扭曲 / 扭转     & “twist”在超对称语境作“扭转” \\
S-duality                   & S-对偶              & \\
duality                     & 对偶性              & \\
electric-magnetic duality   & 电磁对偶            & \\
Montonen--Olive duality     & Montonen--Olive 对偶 & \\
BPS state                   & BPS 态              & \\
BPS/CFT correspondence      & BPS/CFT 对应        & \\
instanton                   & 瞬子                & 不译“瞬子解” \\
monopole                    & 磁单极子            & \\
't Hooft operator           & 't Hooft 算符       & \\
Wilson operator / line      & Wilson 算符 / 圈   & \\
line operator               & 线算符              & \\
surface operator            & 面算符              & 原文 co-dimension-two defect \\
defect                      & 缺陷                & \\
domain wall                 & 畴壁                & \\
co-dimension two (four)     & 余二维（余四维）    & 不用“余维度二” \\
Nahm pole                   & Nahm 极点           & \\
Nahm boundary condition     & Nahm 边界条件       & \\
Nekrasov--Shatashvili limit & Nekrasov--Shatashvili 极限 & 缩写 NS 极限 \\
Omega deformation ($\Om$-background) & $\Om$-形变（$\Om$-背景） & \\
partition function          & 配分函数            & \\
instanton partition function & 瞬子配分函数       & \\
vev (vacuum expectation value) & 真空期望值       & \\
moduli space                & 模空间              & \\
Coulomb branch              & Coulomb 分支        & \\
Higgs branch                & Higgs 分支          & \\
Seiberg--Witten theory      & Seiberg--Witten 理论 & \\
class $\CS$ theories        & $\CS$ 类理论        & 保留花体 $\CS$ \\
conformal field theory (CFT) & 共形场论           & 缩写 CFT 保留 \\
conformal block             & 共形块              & 底本作“共形块”，不作“共形块函数” \\
Virasoro algebra            & Virasoro 代数       & \\
Kac--Moody algebra          & Kac--Moody 代数     & \\
Wess--Zumino--Witten model  & Wess--Zumino--Witten 模型 & 缩写 WZW \\
WZW model                   & WZW 模型            & \\
Liouville theory            & Liouville 理论      & \\
chiral algebra              & 手征代数            & \\
chiral ring                 & 手征环              & \\
central charge              & 中心荷              & \\
anomaly                     & 反常                & 不译“异常” \\
anomaly cancellation        & 反常抵消            & \\
trace anomaly               & 迹反常              & \\
operator product expansion (OPE) & 算符乘积展开   & 缩写 OPE 保留 \\
correlation function        & 关联函数            & \\
normal ordering             & 正规序              & \\
vertex operator             & 顶点算符            & \\
screening charge            & 屏蔽荷              & \\
Coulomb gas                 & Coulomb 气体        & \\
free field                  & 自由场              & \\
holomorphic / antiholomorphic & 全纯 / 反全纯     & \\
holonomy                    & 和乐（群）          & 此处用“和乐”，不用“完整群” \\
monodromy                   & 单值性              & 微分方程语境作“单值性” \\
monodromy matrix            & 单值矩阵            & \\
flat connection             & 平坦联络            & \\
complex structure           & 复结构              & \\
K\"ahler structure          & K\"ahler 结构       & \\
hyperk\"ahler               & 超 K\"ahler         & 约定：hyperk\"ahler manifold = 超 K\"ahler 流形 \\
hyperk\"ahler rotation      & 超 K\"ahler 旋转    & \\
Riemann surface             & Riemann 面          & \\
genus                       & 亏格                & \\
puncture                    & 穿刺点              & \\
marked point                & 标记点              & \\
Deligne--Mumford compactification & Deligne--Mumford 紧化 & \\
stable curve                & 稳定曲线            & \\
\end{longtable}

\medskip
\noindent
\textbf{数学与代数几何部分：}

\medskip
\noindent
\begin{longtable}{@{}p{0.30\textwidth}p{0.26\textwidth}p{0.36\textwidth}@{}}
\toprule
\textbf{英文原文} & \textbf{本译名} & \textbf{说明} \\
\midrule
\endfirsthead
\toprule
\textbf{英文原文} & \textbf{本译名} & \textbf{说明} \\
\midrule
\endhead
\bottomrule
\endfoot
geometric Langlands correspondence & 几何 Langlands 对应 & 程序/纲领作“Langlands 纲领” \\
Langlands program(me)       & Langlands 纲领       & 英式拼写从原文 \\
Hecke operator              & Hecke 算符          & \\
Hecke eigensheaf            & Hecke 本征层        & 原文 eigensheaf \\
Hecke eigenvalue            & Hecke 本征值        & \\
automorphic form            & 自守形式            & \\
automorphic representation  & 自守表示            & \\
Galois representation        & Galois 表示         & \\
local system                & 局部系统              & \\
$D$-module                   & $D$-模             & \\
$\CD$-module                 & $\CD$-模           & 保留原花体记号 \\
category of $\CD$-modules   & $\CD$-模范畴       & \\
Bun$_G$                     & $\Bun_G$            & 不译“G-丛模空间”全称 \\
principal $G$-bundle        & 主 $G$-丛           & \\
vector bundle               & 向量丛              & \\
line bundle                 & 线丛                & \\
connection                  & 联络                & \\
moduli stack                & 模叠                & \\
stack                       & 叠                  & \\
derived category            & 导出范畴            & \\
coherent sheaf              & 凝聚层              & \\
perverse sheaf              & 反常层              & 固定译法，注意与“反常”anomaly 区分 \\
constructible sheaf         & 可构造层            & \\
D-module cohomology         & $\CD$-模上同调      & \\
de Rham cohomology          & de Rham 上同调      & \\
Dolbeault cohomology        & Dolbeault 上同调    & \\
$\bar\partial$-cohomology   & $\bar\partial$-上同调 & \\
BRST cohomology             & BRST 上同调         & \\
conformal block             & 共形块              & \\
Verma module                & Verma 模            & \\
highest weight              & 最高权              & 表示论语境用“权”，不用“重量” \\
highest-weight module       & 最高权模            & \\
integrable representation   & 可积表示            & \\
Wakimoto module             & Wakimoto 模         & \\
screening operator          & 屏蔽算符            & \\
Knizhnik--Zamolodchikov (KZ) equation & Knizhnik--Zamolodchikov (KZ) 方程 & \\
Ward identity               & Ward 恒等式         & \\
Deutsch--Candelas?          & ---                 & \\
tt$^*$ equations            & tt$^*$ 方程         & 原文 $tt^*$ \\
Picard--Fuchs equation      & Picard--Fuchs 方程  & \\
Isomonodromy               & 等单值性            & \\
Stokes phenomenon           & Stokes 现象         & \\
WKB approximation           & WKB 近似            & \\
zeta-regularisation         & $\ze$-正规化        & \\
Hitchin moduli space        & Hitchin 模空间      & \\
Hitchin fibration           & Hitchin 纤维化      & \\
Higgs bundle / Higgs field  & Higgs 丛 / Higgs 场 & \\
spectral curve              & 谱曲线              & \\
Hitchin system              & Hitchin 系统        & \\
hyperk\"ahler quotient      & 超 K\"ahler 商      & \\
Nahm equations              & Nahm 方程           & \\
\end{longtable}


\end{document}
